\section{Theorem K II: agreement of theta integrals and local matching}\label{sec:pf:thmK}

\providecommand{\thmKAE}{\mathbb{A}_E}
\providecommand{\thmKS}{\mathcal{S}}
\providecommand{\thmKSo}{\mathcal{S}^{\circ}}
\providecommand{\thmKPol}{\mathcal{P}}
\providecommand{\thmKPD}{\operatorname{PD}}
\providecommand{\thmKI}{\mathcal{I}}
\providecommand{\thmKaug}{\mathfrak{a}}
\providecommand{\thmKCT}{\operatorname{CT}}
\providecommand{\thmKst}{\operatorname{st}}
\providecommand{\thmKht}{\operatorname{ht}}
\providecommand{\thmKqht}{\operatorname{qht}}
\providecommand{\thmKcone}{\operatorname{cone}}
\providecommand{\thmKEv}{\operatorname{Ev}}
\providecommand{\thmKpure}{\operatorname{pure}}
\providecommand{\thmKSpl}{\operatorname{Spl}}
\providecommand{\thmKSt}{\operatorname{St}}
\providecommand{\thmKacts}[1]{\mathrel{\stackrel{#1}{\leadsto}}}
\providecommand{\thmKLS}[2]{L^{S_0}(#1,#2)}
\providecommand{\thmKvac}{\mathrm{vac}}

This section proves Theorem~K (\cref{thm:theorem-K}): the two propositions $\sfname{ThetaIntegralsAgree}$ and
$\sfname{LocalMatching}$ of \cref{def:families:K} hold. Both are statements about \emph{all} configurations of
complementary hermitian spaces $V,V'$ (of dimensions $a>a'$) and a skew-hermitian space $W$ of dimension $a+a'$ over a
CM extension $E/F$, together with Weil representations $\omega,\omega'$ of $\Uni(V)\times\Uni(W)$ and
$\Uni(V')\times\Uni(W)$. $\sfname{LocalMatching}$ asserts the existence of local identifications (``matching maps'')
between quotients of the local Weil representations; $\sfname{ThetaIntegralsAgree}$ asserts convergence of the theta
integrals $\thmKI$ and $\thmKI'$ of $V$ and $V'$, the existence of a unique non-zero constant $c$ with
$\thmKI(\Phi)=c\,\thmKI'(\Phi')$ for all matched pairs, a description of the local coinvariants, and two recursion
formulas (``steps'') for $c$ when a hyperbolic plane is removed.

\paragraph{Strategy.}
The proof is an induction on the integer $p+q$ which bounds the real ranks of $\Uni(V)$ and $\Uni(V')$. In the definite
case $p=q=0$ all automorphic quotients are compact and the proportionality of $\thmKI$ and $\thmKI'$ follows from a
Hilbert space argument in which the Siegel--Weil formula enters. If $p+q>0$ one removes a hyperbolic plane from $V$
(if $q=0<p$) or from $V'$ (if $q>0$) and from $W$, and compares the constant terms of $\thmKI(\Phi)$ and
$\thmKI'(\Phi')$ along the maximal parabolic subgroup of $\Uni(W)$ stabilising an isotropic line: each constant term has
two components, distinguished by their characters under the Levi torus; one of them is an ``evaluation'' of the theta
integral of the smaller configuration, the other a ``cone integral''. The local behaviour of cone integrals under the
matching maps produces the partial $L$-values and epsilon factors of the steps; a cusp form argument (``end steps'')
shows that agreement of one component of the constant terms forces agreement of the theta integrals. All the inputs of
this argument are collected in twenty-six propositions, the \emph{leaves} (\cref{def:pf:thmK:leaves}); the deduction of
Theorem~K from the leaves (\cref{sec:pf:thmK:induction}) is given in full.

\paragraph{Organisation.}
\Cref{sec:pf:thmK:setting} defines the data and states the theorem. \Cref{sec:pf:thmK:inputs} states the two analytic
inputs (meromorphic continuation of Siegel Eisenstein series and the Siegel--Weil formula for definite spaces), which are
classical theorems, introduces the auxiliary notions (constant terms, evaluation maps, cone integrals, towers) and states
the leaves. \Cref{sec:pf:thmK:induction} deduces Theorem~K from the leaves. \Cref{sec:pf:thmK:leafproofs} proves the
leaves: each leaf is either proved completely or reduced, by a complete argument, to statements of the formal development
\cite{LeanSource2026} which we state precisely and cite. Complete arguments are given for the uniqueness of the
constant-term decomposition and the derivation of the constant-term leaves from the general constant-term formulas, the
towers, the non-vanishing of the step factors (Gauss sums), the construction of partners, the passage to arbitrary data on
a definite small side, the change of Cartan data, the non-vanishing of theta integrals via the doubling zeta integral, the
base case $p=q=0$, and the reduction of the local factors of the steps and of the end steps to local multiplicity one and
to cuspidality statements. \Cref{sec:pf:thmK:deferred} proves the main results used there through their formal statements:
partial $L$-functions at integers, an idelic Fubini theorem and the factorisation of the cone integral, the vanishing of
theta integrals of killed tensors, bump functions and the local cone integral at $v_1$, mixed models, the modulus identity and
the spherical functions at inert places, the unramified doubling integral, local multiplicity one (archimedean and
ramified), matching maps at split places, and the existence of Weil representations; the end steps are reduced to fifteen
statements in \cref{sec:pf:thmK:end}. The statements still cited are deeper ingredients: the general constant-term formulas,
some local computations, the cusp form statements of the end steps, the local theta correspondence and the analytic
estimates. \Cref{sec:pf:thmK:final} assembles the proof, and
\cref{sec:pf:thmK:remarks} records peculiarities of the formal statements.

% ===========================================================================
\subsection{Data, theta integrals and the statement}\label{sec:pf:thmK:setting}

\paragraph{Conventions.}
Throughout, $F$ and $E$ are number fields, $E$ is an $F$-algebra and $\sigma\colon E\to E$ is a homomorphism of
$F$-algebras. We write $\A=\A_F$, $\thmKAE=\A\otimes_FE$, $c_\sigma=\mathrm{id}\otimes\sigma$ on $\thmKAE$, $F_v$ and
$\OO_v$ for the completion and its integers at a finite place $v$ of $F$ (a height one prime of $\OO_F$), $\A_f$ for the
finite adeles and $F_\infty=F\otimes\R$. For a finite set $\kappa$ let $\thmKS(F_v^\kappa)$ be the space of locally
constant compactly supported functions on $F_v^\kappa$ (and likewise $\thmKS(\A_f^\kappa)$), and let the space of
\emph{test functions} $\thmKS(\A^\kappa)$ be the span of the functions $f\otimes\varphi\colon x\mapsto
f(x_\infty)\varphi(x_f)$ with $f$ a Schwartz function on $F_\infty^\kappa$ and $\varphi\in\thmKS(\A_f^\kappa)$. For a
Schwartz function $f$ and functions $\varphi_v$ on $F_v^\kappa$ ($v$ finite) put
$\thmKpure(f,\varphi)(x)=f(x_\infty)\prod_v\varphi_v(x_v)$ (a finite product when $\varphi_v(x_v)=1$ for almost all
$v$); for a place $u$ and a function $z$ on $F_u^\kappa$, $\varphi[u\mapsto z]$ is the family $\varphi$ with its
component at $u$ replaced by $z$. Sums $\sum_{\xi}$ over infinite sets are unconditional sums and integrals are Bochner
integrals; both are given the value $0$ when the sum or the integrand is not summable or integrable.

\emph{Local action.} Let $T$ be an operator on functions on $\A^\kappa$ and $v$ a finite place. For functions
$\varphi,\varphi_1$ on $F_v^\kappa$ we write $T\colon\varphi\thmKacts{v}\varphi_1$ (``$T$ acts at $v$ by
$\varphi\mapsto\varphi_1$'') if $T(\varphi\otimes_v r)=\varphi_1\otimes_v r$ for every function $r$ on $\A^\kappa$
which is invariant under translation by vectors supported at $v$ and such that
$(\varphi\otimes_v r)(x)=\varphi(x_v)r(x)$ is a test function. The archimedean analogue (with $F_\infty$ and Schwartz
functions) and the derivative analogue (for a one-parameter family $T_t$, with $\frac{d}{dt}|_{t=0}T_t$) are defined in
the same way.

\begin{definition}[Base data]\label{pf:thmK:def:base}
A \emph{base datum} $B=(\psi,(\tau_w)_{w\mid\infty},m_\alpha,m_\beta,P,(\tau_v)_{v\in P},v_1,\tau_1)$ consists of an
additive character $\psi\colon\A\to\C^\times$, for each infinite place $w$ of $F$ an embedding $\tau_w\colon E\to\C$
above $w$, functions $m_\alpha,m_\beta$ from the infinite places to $\Z$, a finite set $P$ of finite places with
$F$-algebra maps $\tau_v\colon E\to F_v$, and a finite place $v_1$ with an $F$-algebra map $\tau_1\colon E\to F_{v_1}$.
It is \emph{good} if: $E/F$ is a CM extension with complex conjugation $\sigma$; $E/\Q$ is Galois; $\OO_E/\OO_F$ is
unramified; $\psi$ is continuous, trivial on $F$, and $\psi(a)=\exp(-2\pi i\Tr_{F_\infty/\R}(a_\infty))$ for adeles $a$
supported at the infinite places; for every finite place $v$ inert in $E$ one has $2\notin v$, $\Nm v>3$ and $\psi$ is
unramified at $v$ (i.e.\ $\psi_v(x\OO_v)=1$ iff $x\in\OO_v$); and $\psi$ is unramified at every $v\in P$.
\formal{GlobalTheory.IndefiniteRowSideFullHecke.WeilTensorModels.Base.IsGood}
\end{definition}

\begin{definition}[Spaces and unitary groups]\label{pf:thmK:def:spaces}
A \emph{space} of dimension $n$ is a tuple $X=(J,L,\theta,\mathrm{bas})$ with $J\in M_n(E)$, $L$ an $\OO_E$-submodule
of $E^n$, matrices $\theta_w\in M_n(\C)$ for the infinite places and $\mathrm{bas}_v\in\GL_n(F_v)$ for the finite
places. Put $\langle x,y\rangle_J=x^{\mathsf T}J\sigma(y)$. For a base datum $B$, $X$ is \emph{hermitian of signature
$(r,s)$}, $X\in\mathrm{Herm}_B(r,s)$, if $\sigma(J_{ji})=J_{ij}$; for every infinite $w$ the complex matrix $J^{\tau_w}$
has signature $(r,s)$; $L$ is a lattice and $y\in L$ iff $\langle y,x\rangle_J$ is integral for all $x\in L$; each
$\theta_w$ is a Cartan involution of $J^{\tau_w}$ ($\theta_w^2=1$, $\theta_w^{\mathsf T}J^{\tau_w}\overline{\theta_w}=
J^{\tau_w}$, $\Re((\theta_wx)^{\mathsf T}J^{\tau_w}\bar x)>0$ for $x\ne0$); and for $v\in P$, $\mathrm{bas}_v$ is a
\emph{split basis} of $L$ at $v$ (a vector $u\in F_v^n$ is integral iff $\mathrm{bas}_vu$ lies in the $\tau_v$-component of
$L\otimes\OO_v$). It is \emph{skew-hermitian of signature $(r,s)$}, $X\in\mathrm{Skew}_B(r,s)$, if
$\sigma(J_{ji})=-J_{ij}$, the matrices $iJ^{\tau_w}$ have signature $(r,s)$, $L$ is a lattice with $y\in L$ iff
$\Tr_{E/\Q}(o\langle y,x\rangle_J)\in\Z$ for all $x\in L$, $o\in\OO_E$, the $\theta_w$ are Cartan involutions of
$iJ^{\tau_w}$, and the $\mathrm{bas}_v$ ($v\in P$) are split bases.

For $J\in M_n(E)$ let $\Uni(J)(\A)=\{g\in\GL_n(\thmKAE):g^{\mathsf T}Jc_\sigma(g)=J\}$, let $\Gamma_J\subset\Uni(J)(\A)$
be the subgroup of elements coming from $\GL_n(E)$, let $\thmKht(g)$ be the adelic height of $g$ (the product over all
places of $\max(1,\|g\|_v,\|g^{-1}\|_v)$) and $\thmKqht(g)=\inf_{\gamma\in\Gamma_J}\thmKht(\gamma g)$. We write
$\Uni(X)$ for $\Uni(J_X)$; $\Uni(X)(F_v)$ (resp.\ $\Uni(X)(F_\infty)$) is the set of $g$ with $g-1$ supported at $v$
(resp.\ at the infinite places); $K_\infty(X)$ is the set of $k\in\Uni(X)(F_\infty)$ whose components commute with all
$\theta_w$; a \emph{one-parameter subgroup} of $X$ is a continuous homomorphism $\R\to\Uni(X)(F_\infty)$.
\formal{GlobalTheory.IndefiniteRowSideFullHecke.WeilTensorModels.Spc.IsHerm}
\formal{GlobalTheory.IndefiniteRowSideFullHecke.WeilTensorModels.Spc.IsSkew}
\end{definition}

\begin{definition}[Models, Weil representations, Weil pairs]\label{pf:thmK:def:models}
Let $X,W$ be spaces of dimensions $n,d$. On $M_{n\times d}(\thmKAE)$ consider the alternating $\A$-bilinear form
$\langle x,y\rangle=\Tr_{\thmKAE/\A}\bigl(-\sum_{i,j,k,l}x_{ij}(J_X)_{ik}(J_W)_{jl}c_\sigma(y_{kl})\bigr)$. A \emph{model}
$M=(\mathbf e,(\thmKvac_v)_v,\mu,\mathrm{dom})$ for $(X,W)$ consists of Darboux coordinates
$\mathbf e\colon M_{n\times d}(\thmKAE)\xrightarrow{\sim}\A^{n\times d}\times\A^{n\times d}$ (so that $\langle w,w'\rangle
=\mathbf e(w)_1\cdot\mathbf e(w')_2-\mathbf e(w')_1\cdot\mathbf e(w)_2$), functions $\thmKvac_v$ on $F_v^{n\times d}$, a
measure $\mu$ on $\Uni(X)(\A)$ and a subset $\mathrm{dom}\subset\Uni(X)(\A)$. With
$\rho_\psi(w)\Phi(x)=\psi(x\cdot w_2+\tfrac12w_1\cdot w_2)\Phi(x+w_1)$, $(w_1,w_2)=\mathbf e(w)$, the model is
\emph{good} if $\mathbf e$ maps rational matrices to rational vectors; each $\thmKvac_v$ is a \emph{vacuum}, i.e.\ a
non-zero element of $\thmKS(F_v^{n\times d})$ such that, for some $\vartheta\in\OO_E$ with $\vartheta+\sigma(\vartheta)=1$
and every element $\gamma$ of the lattice generated by the matrices $a\,xy^{\mathsf T}$ ($a\in\A$ supported and integral at
$v$, $x\in L_X$, $y\in L_W$), the operator $\rho_\psi(\gamma)$ acts at $v$ on $\thmKvac_v$ by multiplication by
$\psi(\tfrac12\langle\gamma,\gamma\rangle_\vartheta)$, where $\langle\cdot,\cdot\rangle_\vartheta$ is the form above with $-1$
replaced by $\vartheta$; $\thmKvac_v=\one_{\OO_v^{n\times
d}}$ for almost all $v$; $\mu$ is a Haar measure; and $\mathrm{dom}$ is a measurable fundamental domain for $\Gamma_X$.

A \emph{representation} for $(X,W)$ is a monoid homomorphism $\omega$ from $\Uni(X)(\A)\times\Uni(W)(\A)$ to the
endomorphisms of the space of functions on $\A^{n\times d}$. For Hecke characters $\chi_X,\chi_W$ (below), $\omega$ is a
\emph{Weil representation for $M$ with splitting characters $(\chi_X,\chi_W)$} if: its restriction to $\Uni(W)(\A)$
preserves test functions and the linear forms on $\thmKS(\A^{n\times d})$ which are continuous in the archimedean
variable, is unitary and continuous on test functions, intertwines the Heisenberg operators
($\omega(h)\rho_\psi(x)=\rho_\psi(xh^{\mathsf T})\omega(h)$), and extends compatibly to a representation of the doubled
group $\Uni(W\oplus-W)(\A)$ with the same properties which satisfies Siegel's rule for the diagonal Lagrangian with the
character $\chi_X(\det t)\,|\Nm\det t|^{n/2}$; the same holds for the restriction to $\Uni(X)(\A)$ with the roles of
$X,W$ exchanged, the character $\chi_W$ and the exponent $d/2$; and $g\mapsto l(\omega(g)\Phi)$ is continuous for such
linear forms $l$ and test functions $\Phi$.

A \emph{Weil pair} for $(X,W)$ of signature $(r,s;r',s')$ with characters $(\chi_X,\chi_W)$ is a pair $(M,\omega)$ such
that $B$ is good, $X\in\mathrm{Herm}_B(r,s)$, $W\in\mathrm{Skew}_B(r',s')$, $M$ is good, $\chi_X,\chi_W$ are Hecke
characters with $\chi_X|_{\A^\times}=\eta^n$, $\chi_W|_{\A^\times}=\eta^d$, and $\omega$ is a Weil representation for $M$
with splitting characters $(\chi_X,\chi_W)$.
\formal{GlobalTheory.IndefiniteRowSideFullHecke.WeilTensorModels.Model.IsGood}
\formal{GlobalTheory.IndefiniteRowSideFullHecke.WeilTensorModels.Model.IsWeil}
\formal{GlobalTheory.ThK.IsPair}
\end{definition}

\begin{definition}[Fock space, polynomial data, the spaces $\thmKSo$]\label{pf:thmK:def:fock}
Let $M$ be a model for $(X,W)$. The \emph{Fock space} $\thmKPol(M)$ is the smallest space of Schwartz functions on
$F_\infty^{n\times d}$ which contains the Gaussians of $M$ (Schwartz functions on which the derivatives of the
Heisenberg lines $t\mapsto\rho_\psi(tu)$, $u$ archimedean, and $t\mapsto\rho_\psi(tu')$ with
$u'^{\tau_w}=\pm i\,\theta^X_wu^{\tau_w}(\theta^W_w)^{\mathsf T}$ differ by the factor $i$) and is stable under the
derivatives of the archimedean Heisenberg lines. The space of \emph{polynomial data} $\thmKPD(M)$ is the span of the
functions $\thmKpure(f,\varphi)$ with $f\in\thmKPol(M)$, $\varphi_v\in\thmKS(F_v^{n\times d})$ and
$\varphi_v=\thmKvac_v$ for almost all $v$. For a finite place $v$ and $t\colon E\to F_v$, $\thmKSo_v(t)$ is the set of
$\varphi\in\thmKS(F_v^{n\times d})$ whose image under some linear model at $v$ of conjugate type (a bijection of
$\thmKS$ onto the locally constant compactly supported functions on $M_{n\times d}(F_v)$ which turns the Heisenberg action
into translations and multiplications by characters) is supported on matrices of rank $n$; and $\thmKSo(M)$ is the span of
the functions $\thmKpure(f,\varphi)$ with $f$ Schwartz, $\varphi_v\in\thmKS(F_v^{n\times d})$, $\varphi_v=\thmKvac_v$
for almost all $v$ and $\varphi_{v_1}\in\thmKSo_{v_1}(\tau_1)$.
\formal{GlobalTheory.IndefiniteRowSideFullHecke.WeilTensorModels.Model.Pol}
\formal{GlobalTheory.ThK.PolData}
\formal{GlobalTheory.IndefiniteRowSideFullHecke.WeilTensorModels.Model.So}
\end{definition}

\begin{definition}[Theta functions and theta integrals]\label{pf:thmK:def:theta}
For a representation $\omega$ for $(X,W)$, a function $\Phi$ on $\A^{n\times d}$, $g\in\Uni(X)(\A)$ and
$h\in\Uni(W)(\A)$ put
\begin{equation}\label{pf:thmK:eq:theta}
\theta_\omega(\Phi)(g,h)=\sum_{\xi\in F^{n\times d}}(\omega(g,h)\Phi)(\xi),\qquad
\thmKI_{M,\omega}(\Phi)(h)=\int_{\mathrm{dom}}\theta_\omega(\Phi)(g,h)\,d\mu(g).
\end{equation}
We say that \emph{the theta integral of $\Phi$ converges} if $g\mapsto\theta_\omega(\Phi)(g,h)$ is integrable on
$\mathrm{dom}$ for every $h$. Note that $\thmKI_{M,\omega}$ depends on $M$ only through $(\mu,\mathrm{dom})$.
\formal{GlobalTheory.IndefiniteRowSideFullHecke.WeilTensorModels.Model.I}
\end{definition}

\begin{definition}[Characters, $L$-values, step factors]\label{pf:thmK:def:chars}
A character is a homomorphism $\chi\colon\thmKAE^\times\to S^1$; it is \emph{Hecke} if it is continuous and trivial on
$E^\times$. We write $\chi|_{\A^\times}=\eta^n$ if $\chi(x)=1$ for $x\in F^\times\Nm(\thmKAE^\times)$ and $\chi(x)=(-1)^n$
for the other $x\in\A^\times$. $\chi$ has \emph{infinity type} $m$ if $\chi(z)=\prod_w(z_w/|z_w|)^{m(w)}$ for ideles
$z$ with trivial finite components (computed through $\tau_w$). For a finite set $S$ of places of $F$,
$L^S(s,\chi)=\prod_{w}L_w(s,\chi)$ is the product over the finite places $w$ of $E$ not above $S$ of
$L_w(s,\chi)=(1-\chi_w(\varpi_w)\Nm w^{-s})^{-1}$ if $\chi$ is unramified at $w$ and $1$ otherwise. For a fractional
ideal $\mathfrak a$ of $E$, $\chi(\mathfrak a)=\prod_w\chi_w(\varpi_w)^{\ord_w\mathfrak a}$. For a character $\nu_v$ of
$F_v^\times$ with conductor exponent $c=c(\nu_v)$ and a uniformiser $\varpi$, the Gauss sum and epsilon factor are
\begin{equation}\label{pf:thmK:eq:gauss}
\begin{gathered}
g(\psi_v,\nu_v)=\sum_{y\in(\OO_v/\varpi^c)^\times}\nu_v(y)^{-1}\psi_v(y\varpi^{-c}),\\
\eps(\psi_v,\nu_v,z)=\Nm v^{-cz}\,g(\psi_v,\nu_v)\,\nu_v(\varpi)^c\qquad(z\in\Z).
\end{gathered}
\end{equation}
A \emph{character datum} $X=(\alpha,\beta,S_0,Q)$ consists of two characters and two finite sets of finite places; put
$\nu=\alpha/\beta$ and $T=S_0\cup Q$. A finite place $v$ is \emph{good} for $X$ if $v\notin T$ and the rational prime
below $v$ splits completely in $E$. For $v\in P$ let $\nu_v$ be the character of $F_v^\times$ obtained from $\nu$ through
the place of $E$ determined by $\tau_v$. For $e_1,e_2,s\in\{\pm1\}$, fractional ideals $\mathfrak b,\mathfrak c$ of $E$ and
$z\in\Z$, the \emph{step factor} is
\begin{equation}\label{pf:thmK:eq:stepfac}
\thmKst(e_1,e_2,s;\mathfrak b,\mathfrak c;z)=\nu(\mathfrak b)^{e_1}\nu(\mathfrak c)^{e_2}
\prod_{v\in P\cap T}\nu_v(s)\,\eps(\psi_v,\nu_v,z)
\end{equation}
(it depends on $B$ and $X$, which will be clear from the context). The datum $X$ is \emph{good} for $(B,a,a')$ if
$\alpha,\beta$ are Hecke characters with $\alpha|_{\A^\times}=\eta^{a'}$, $\beta|_{\A^\times}=\eta^{a}$ and infinity types
$m_\alpha,m_\beta$; $S_0\cap Q=\emptyset$; $T\subseteq P$; $\nu$ is ramified at every place of $E$ above $T$;
$\alpha,\beta$ are unramified outside $T$; and $v_1$ is good for $X$.
\formal{GlobalTheory.IndefiniteRowSideFullHecke.WeilTensorModels.stepFac}
\formal{GlobalTheory.IndefiniteRowSideFullHecke.WeilTensorModels.Chars.IsGood}
\end{definition}

\begin{definition}[Configurations and triples]\label{pf:thmK:def:geo}
A \emph{configuration} $G=(p,q,V,V',W,M,M')$ of type $(a,a')$ consists of integers $p,q\ge0$, spaces $V,V',W$ of
dimensions $a,a',a+a'$, a model $M$ for $(V,W)$ and a model $M'$ for $(V',W)$. It is \emph{good} for $B$ if: $a'<a$,
$a+a'$ is even and $a+a'<m_\alpha(w)-m_\beta(w)$ for all infinite $w$; $2(p+q)+1<\lfloor(a-a')/2\rfloor$ and
$2(p+q)+1<a'$; $V\in\mathrm{Herm}_B(a-p,p)$, $V'\in\mathrm{Herm}_B(q,a'-q)$, $W\in\mathrm{Skew}_B(a+a'-(p+q),p+q)$;
and $M,M'$ are good. The data $(B,G,X,\omega,\omega')$ form a \emph{triple} if $B$ is good, $G$ is good for $B$, $X$ is
good for $(B,a,a')$, $\omega$ is a Weil representation for $M$ with splitting characters $(\beta,\beta/\alpha)$ and
$\omega'$ one for $M'$ with splitting characters $(\alpha,\alpha/\beta)$. We write $\thmKI=\thmKI_{M,\omega}$ and
$\thmKI'=\thmKI_{M',\omega'}$.
\formal{GlobalTheory.IndefiniteRowSideFullHecke.WeilTensorModels.Geo.IsGood}
\formal{GlobalTheory.IndefiniteRowSideFullHecke.WeilTensorModels.IsTriple}
\end{definition}

\begin{lemma}[The Weil pairs of a triple]\label{pf:thmK:pairs}
For a triple, $(M,\omega)$ is a Weil pair for $(V,W)$ of signature $(a-p,p;a+a'-(p+q),p+q)$ with characters
$(\beta,\beta/\alpha)$, and $(M',\omega')$ is a Weil pair for $(V',W)$ of signature $(q,a'-q;a+a'-(p+q),p+q)$ with
characters $(\alpha,\alpha/\beta)$; we call them the Weil pairs of $V$ and of $V'$. Moreover $(a+a')+2\min(a-p,p)<2a$.
\formallabel{PositivityTheory.ThKG.pairV}\formallabel{PositivityTheory.ThKG.pairV'}
\formallabel{PositivityTheory.ThKG.weilCondV}
\end{lemma}

\begin{proof}
The quotient of two Hecke characters is a Hecke character, and $(\beta/\alpha)|_{\A^\times}=\eta^{a}\eta^{-a'}=\eta^{a+a'}$
since $\eta^2=1$; similarly for $\alpha/\beta$. The inequality follows from $2(p+q)+1<\lfloor(a-a')/2\rfloor$.
\end{proof}

\begin{definition}[Matching maps]\label{pf:thmK:def:maps}
Let $(B,G,X,\omega,\omega')$ be a triple. For a finite place $v$ the \emph{augmentation module} $\thmKaug_v(\omega)$ is the
span of the differences $\varphi_1-\varphi$ with $\varphi,\varphi_1\in\thmKS(F_v^{a\times(a+a')})$ and
$\omega(g,1)\colon\varphi\thmKacts{v}\varphi_1$ for some $g\in\Uni(V)(F_v)$; similarly $\thmKaug_v(\omega')$.
\begin{enumerate}
\item \emph{Real maps} $R=(\Pi_\infty,p,p')$: a complex vector space with linear maps $p,p'$ from the archimedean
Schwartz spaces of $F_\infty^{a\times(a+a')}$, $F_\infty^{a'\times(a+a')}$ to $\Pi_\infty$. They are \emph{matching}
if $p(\thmKPol(M))=\Pi_\infty$, $p$ kills the derivatives of $t\mapsto\omega(\gamma(t),1)f$ ($f\in\thmKPol(M)$, $\gamma$ a
one-parameter subgroup of $V$) and $p(\omega(k,1)f)=p(f)$ for $k\in K_\infty(V)$; the same for $p'$; if $p(f)=p'(f')$
($f\in\thmKPol(M)$, $f'\in\thmKPol(M')$) then $p(f_1)=p'(f'_1)$ for the derivatives $f_1,f'_1$ of $f,f'$ along the same
one-parameter subgroup of $W$ and for the translates by the same $k\in K_\infty(W)$; and $\Pi_\infty$ is non-zero and every
subspace stable under the action of the one-parameter subgroups of $W$ and of $K_\infty(W)$ induced through $p$ is $0$ or
$\Pi_\infty$.
\item \emph{Finite maps} $R_f=(\Pi_v,p_v,p'_v)_v$: for every finite $v$ a complex vector space with linear maps
$p_v,p'_v$ from the functions on $F_v^{a\times(a+a')}$, $F_v^{a'\times(a+a')}$. They are \emph{matching} if for every
$v$: $p_v,p'_v$ map $\thmKS$ onto $\Pi_v$ and vanish on $\thmKaug_v(\omega)$, $\thmKaug_v(\omega')$; if
$p_v(\varphi)=p'_v(\varphi')$ and $h\in\Uni(W)(F_v)$ acts by $\varphi\mapsto\varphi_1$, $\varphi'\mapsto\varphi'_1$
under $\omega(1,h)$, $\omega'(1,h)$, then $p_v(\varphi_1)=p'_v(\varphi'_1)$; $\Pi_v\ne0$ and every subspace stable under
the induced action of $\Uni(W)(F_v)$ is $0$ or $\Pi_v$; for $v\notin T$, $p_v(\thmKvac_v)=p'_v(\thmKvac'_v)\ne0$; and for
$v\in P\cap T$ the \emph{zeta condition} holds: there are linear models $\iota,\iota'$ at $v$ adapted to the bases
$\mathrm{bas}_v$, sending the vacua to the characteristic functions of the integral matrices, such that
$p_v(\varphi)=p'_v(\varphi')$ iff two explicit zeta integrals of Godement--Jacquet type \cite{GodementJacquet1972}
attached to $(\iota\varphi,\nu_v,\beta_v)$ and $(\iota'\varphi',\nu_v,\alpha_v)$ agree as functions on
$\GL_{a+a'}(F_v)$.
\end{enumerate}
The data $(B,G,X,\omega,\omega',R,R_f)$ are \emph{matched} if they form a triple and $R,R_f$ are matching.
\formal{GlobalTheory.IndefiniteRowSideFullHecke.WeilTensorModels.IsRealMaps}
\formal{GlobalTheory.IndefiniteRowSideFullHecke.WeilTensorModels.IsFinMaps}
\formal{GlobalTheory.IndefiniteRowSideFullHecke.WeilTensorModels.IsZetaAt}
\end{definition}

\begin{definition}[Matched functions and constants]\label{pf:thmK:def:match}
Let $R,R_f$ be real and finite maps. Functions $\Phi$ on $\A^{a\times(a+a')}$ and $\Phi'$ on $\A^{a'\times(a+a')}$ are
\emph{polynomially matched}, written $\Phi\approx\Phi'$, if there are a finite set $T'\supseteq T$ of finite places and
decompositions $\Phi=\sum_i\thmKpure(f_i,\varphi_i)$, $\Phi'=\sum_j\thmKpure(f'_j,\varphi'_j)$ with $f_i\in\thmKPol(M)$,
$f'_j\in\thmKPol(M')$, all local components in $\thmKS$ and equal to the vacua outside $T'$, such that
\begin{equation}\label{pf:thmK:eq:polmatch}
\sum_ip(f_i)\otimes\bigotimes_{v\in T'}p_v(\varphi_{i,v})=\sum_jp'(f'_j)\otimes\bigotimes_{v\in T'}p'_v(\varphi'_{j,v})
\quad\text{in }\Pi_\infty\otimes\textstyle\bigotimes_{v\in T'}\Pi_v.
\end{equation}
They are \emph{matched}, written $\Phi\sim\Phi'$, if $\Phi'\in\thmKSo(M')$ and there are decompositions
$\Phi=\sum_if_i\otimes g_i$, $\Phi'=\sum_jf'_j\otimes g'_j$ ($f$ Schwartz, $g\in\thmKS(\A_f)$) and sequences
$s_{k,i}\to f_i$, $s'_{k,j}\to f'_j$ in the Schwartz topology such that for every $k$ the functions
$\Phi_k=\sum_is_{k,i}\otimes g_i$, $\Phi'_k=\sum_js'_{k,j}\otimes g'_j$ satisfy $\Phi'_k\in\thmKSo(M')$ and
$\Phi_k\approx\Phi'_k$. A number $c\in\C$ is a \emph{constant} if $\thmKI(\Phi)(h)=c\,\thmKI'(\Phi')(h)$ for all
$\Phi\sim\Phi'$ and all $h\in\Uni(W)(\A)$.
\formal{GlobalTheory.IndefiniteRowSideFullHecke.WeilTensorModels.PolMatch}
\formal{GlobalTheory.IndefiniteRowSideFullHecke.WeilTensorModels.Match}
\formal{GlobalTheory.IndefiniteRowSideFullHecke.WeilTensorModels.IsConst}
\end{definition}

\begin{definition}[Decay, one constant, coinvariants]\label{pf:thmK:def:decay}
A triple satisfies the \emph{decay condition} if the theta integral of every test function $\Phi$ converges, and for every
$\Phi'\in\thmKSo(M')$: for every $N$ there are $C,k$ with $|\theta_{\omega'}(\Phi')(g,h)|\,\thmKqht(g)^N\le
C\thmKht(h)^k$ for all $g,h$, the theta integral of $\Phi'$ converges, and $|\thmKI'(\Phi')(h)|\le C\thmKht(h)^N$ for
some $C,N$. Matched data satisfy the \emph{one-constant condition} if there is $\Phi'\in\thmKSo(M')$ with
$\thmKI'(\Phi')\ne0$, there is exactly one constant, and every constant is non-zero. They satisfy the \emph{coinvariant
condition} if for every place $v$ good for $X$: (i) $\ker p_v\cap\thmKS=\thmKaug_v(\omega)$ and $\ker p'_v\cap\thmKS=
\thmKaug_v(\omega')$; (ii) for every $t\colon E\to F_v$, $p'_v$ maps $\thmKSo_v(t)$ onto $\Pi_v$; (iii) every linear form on
the functions on $F_v^{a'\times(a+a')}$ vanishing on $\thmKaug_v(\omega')$ factors on $\thmKS$ through $p'_v$; and
moreover $\thmKSo(M')$ is stable under all operators $\omega'(\gamma)$.
\formal{GlobalTheory.IndefiniteRowSideFullHecke.WeilTensorModels.Decay}
\formal{GlobalTheory.IndefiniteRowSideFullHecke.WeilTensorModels.OneConst}
\formal{GlobalTheory.IndefiniteRowSideFullHecke.WeilTensorModels.Coinv}
\end{definition}

\begin{definition}[Adapted planes and step relations]\label{pf:thmK:def:steps}
A \emph{plane} $H=(x,y,\iota,\mathfrak b)$ of dimensions $(n,n_1)$ consists of $x,y\in E^n$,
$\iota\in M_{n\times n_1}(E)$ and a fractional ideal $\mathfrak b=\mathfrak b_H$ of $E$. It is \emph{adapted} to spaces
$(Y,Y_1)$ at positions $(i_0,i_1)$ if $n=n_1+2$; $\langle x,x\rangle=\langle y,y\rangle=0$, $\langle x,y\rangle=1$; the
image of $\iota$ is orthogonal to $x,y$ and $\iota^{\mathsf T}J_Y\sigma(\iota)=J_{Y_1}$; $u\in L_{Y_1}$ iff $\iota u\in
L_Y$, and every $z\in L_Y$ is $sx+\iota u+ty$ with $sx,ty\in L_Y$ and $u\in L_{Y_1}$; $\mathfrak b=\{s:sx\in L_Y\}$;
$\theta^Y_w\iota=\iota\theta^{Y_1}_w$ and $y^{\tau_w}\in\C^\times\theta^Y_wx^{\tau_w}$ for all $w$; and for $v\in P$ the
columns $i_0,i_1$ of $\mathrm{bas}_v(Y)$ are $\tau_v(x),\tau_v(y)$ and the other columns are, in order, those of
$\iota^{\tau_v}\mathrm{bas}_v(Y_1)$.

For $e_1,e_2,s\in\{\pm1\}$ and fixed $F,E,\sigma$:
\begin{enumerate}
\item the \emph{$\beta$-step relation} $\thmKSt_\beta(e_1,e_2,s)$: for all $B$, $a,b$, configurations $G$ of type
$(a,b+2)$ and $G_1$ of type $(a,b)$ with $V(G_1)=V(G)$, $p(G_1)=p(G)$, $q(G_1)+1=q(G)$, planes $H$ adapted to
$(V'(G),V'(G_1))$ at $(0,b+1)$ and $H'$ adapted to $(W(G),W(G_1))$ at $(a,a+b+1)$, and all real maps $R,R_1$, there is
$C\ne0$ such that for all character data $X$, all matched data $(\omega,\omega',R_f)$ for $(G,R)$ and
$(\omega_1,\omega'_1,R_{f,1})$ for $(G_1,R_1)$ and all constants $c$ for $G$ and $c_1$ for $G_1$, there is $u$ with
\begin{equation}\label{pf:thmK:eq:stepbeta}
u\,C\,\thmKst(e_1,e_2,s;\mathfrak b_H,\mathfrak b_{H'};-k)=1,\qquad c\,\thmKLS{k+1}{\nu}=u\,c_1,\qquad
k=\lfloor(a-(b+2))/2\rfloor;
\end{equation}
\item the \emph{$\alpha$-step relation} $\thmKSt_\alpha(e_1,e_2,s)$: for all $B$, $b,a'$, configurations $G$ of type
$(b+2,a')$ and $G_1$ of type $(b,a')$ with $V'(G_1)=V'(G)$, $q(G)=q(G_1)=0$, $p(G_1)+1=p(G)$, planes $H$ adapted to
$(V(G),V(G_1))$ at $(0,b+1)$ and $H'$ adapted to $(W(G),W(G_1))$ at $(0,b+1)$, and all $R,R_1$, there is $C\ne0$ such
that, with the same quantifiers,
\begin{equation}\label{pf:thmK:eq:stepalpha}
c=C\,\thmKst(e_1,e_2,s;\mathfrak b_H,\mathfrak b_{H'};1-k)\,\thmKLS{k}{\nu}\,c_1,\qquad k=\lfloor(b+2-a')/2\rfloor.
\end{equation}
\end{enumerate}
\formal{GlobalTheory.IndefiniteRowSideFullHecke.WeilTensorModels.Plane.IsAdapted}
\formal{GlobalTheory.IndefiniteRowSideFullHecke.WeilTensorModels.StepBeta}
\formal{GlobalTheory.IndefiniteRowSideFullHecke.WeilTensorModels.StepAlpha}
\end{definition}

\begin{definition}[The two notions of Theorem K]\label{pf:thmK:def:K}
$\sfname{LocalMatching}$ is the statement: for all number fields $F,E$ with $E$ an $F$-algebra, all $\sigma,B,a,a'$,
configurations $G$ of type $(a,a')$, character data $X$ and representations $\omega,\omega'$ forming a triple, there are
real maps $R$ and finite maps $R_f$ such that the data are matched. $\sfname{ThetaIntegralsAgree}$ is the conjunction of:
(i) every triple satisfies the decay condition, and all matched data satisfy the one-constant condition and the
coinvariant condition; (ii) there are $e_1,e_2,s\in\{\pm1\}$ with $\thmKSt_\beta(e_1,e_2,s)$ for all $F,E,\sigma$; (iii)
there are $e_1,e_2,s\in\{\pm1\}$ with $\thmKSt_\alpha(e_1,e_2,s)$ for all $F,E,\sigma$.
\formal{GlobalTheory.IndefiniteRowSideFullHecke.WeilTensorModels.LocalMatching}
\formal{GlobalTheory.IndefiniteRowSideFullHecke.WeilTensorModels.ThetaIntegralsAgree}
\end{definition}

\begin{theorem}[\cref{thm:theorem-K}]\label{pf:thmK:main}
$\sfname{ThetaIntegralsAgree}$ and $\sfname{LocalMatching}$ hold.
\formallabel{UpperTheory.LiteratureTheorems.theta_integrals_agree_and_local_matching_of_literature}
\end{theorem}

The formal statement has a hypothesis (a proof of three propositions on Galois representations and on the local Langlands
correspondence) which is not used (\cref{pf:thmK:rem:unused}). The proof is completed in \cref{sec:pf:thmK:final}.


% ===========================================================================
\subsection{Analytic inputs, auxiliary notions and the leaves}\label{sec:pf:thmK:inputs}

\subsubsection{Siegel Eisenstein series and the Siegel--Weil formula}\label{sec:pf:thmK:sw}
For a space $Y$ of dimension $n$ let $G_2(Y)=\Uni(Y\oplus-Y)(\A)$ be the doubled group, $\Delta$ the diagonal Lagrangian
and, for $p\in G_2(Y)$ and $t\in\GL_n(\thmKAE)$, write $p\in P_\Delta$ with Levi part $t$ if $p\Delta=\Delta t$. Frame
data $(b,K,a)$ consist of a basis $b$ of $E^{2n}$ in which the doubled form is $\left(\begin{smallmatrix}0&1\\
\pm1&0\end{smallmatrix}\right)$ and whose first $n$ vectors span $\Delta$, the subgroup $K$ of elements whose matrices
in $b$ are unitary at the infinite places and, together with their inverses, integral at the finite places, and a
function $a$ with $G_2(Y)=P_\Delta K$ and $a(pk)=|\det t|_{\thmKAE}$. A \emph{section} (for a Hecke character $\chi$ and
$s_1\in\C$) is a continuous right $K$-finite function $f$ on $G_2(Y)$ with $f(px)=\chi(\det t)|\det t|^{s_1+n/2}f(x)$;
its flat deformation is $f_s=a^{s-s_1}f$, and its Eisenstein series is $\sum_{\gamma\in P_\Delta(F)\backslash
G_2(Y)(F)}f_s(\gamma x)$.

\begin{littheorem}[Meromorphic continuation; \cite{Langlands1976,MoeglinWaldspurger1995}]\label{pf:thmK:eis}
Let $B$ be good, $n\ge1$, $Y$ hermitian or skew-hermitian of dimension $n$, $\chi$ a Hecke character with
$\chi|_{\A^\times}=\eta^{c}$ for some $c\in\N$, $(b,K,a)$ frame data and $f$ a section. There is $\hat E(s,x)$ such that the Eisenstein
series converges to $\hat E(s,x)$ for $\Re s>t$ (some $t$) and is holomorphic there, and for every $s_0$ there are $k\ge0$,
a neighbourhood $U$ of $s_0$ and a family $H(s)$ of functions of moderate growth with $H(s)=(s-s_0)^k\hat E(s,\cdot)$ on
$U\smallsetminus\{s_0\}$ such that $s\mapsto\mu(H(s))$ is holomorphic on $U$ for every growth functional $\mu$.
\formal{GlobalTheory.ThK.Dbl.EisBL}
\formallabel{PositivityTheory.RgEis.Hd.eisBL}\formallabel{UpperTheory.LiteratureTheorems.eis_bl_of_literature}
\end{littheorem}

This is the meromorphic continuation of Siegel Eisenstein series, in a weak (functional-wise) form. In the development it
is proved from the convergence in a half plane and a multiplicativity argument for the constant terms; we use it as a
classical theorem.

\begin{littheorem}[Siegel--Weil formula for definite spaces; \cite{Weil1965,KudlaRallis1988,Ichino2004}]\label{pf:thmK:sw:thm}
Let $B$ be good, $m,n\ge1$, and let $A$, $Y$ be spaces of dimensions $m,n$, one hermitian and the other
skew-hermitian, with $A$ of signature $(r,s)$, $\min(r,s)=0$. Let $\chi$ be a Hecke character with
$\chi|_{\A^\times}=\eta^m$, and let $(\mathbf d,\omega,S,l)$ be a doubling datum: rational Darboux coordinates for
$A\otimes(Y\oplus-Y)$, a Weil representation $\omega$ of $G_2(Y)$ satisfying Siegel's rule for $\Delta$ with the
character $\chi$, a representation $S$ of $\Uni(A)(\A)$ by linear substitutions, and an admissible linear form $l$
invariant under $S$ and under the Heisenberg translations by $z\Delta^{\mathsf T}$, with
$\sum_{y}l(\rho_\psi(y,-y)\Phi)=\sum_\xi\Phi(\xi)$ for test functions $\Phi$. Let $(b,K,a)$ be frame data, $\mu$ a Haar
measure and $D$ a fundamental domain for $\Gamma_A$. Then for every test function $\Phi$ whose $\omega(K)$-translates span
a finite-dimensional space, the Eisenstein series of the flat deformation of the section $x\mapsto l(\omega(x)\Phi)$ has,
at every $x$, a meromorphic continuation which is holomorphic at $s_0=(m-n)/2$ with value
\begin{equation}\label{pf:thmK:eq:sw}
E(x,s_0)=\kappa\,\frac1{\mu(D)}\int_D\sum_\xi\bigl(\omega(x)S(g)\Phi\bigr)(\xi)\,d\mu(g),\qquad
\kappa=\begin{cases}1&n<m,\\2&n\ge m.\end{cases}
\end{equation}
\formal{GlobalTheory.ThK.Dbl.SiegelWeilStd}
\formallabel{PositivityTheory.SiegelWeilFormula.Fin.siegelWeilStd}
\formallabel{UpperTheory.LiteratureTheorems.siegel_weil_std_of_literature}
\end{littheorem}

Since $A$ is definite at the infinite places, it is anisotropic, and \eqref{pf:thmK:eq:sw} is the classical Siegel--Weil
formula for the dual pair $(\Uni(A),\Uni(Y\oplus-Y))$ with its usual constant $\kappa\in\{1,2\}$
\cite{Weil1965,KudlaRallis1988,Ichino2004}. In the development it is proved (via Poisson summation, a splitting of the
constant terms and a comparison of Whittaker--Fourier coefficients); we use it as a classical theorem.

\subsubsection{Auxiliary notions}\label{sec:pf:thmK:aux}
Fix $B$, a configuration $G$ of type $(a,a')$, a character datum $X$ and representations $\omega,\omega'$; when matched
data are given, $R,R_f$ denote the matching maps.

\begin{definition}[Polynomial constants]\label{pf:thmK:def:polconst}
$c\in\C$ is a \emph{polynomial constant} if $\thmKI(\Phi)(h)=c\,\thmKI'(\Phi')(h)$ for all $\Phi$, all
$\Phi'\in\thmKSo(M')$ with $\Phi\approx\Phi'$ and all $h$.
\formal{GlobalTheory.ThK.PolConst}
\end{definition}

\begin{definition}[Constant terms along a plane; splittings]\label{pf:thmK:def:ct}
Let $H'=(x,y,\iota,\mathfrak b)$ be a plane adapted to $(W,W_1)$. Let $N(H')\subseteq\Uni(W)(\A)$ be the subgroup
generated by the $u$ with $ux=x$ and $u\iota z\in\iota z+\thmKAE x$ for all $z$. For functions $f,c$ on $\Uni(W)(\A)$,
$\thmKCT_{H'}(f)=c$ means: $c(h)=\int_{\mathcal F}f(uh)\,du$ for every Haar measure on $N(H')$ and every fundamental
domain $\mathcal F$ of $N(H')\cap\Gamma_W$ of volume one. An element $m\in\Uni(W)(\A)$ is a \emph{Levi element}
$m=m(t,h_1)$ ($t\in\thmKAE^\times$, $h_1\in\Uni(W_1)(\A)$) if $mx=tx$, $my=c_\sigma(t)^{-1}y$ and $m\iota=\iota h_1$. A
function $f$ \emph{has law} $(\chi,z)$ if $f(mh)=\chi(t)|t|_{\thmKAE}^zf(h)$ for all Levi elements $m=m(t,1)$ and all
$h$. For a configuration of type $(a,a')$, a \emph{splitting} of a function $f$ on $\Uni(W)(\A)$ along $H'$ is a pair
$(f_\alpha,f_\beta)$ of functions with
\begin{equation}\label{pf:thmK:eq:split2}
\thmKCT_{H'}(f)=f_\alpha+f_\beta,\qquad f_\alpha\text{ of law }(\alpha,a'/2),\qquad f_\beta\text{ of law }(\beta,a/2);
\end{equation}
we write $(f_\alpha,f_\beta)\in\thmKSpl_{H'}(f)$, or $\thmKSpl(f)$. Clearly $(cf_\alpha,cf_\beta)\in\thmKSpl(cf)$ if
$(f_\alpha,f_\beta)\in\thmKSpl(f)$, and $(0,0)\in\thmKSpl(0)$.
\formal{GlobalTheory.ThK.IsCT}\formal{GlobalTheory.ThK.IsLevi}\formal{GlobalTheory.ThK.Split2}
\end{definition}

\begin{definition}[Evaluation maps, stabilisers, cone integrals]\label{pf:thmK:def:ev}
Let $M$ be a model for $(X,W)$ and $M_1$ one for $(X_1,W_1)$, $\iota\in M_{n\times n_1}(E)$, $\iota'\in
M_{d\times d_1}(E)$, and $I$ a set of matrices in $M_{n\times d}(\thmKAE)$. A linear map $e$ from functions on
$\A^{n\times d}$ to functions on $\A^{n_1\times d_1}$ is an \emph{evaluation map}, $e\in\thmKEv(M,M_1;\iota,\iota';I)$,
if it preserves test functions; $e(\rho_\psi(\iota z\iota'^{\mathsf T})\Phi)=\rho_\psi(z)e(\Phi)$ and
$e(\rho_\psi(u)\Phi)=e(\Phi)$ for test functions $\Phi$, $z\in M_{n_1\times d_1}(\thmKAE)$ and $u\in I$; and there are a
non-zero continuous linear map $e_\infty$ of Schwartz spaces and linear maps $e_v$ with $e_v(\thmKvac_v)=\thmKvac^{M_1}_v$
and $e_v(\thmKS)=\thmKS$ such that $e(\thmKpure(f,\varphi))=\thmKpure(e_\infty f,(e_v\varphi_v)_v)$. For planes $H$ in $X$
and $H'$ in $W$ put $I_0(H')=\{z\,x_{H'}^{\mathsf T}\}$, $I_A(H,H')=I_0(H')\cup\{x_H(\iota_{H'}z)^{\mathsf T}\}$ and
$I_B(H,H')=\{x_Hz^{\mathsf T}\}\cup\{(\iota_Hz)x_{H'}^{\mathsf T}\}$ ($z$ ranging over adelic vectors).

A subgroup $K\subseteq\Uni(X)(\A)$ is the \emph{stabiliser} of $X$ if $g\in K$ iff the components of $g$ commute with all
$\theta_w$ and $g,g^{-1}$ preserve the closure of the group generated by the vectors $ax$ ($x\in L_X$, $a\in\A$ finite and
integral). For such $K$ with a Haar measure $\mu_K$, a Haar measure $\mu$ on $\thmKAE^\times$, a map
$u\colon\thmKAE^\times\to M_{n\times d}(\thmKAE)$, a character $\nu$ and $z\in\N$, the \emph{cone integral} is
\begin{equation}\label{pf:thmK:eq:cone}
\thmKcone(\Phi)(y)=\int_K\int_{\thmKAE^\times}e\bigl(\rho_\psi(u(l))\,\omega(k,1)\Phi\bigr)(y)\,\nu(l)\,|l|^z_{\thmKAE}
\,d\mu(l)\,d\mu_K(k).
\end{equation}
We use $u=u_A(l)=l\,x_Hy_{H'}^{\mathsf T}$ and $u=u_B(l)=c_\sigma(l)\,y_Hx_{H'}^{\mathsf T}$.
\formal{GlobalTheory.ThK.IsEv}\formal{GlobalTheory.ThK.IsStab}\formal{GlobalTheory.ThK.cone}
\end{definition}

\begin{definition}[Re-choices and step configurations]\label{pf:thmK:def:cfg}
For spaces $X,Y$ of the same dimension let $X[Y]$ be $X$ with its Cartan involutions and local bases replaced by those of
$Y$; for a configuration $G$, $G[V_2,V'_2,W_2]$ is $G$ with $V,V',W$ replaced by $V[V_2]$, $V'[V'_2]$, $W[W_2]$ and the same
models, integers and theta integrals; we call it a \emph{re-choice} of $G$. For $G$ of type $(a,b+2)$ and $G_1$ of type
$(a,b)$, $(G,G_1,H,H')$ is a \emph{$\beta$-step configuration} if $V(G_1)=V(G)$, $p(G_1)=p(G)$, $q(G_1)+1=q(G)$, $H$ is
adapted to $(V'(G),V'(G_1))$ at $(0,b+1)$ and $H'$ to $(W(G),W(G_1))$ at $(a,a+b+1)$. For $G$ of type $(b+2,a')$ and $G_1$
of type $(b,a')$, it is an \emph{$\alpha$-step configuration} if $V'(G_1)=V'(G)$, $q(G)=q(G_1)=0$, $p(G_1)+1=p(G)$, $H$ is
adapted to $(V(G),V(G_1))$ at $(0,b+1)$ and $H'$ to $(W(G),W(G_1))$ at $(0,b+1)$. In both cases
$p(G_1)+q(G_1)=p(G)+q(G)-1$.
\formal{GlobalTheory.ThK.reG}\formal{GlobalTheory.ThK.CfgB}\formal{GlobalTheory.ThK.CfgA}
\end{definition}

In an $\alpha$-step configuration we write $\thmKcone_A$ for the cone integral \eqref{pf:thmK:eq:cone} on the $V$-side with
$u=u_A$, $\nu=\alpha/\beta$ and $z=k_A:=\lfloor(b+2-a')/2\rfloor$; in a $\beta$-step configuration, $\thmKcone_B$ is the cone
integral on the $V'$-side with $u=u_B$, $\nu$ and $z=k_B+1$, $k_B:=\lfloor(a-(b+2))/2\rfloor$. The theta integrals of $G_1$
are denoted $\thmKI_1,\thmKI'_1$, and its models $M_1,M'_1$.

\begin{definition}[Non-degenerate data]\label{pf:thmK:def:claim}
Matched data are \emph{non-degenerate} if there is a non-zero polynomial constant, there is $\Phi\in\thmKPD(M)$ with
$\thmKI(\Phi)\ne0$, and there is $\Phi'\in\thmKPD(M')\cap\thmKSo(M')$ with $\thmKI'(\Phi')\ne0$. For $n\in\N$ let
$\mathcal B(n)$ be the statement that all matched data (over all $F,E,\sigma,B,a,a',\dots$) with $p+q<n$ are
non-degenerate. Clearly $\mathcal B(n)$ implies $\mathcal B(m)$ for $m\le n$.
\formal{PositivityTheory.ThKG.Claim}\formal{PositivityTheory.ThKG.Below}
\end{definition}

\subsubsection{The leaves}\label{sec:pf:thmK:leavesdef}

\begin{definition}[The leaves]\label{def:pf:thmK:leaves}
The \emph{leaves} are the following twenty-six statements, each universally quantified over all number fields $F,E$,
$\sigma$, base data, configurations, character data, representations and matching maps (``a triple'', ``matched data''
and ``a Weil pair'' refer to \cref{pf:thmK:def:geo,pf:thmK:def:maps,pf:thmK:def:models}).
\begin{enumerate}[label=(L\arabic*),leftmargin=*]
\item\label{pf:thmK:L:evzero} \emph{Evaluation along $W$.} For a Weil pair for $(X,W)$, a plane $H'$ adapted to $(W,W_1)$
and a good model $M_1$ for $(X,W_1)$ there is $e\in\thmKEv(M,M_1;1,\iota_{H'};I_0(H'))$, and every such $e$ maps
$\thmKPD(M)$ onto $\thmKPD(M_1)$.
\item\label{pf:thmK:L:evcone} \emph{Evaluation along $X$ and $W$.} For a Weil pair for $(X,W)$, planes $H$ adapted to
$(X,X_1)$ and $H'$ adapted to $(W,W_1)$ and a good model $M_1$ for $(X_1,W_1)$, the sets
$\thmKEv(M,M_1;\iota_H,\iota_{H'};I_A(H,H'))$ and $\thmKEv(M,M_1;\iota_H,\iota_{H'};I_B(H,H'))$ are non-empty.
\item\label{pf:thmK:L:stab} \emph{Stabilisers.} For good $B$ and $X\in\mathrm{Herm}_B(r,s)$ the stabiliser $K$ of $X$
exists and is compact, and there are Haar measures on $K$ and on $\thmKAE^\times$.
\item\label{pf:thmK:L:towerA} \emph{$\alpha$-towers.} For a triple with $G$ of type $(b+2,a')$ and $q=0<p$ there are a
re-choice $G_2=G[V_2,V'_2,W_2]$ such that $(B,G_2,X,\omega,\omega')$ is a triple, a configuration $G_1$ of type $(b,a')$,
planes $H,H'$ such that $(G_2,G_1,H,H')$ is an $\alpha$-step configuration, and $\omega_1,\omega'_1$ such that
$(B,G_1,X,\omega_1,\omega'_1)$ is a triple.
\item\label{pf:thmK:L:towerB} \emph{$\beta$-towers.} The same for $G$ of type $(a,b+2)$ with $q\ge1$, with $G_1$ of type
$(a,b)$ and a $\beta$-step configuration.
\item\label{pf:thmK:L:facne} \emph{Non-vanishing of step factors.} For a triple,
$\thmKst(e_1,e_2,s;\mathfrak b,\mathfrak c;z)\ne0$ for all $e_1,e_2,s,\mathfrak b,\mathfrak c,z$, and
$\thmKLS{k}{\nu}\ne0$ for all integers $k\ge2$.
\item\label{pf:thmK:L:matchex} \emph{Partners.} For matched data: (a) every $\Phi\in\thmKPD(M)$ has some
$\Phi'\in\thmKSo(M')$ with $\Phi\approx\Phi'$; (b) every $\Phi'\in\thmKPD(M')$ has some $\Phi$ with $\Phi\approx\Phi'$; (c) if
$\Phi\approx\Phi'$ and $\Phi'\in\thmKSo(M')$ then $\Phi\in\thmKPD(M)$, $\Phi'\in\thmKPD(M')$ and $\Phi\sim\Phi'$.
\item\label{pf:thmK:L:levi} \emph{Levi elements.} For good $B$, $W\in\mathrm{Skew}_B(r,s)$ and $H'$ adapted to $(W,W_1)$,
every $h_1\in\Uni(W_1)(\A)$ admits a Levi element $m(1,h_1)$.
\item\label{pf:thmK:L:poldense} \emph{Density.} For a Weil pair, $\thmKPol(M)$ is dense in the Schwartz space.
\item\label{pf:thmK:L:weilconv} \emph{Weil's convergence.} For a Weil pair with $d+2\min(r,s)<2n$ or $\min(r,s)=0$: the theta
integral of every test function $\Phi$ converges and $\int_{\mathrm{dom}}|\theta_\omega(\Phi)(g,h)|\,d\mu(g)\le
C\thmKht(h)^k$ for some $C,k$; and for $\varphi\in\thmKS(\A_f^{n\times d})$ and every $h$, $f\mapsto\thmKI(f\otimes
\varphi)(h)$ is continuous on the Schwartz space.
\item\label{pf:thmK:L:socont} \emph{Continuity on $\thmKSo$.} For a Weil pair, finitely many $\varphi_j\in\thmKS(\A_f^{n\times
d})$ and Schwartz-convergent sequences $u_{k,j}\to f_j$ with $\sum_ju_{k,j}\otimes\varphi_j\in\thmKSo(M)$ for all $k$ and
$\sum_jf_j\otimes\varphi_j\in\thmKSo(M)$, one has $\thmKI(\sum_ju_{k,j}\otimes\varphi_j)(h)\to
\thmKI(\sum_jf_j\otimes\varphi_j)(h)$ for every $h$.
\item\label{pf:thmK:L:splitex} \emph{Splittings.} For a triple and $H'$ adapted to $(W,W_1)$: every $\thmKI(\Phi)$,
$\Phi\in\thmKPD(M)$, and every $\thmKI'(\Phi')$, $\Phi'\in\thmKPD(M')\cap\thmKSo(M')$, has exactly one splitting along $H'$.
\item\label{pf:thmK:L:ctbz} \emph{Evaluation part, big side.} For a $\beta$-step configuration $(G,G_1,H,H')$ and
$e\in\thmKEv(M,M_1;1,\iota_{H'};I_0(H'))$ there is $C\ne0$ such that for all triples on $G$ and on $G_1$ (with the same
$X$), all $\Phi\in\thmKPD(M)$, all $(f_\alpha,f_\beta)\in\thmKSpl(\thmKI(\Phi))$, all Levi elements $m=m(1,h_1)$ and all
$h$: $f_\beta(mh)=C\,\thmKI_1(e(\omega(1,h)\Phi))(h_1)$.
\item\label{pf:thmK:L:ctsz} \emph{Evaluation part, small side.} For an $\alpha$-step configuration and an evaluation map
$e'\in\thmKEv(M',M'_1;1,\iota_{H'};I_0(H'))$ there is $C\ne0$ with the following property: for all triples on $G$ and
$G_1$ and all $\Phi'\in\thmKPD(M')$ there is a function $f'_\alpha$ with $(f'_\alpha,0)\in\thmKSpl(\thmKI'(\Phi'))$ such
that, for all Levi elements $m=m(1,h_1)$ and all $h$, $f'_\alpha(mh)=C\,\thmKI'_1(e'(\omega'(1,h)\Phi'))(h_1)$.
\item\label{pf:thmK:L:ctbc} \emph{Cone part, big side.} For an $\alpha$-step configuration,
$e\in\thmKEv(M,M_1;\iota_H,\iota_{H'};I_A(H,H'))$, the stabiliser $K$ of $V$ and Haar measures $\mu_K,\mu$, there is $C\ne0$
such that for all triples on $G,G_1$, $\Phi\in\thmKPD(M)$, $(f_\alpha,f_\beta)\in\thmKSpl(\thmKI(\Phi))$, Levi elements
$m=m(1,h_1)$ and $h$: $f_\alpha(mh)=C\,\thmKI_1(\thmKcone_A(\omega(1,h)\Phi))(h_1)$.
\item\label{pf:thmK:L:ctsc} \emph{Cone part, small side.} For a $\beta$-step configuration,
$e'\in\thmKEv(M',M'_1;\iota_H,\iota_{H'};I_B(H,H'))$, the stabiliser $K$ of $V'$ and Haar measures, there is $C\ne0$ such
that for all triples on $G,G_1$, $\Phi'\in\thmKPD(M')\cap\thmKSo(M')$, $(f'_\alpha,f'_\beta)\in\thmKSpl(\thmKI'(\Phi'))$,
$m=m(1,h_1)$ and $h$: $f'_\beta(mh)=C\,\thmKI'_1(\thmKcone_B(\omega'(1,h)\Phi'))(h_1)$.
\item\label{pf:thmK:L:anall} \emph{Definite small side.} For matched data with $q=0$ and a polynomial constant $c$,
$\thmKI(\Phi)=c\,\thmKI'(\Phi')$ for \emph{all} $\Phi\approx\Phi'$.
\item\label{pf:thmK:L:polre} \emph{Change of Cartan data.} For matched data $(\omega,\omega',R,R_f)$ on $G$ and matched data
$(\omega,\omega',R_2,R_{f,2})$ on a re-choice $G_2=G[V_2,V'_2,W_2]$: if $G_2$ has a non-zero polynomial constant, so does
$G$.
\item\label{pf:thmK:L:nonzeroan} \emph{Non-vanishing.} For a triple, if $p=0$ there is $\Phi\in\thmKPD(M)$ with
$\thmKI(\Phi)\ne0$, and if $q=0$ there is $\Phi'\in\thmKPD(M')\cap\thmKSo(M')$ with $\thmKI'(\Phi')\ne0$.
\item\label{pf:thmK:L:polzero} \emph{The definite case.} For matched data with $p=q=0$ there is a non-zero polynomial
constant.
\item\label{pf:thmK:L:locA} \emph{Local factors of $\alpha$-steps.} There are $e_1,e_2,s\in\{\pm1\}$ such that for an
$\alpha$-step configuration, real maps $R,R_1$, $e\in\thmKEv(M,M_1;\iota_H,\iota_{H'};I_A)$,
$e'\in\thmKEv(M',M'_1;1,\iota_{H'};I_0)$, $K,\mu_K,\mu$ as in \ref{pf:thmK:L:ctbc}, there is $C\ne0$ such that for all
matched data on $G$ (with $R$) and $G_1$ (with $R_1$) satisfying $\mathcal B(p+q)$, and all $\Phi\approx\Phi'$ on $G$:
\[
\thmKcone_A\Phi\ \approx\ C\,\thmKst(e_1,e_2,s;\mathfrak b_H,\mathfrak b_{H'};1-k_A)\,\thmKLS{k_A}{\nu}\,e'\Phi'\quad
\text{on }G_1.
\]
\item\label{pf:thmK:L:locB} \emph{Local factors of $\beta$-steps.} There are $e_1,e_2,s$ such that for a $\beta$-step
configuration, $R,R_1$, $e\in\thmKEv(M,M_1;1,\iota_{H'};I_0)$, $e'\in\thmKEv(M',M'_1;\iota_H,\iota_{H'};I_B)$ and
$K,\mu_K,\mu$ as in \ref{pf:thmK:L:ctsc}, there is $C\ne0$ such that for matched data on $G,G_1$ with $\mathcal B(p+q)$ and
$\Phi\approx\Phi'$ with $\Phi'\in\thmKSo(M')$: $\thmKcone_B\Phi'\in\thmKSo(M'_1)$ and
\[
C\,\thmKst(e_1,e_2,s;\mathfrak b_H,\mathfrak b_{H'};-k_B)\,\thmKLS{k_B+1}{\nu}\,e\Phi\ \approx\ \thmKcone_B\Phi'\quad
\text{on }G_1.
\]
\item\label{pf:thmK:L:endA} \emph{$\alpha$-end step.} For matched data with $q=0<p$ and $\mathcal B(p+q)$, a plane $H'$
adapted to $(W,W_1)$ and $c\in\C$: if $f_\alpha(m)=c\,f'_\alpha(m)$ for all $\Phi\approx\Phi'$ with $\Phi'\in\thmKSo(M')$,
all $(f_\alpha,f_\beta)\in\thmKSpl(\thmKI(\Phi))$, $(f'_\alpha,f'_\beta)\in\thmKSpl(\thmKI'(\Phi'))$ and all Levi elements
$m=m(1,h_1)$, then $c$ is a polynomial constant.
\item\label{pf:thmK:L:endB} \emph{$\beta$-end step.} The same with $q\ge1$ and $f_\beta(m)=c\,f'_\beta(m)$.
\item\label{pf:thmK:L:kreal} \emph{Archimedean matching.} Every triple admits matching real maps.
\item\label{pf:thmK:L:kfin} \emph{Non-archimedean matching.} Every triple admits matching finite maps.
\end{enumerate}
Besides the leaves we use the statements \textup{(Dec)}: every triple satisfies the decay condition; and \textup{(Coi)}:
all matched data satisfy the coinvariant condition.
\formal{PositivityTheory.ThKG.Leaves}\formal{GlobalTheory.ThK.KDecay}\formal{GlobalTheory.ThK.KCoinv}
\end{definition}


% ===========================================================================
\subsection{From the leaves to Theorem K}\label{sec:pf:thmK:induction}

Throughout this subsection we assume the leaves. We use freely that a polynomially matched pair $\Phi\approx\Phi'$
consists of polynomial data, $\Phi\in\thmKPD(M)$ and $\Phi'\in\thmKPD(M')$ (each pure tensor in
\eqref{pf:thmK:eq:polmatch} lies in $\thmKPD$), and that $\thmKI(c\Phi)=c\,\thmKI(\Phi)$.

\begin{lemma}\label{pf:thmK:parts}
(i) If every triple satisfies the decay condition, all matched data satisfy the one-constant condition and the
coinvariant condition, and the step relations $\thmKSt_\beta(e_1,e_2,s)$, $\thmKSt_\alpha(e'_1,e'_2,s')$ hold for some signs
and all $F,E,\sigma$, then $\sfname{ThetaIntegralsAgree}$ holds. (ii) If every triple admits matching real maps and
matching finite maps, then $\sfname{LocalMatching}$ holds; in particular, under the leaves every triple admits matching
maps. (iii) For every triple and every test function $\Phi$ the theta integral of $\Phi$ converges; consequently $\thmKI$
is linear on test functions.
\formallabel{GlobalTheory.ThK.agree_of_parts}\formallabel{GlobalTheory.ThK.matching_of_parts}
\formallabel{PositivityTheory.ThKG.Leaves.matched}\formallabel{PositivityTheory.ThKG.Leaves.conv}
\end{lemma}

\begin{proof}
(i) and (ii) are the definitions (\cref{pf:thmK:def:K}) rearranged; for (ii) combine the real maps of
\ref{pf:thmK:L:kreal} with the finite maps of \ref{pf:thmK:L:kfin}. (iii) The Weil pair of $V$ satisfies the hypothesis of
\ref{pf:thmK:L:weilconv} by \cref{pf:thmK:pairs}. Linearity follows since the theta series of test functions are summable
and the integrals in \eqref{pf:thmK:eq:theta} converge.
\end{proof}

\subsubsection{The two key lemmas}

\begin{lemma}[Key lemma for $\beta$-steps]\label{pf:thmK:keyB}
There are $e_1,e_2,s\in\{\pm1\}$ with the following property. Let $(G,G_1,H,H')$ be a $\beta$-step configuration, $G$ of
type $(a,b+2)$, and let $R,R_1$ be real maps. Then there is $C\ne0$ such that for all matched data $(\omega,\omega',R_f)$ on
$G$ and $(\omega_1,\omega'_1,R_{f,1})$ on $G_1$ (with $R$, $R_1$ and the same $X$) satisfying $\mathcal B(p+q)$:
\begin{enumerate}
\item for every polynomial constant $c_1$ of $G_1$, every $\Phi\approx\Phi'$ with $\Phi'\in\thmKSo(M')$, all splittings
$(f_\alpha,f_\beta)\in\thmKSpl(\thmKI(\Phi))$, $(f'_\alpha,f'_\beta)\in\thmKSpl(\thmKI'(\Phi'))$ along $H'$ and every Levi
element $m=m(1,h_1)$,
\begin{equation}\label{pf:thmK:eq:keyB}
\bigl(C\,\thmKst(e_1,e_2,s;\mathfrak b_H,\mathfrak b_{H'};-k_B)\,\thmKLS{k_B+1}{\nu}\bigr)\,f_\beta(m)=c_1\,f'_\beta(m);
\end{equation}
\item there is $\Phi\in\thmKPD(M)$ such that every splitting $(f_\alpha,f_\beta)\in\thmKSpl(\thmKI(\Phi))$ has
$f_\beta(m)\ne0$ for some Levi element $m=m(1,h_1)$.
\end{enumerate}
\formallabel{PositivityTheory.ThKG.keyB_of}\formallabel{PositivityTheory.ThKG.Leaves.keyB}
\end{lemma}

\begin{proof}
Take $(e_1,e_2,s)$ from \ref{pf:thmK:L:locB}. Fix $G,G_1,H,H',R,R_1$. If there are no $X,\omega,\omega',\omega_1,
\omega'_1$ with triples on both $G$ and $G_1$, the statement is vacuous and we take $C=1$. Otherwise fix such data; they
are used only to produce the following objects, which do not depend on them. The Weil pair of $V$ and
\ref{pf:thmK:L:evzero} (for $W_1=W(G_1)$, the plane $H'$ and the model $M_1$ of $(V,W_1)$) give
$e\in\thmKEv(M,M_1;1,\iota_{H'};I_0(H'))$, and every such $e$ maps $\thmKPD(M)$ onto $\thmKPD(M_1)$. The Weil pair of $V'$
and \ref{pf:thmK:L:evcone} give $e'\in\thmKEv(M',M'_1;\iota_H,\iota_{H'};I_B(H,H'))$. \ref{pf:thmK:L:stab} gives the
compact stabiliser $K$ of $V'$ and Haar measures $\mu_K,\mu$. Let $C_0\ne0$, $C_1\ne0$, $C_2\ne0$ be the constants given by
\ref{pf:thmK:L:ctbz} (for $e$), \ref{pf:thmK:L:ctsc} (for $e',K,\mu_K,\mu$) and \ref{pf:thmK:L:locB} (for
$R,R_1,e,e',K,\mu_K,\mu$), and put $C=C_2C_1/C_0$. Write $\kappa_B=\thmKst(e_1,e_2,s;\mathfrak b_H,\mathfrak
b_{H'};-k_B)\,\thmKLS{k_B+1}{\nu}$.

(1) Let the data be as in (1). By \ref{pf:thmK:L:ctbz} and \ref{pf:thmK:L:ctsc} with $h=1$,
$f_\beta(m)=C_0\,\thmKI_1(e\Phi)(h_1)$ and $f'_\beta(m)=C_1\,\thmKI'_1(\thmKcone_B\Phi')(h_1)$. By \ref{pf:thmK:L:locB}, the
function $\thmKcone_B\Phi'$ lies in $\thmKSo(M'_1)$ and $C_2\kappa_B\,e\Phi\approx\thmKcone_B\Phi'$ on $G_1$; since $c_1$ is a
polynomial constant of $G_1$, we get $C_2\kappa_B\,\thmKI_1(e\Phi)(h_1)=c_1\,\thmKI'_1(\thmKcone_B\Phi')(h_1)$. Hence
\[
C\kappa_B\,f_\beta(m)=\frac{C_2C_1}{C_0}\kappa_B\,C_0\,\thmKI_1(e\Phi)(h_1)=C_1c_1\,\thmKI'_1(\thmKcone_B\Phi')(h_1)
=c_1f'_\beta(m).
\]
(2) Since $p(G_1)+q(G_1)<p+q$, $\mathcal B(p+q)$ applied to the matched data on $G_1$ gives $\Phi_1\in\thmKPD(M_1)$ and $h_1$
with $\thmKI_1(\Phi_1)(h_1)\ne0$. As $e$ maps $\thmKPD(M)$ onto $\thmKPD(M_1)$, $\Phi_1=e\Phi$ with $\Phi\in\thmKPD(M)$. By
\ref{pf:thmK:L:levi} there is a Levi element $m=m(1,h_1)$, and for every splitting $(f_\alpha,f_\beta)$ of $\thmKI(\Phi)$,
\ref{pf:thmK:L:ctbz} gives $f_\beta(m)=C_0\thmKI_1(\Phi_1)(h_1)\ne0$.
\end{proof}

\begin{lemma}[Key lemma for $\alpha$-steps]\label{pf:thmK:keyA}
There are $e_1,e_2,s\in\{\pm1\}$ with the following property. Let $(G,G_1,H,H')$ be an $\alpha$-step configuration, $G$
of type $(b+2,a')$, and let $R,R_1$ be real maps. Then there is $C\ne0$ such that for all matched data on $G$ and $G_1$
(with $R,R_1$, the same $X$) satisfying $\mathcal B(p+q)$, every polynomial constant $c_1$ of $G_1$, and
$\kappa=C\,\thmKst(e_1,e_2,s;\mathfrak b_H,\mathfrak b_{H'};1-k_A)\,\thmKLS{k_A}{\nu}\,c_1$:
\begin{enumerate}
\item for every $\Phi\approx\Phi'$ there is $f'_\alpha$ with $(f'_\alpha,0)\in\thmKSpl(\thmKI'(\Phi'))$ such that
$f_\alpha(m)=\kappa\,f'_\alpha(m)$ for every $(f_\alpha,f_\beta)\in\thmKSpl(\thmKI(\Phi))$ and every Levi element
$m=m(1,h_1)$;
\item there are $\Phi\approx\Phi'$, $f'_\alpha$ as in (1) and a Levi element $m=m(1,h_1)$ with $f'_\alpha(m)\ne0$.
\end{enumerate}
\formallabel{PositivityTheory.ThKG.keyA_of}\formallabel{PositivityTheory.ThKG.Leaves.keyA}
\end{lemma}

\begin{proof}
Take $(e_1,e_2,s)$ from \ref{pf:thmK:L:locA}, fix $G,G_1,H,H',R,R_1$, and as before take $C=1$ if there are no triples on
$G$ and $G_1$. Otherwise fix triples to construct: by the Weil pair of $V'$ and \ref{pf:thmK:L:evzero} (note
$V'(G_1)=V'(G)$), $e'\in\thmKEv(M',M'_1;1,\iota_{H'};I_0(H'))$, mapping $\thmKPD(M')$ onto $\thmKPD(M'_1)$; by the Weil pair of
$V$ and \ref{pf:thmK:L:evcone}, $e\in\thmKEv(M,M_1;\iota_H,\iota_{H'};I_A(H,H'))$; by \ref{pf:thmK:L:stab}, the stabiliser
$K$ of $V$ and $\mu_K$, $\mu$. Let $C_0,C_1,C_2$ be the non-zero constants of \ref{pf:thmK:L:ctsz} (for $e'$),
\ref{pf:thmK:L:ctbc} (for $e,K,\mu_K,\mu$) and \ref{pf:thmK:L:locA}, and $C=C_2C_1/C_0$. Write
$\kappa_A=\thmKst(e_1,e_2,s;\mathfrak b_H,\mathfrak b_{H'};1-k_A)\,\thmKLS{k_A}{\nu}$, so $\kappa=C\kappa_Ac_1$.

We first show: for every $\Phi\approx\Phi'$ there is $f'_\alpha$ as in (1) which moreover satisfies
$f'_\alpha(m)=C_0\,\thmKI'_1(e'\Phi')(h_1)$ for all Levi elements $m=m(1,h_1)$. Indeed \ref{pf:thmK:L:ctsz} gives
$f'_\alpha$ with $(f'_\alpha,0)\in\thmKSpl(\thmKI'(\Phi'))$ and, at $h=1$, the last formula. Let
$(f_\alpha,f_\beta)\in\thmKSpl(\thmKI(\Phi))$ and $m=m(1,h_1)$. By \ref{pf:thmK:L:ctbc}, $f_\alpha(m)=C_1\,
\thmKI_1(\thmKcone_A\Phi)(h_1)$. By \ref{pf:thmK:L:locA}, $\thmKcone_A\Phi\approx C_2\kappa_A\,e'\Phi'$ on $G_1$; since
$q(G_1)=0$, \ref{pf:thmK:L:anall} applies on $G_1$ to this pair (whose second member need not lie in $\thmKSo$) and gives
$\thmKI_1(\thmKcone_A\Phi)(h_1)=c_1C_2\kappa_A\,\thmKI'_1(e'\Phi')(h_1)$. Therefore
\[
f_\alpha(m)=C_1c_1C_2\kappa_A\,\thmKI'_1(e'\Phi')(h_1)=\frac{C_2C_1}{C_0}\kappa_Ac_1\,\bigl(C_0\,\thmKI'_1(e'\Phi')(h_1)\bigr)
=\kappa\,f'_\alpha(m).
\]
This proves (1). For (2), $\mathcal B(p+q)$ applied to $G_1$ gives $\Phi'_1\in\thmKPD(M'_1)\cap\thmKSo(M'_1)$ and $h_1$ with
$\thmKI'_1(\Phi'_1)(h_1)\ne0$; write $\Phi'_1=e'\Phi'$ with $\Phi'\in\thmKPD(M')$; \ref{pf:thmK:L:matchex}(b) gives $\Phi$ with
$\Phi\approx\Phi'$; take $f'_\alpha$ as just constructed and $m=m(1,h_1)$ from \ref{pf:thmK:L:levi}; then
$f'_\alpha(m)=C_0\thmKI'_1(\Phi'_1)(h_1)\ne0$.
\end{proof}

\subsubsection{The induction on $p+q$}

\begin{lemma}[Change of Fock space]\label{pf:thmK:nzmove}
Let $(B,G,X,\omega,\omega')$ and $(B,G_2,X,\omega,\omega')$ be triples, where $G_2=G[V_2,V'_2,W_2]$ is a re-choice. If
$\thmKI(\Phi)\ne0$ for some $\Phi\in\thmKPD(M_2)$, then $\thmKI(\Phi)\ne0$ for some $\Phi\in\thmKPD(M)$.
\formallabel{PositivityTheory.ThKG.Leaves.nz_move}
\end{lemma}

\begin{proof}
$G_2$ has the same models (hence the same vacua, measure, domain and theta integrals) as $G$, but other Cartan involutions,
hence another Fock space. Suppose $\thmKI(\Phi)=0$ for all $\Phi\in\thmKPD(M)$; we show $\thmKI(\Psi)=0$ for all
$\Psi\in\thmKPD(M_2)$, which is spanned by the $\thmKpure(f,\varphi)$ with $f\in\thmKPol(M_2)$. Fix $\varphi=(\varphi_v)$
with $\varphi_v\in\thmKS$ and $\varphi_v=\thmKvac_v$ for almost all $v$. As the vacua are characteristic functions of
$\OO_v^{a\times(a+a')}$ for almost all $v$, there is $g\in\thmKS(\A_f^{a\times(a+a')})$ with $\thmKpure(f,\varphi)=f\otimes g$
for all Schwartz $f$. By \ref{pf:thmK:L:weilconv} (applicable to the Weil pair of $V$), for every $h$ the map
$f\mapsto\thmKI(f\otimes g)(h)$ is continuous on the Schwartz space; it vanishes on $\thmKPol(M)$ by assumption, and
$\thmKPol(M)$ is dense (\ref{pf:thmK:L:poldense}). Hence it vanishes identically, in particular on $\thmKPol(M_2)$. By
linearity of $\thmKI$ on test functions, $\thmKI$ vanishes on $\thmKPD(M_2)$.
\end{proof}

\begin{proposition}[Existence of a polynomial constant]\label{pf:thmK:constof}
Matched data $(\omega,\omega',R,R_f)$ on $G$ satisfying $\mathcal B(p+q)$ have a non-zero polynomial constant.
\formallabel{PositivityTheory.ThKG.Leaves.const_of}
\end{proposition}

\begin{proof}
\emph{Case $p=q=0$:} this is \ref{pf:thmK:L:polzero}.

\emph{Case $q=0<p$.} Since $\lfloor(a-a')/2\rfloor>2(p+q)+1$, $a\ge2$; write $a=b+2$. \ref{pf:thmK:L:towerA} gives a re-choice
$G_2$ carrying the triple $(B,G_2,X,\omega,\omega')$, a triple on $G_1$ of type $(b,a')$, and planes such that
$(G_2,G_1,H,H')$ is an $\alpha$-step configuration. By \cref{pf:thmK:parts}(ii) there are matching maps $R_2,R_{f,2}$ on
$G_2$ and $R_1,R_{f,1}$ on $G_1$. By \ref{pf:thmK:L:polre} it suffices to find a non-zero polynomial constant of $G_2$ (with
$R_2,R_{f,2}$). Let $C$ and $(e_1,e_2,s)$ be given by \cref{pf:thmK:keyA} for $(G_2,G_1,H,H')$ and $R_2,R_1$. As
$p(G_1)+q(G_1)=p+q-1$, $\mathcal B(p+q)$ gives a non-zero polynomial constant $c_1$ of $G_1$. Put
$c=C\,\thmKst(e_1,e_2,s;\mathfrak b_H,\mathfrak b_{H'};1-k_A)\,\thmKLS{k_A}{\nu}\,c_1$. Since $p\ge1$ we have
$k_A=\lfloor(a-a')/2\rfloor>2(p+q)+1\ge2$, so $c\ne0$ by \ref{pf:thmK:L:facne}. We apply \ref{pf:thmK:L:endA} to the matched
data on $G_2$, the plane $H'$ and $c$. Let $\Phi\approx\Phi'$ with $\Phi'\in\thmKSo(M')$, splittings $(f_\alpha,f_\beta)$ of
$\thmKI(\Phi)$ and $(f'_\alpha,f'_\beta)$ of $\thmKI'(\Phi')$, and $m=m(1,h_1)$. \Cref{pf:thmK:keyA}(1) (where $\kappa=c$) gives
$f''_\alpha$ with $(f''_\alpha,0)\in\thmKSpl(\thmKI'(\Phi'))$ and $f_\alpha(m)=c\,f''_\alpha(m)$. Since $\Phi'\in\thmKPD(M')
\cap\thmKSo(M')$, the uniqueness in \ref{pf:thmK:L:splitex} gives $(f'_\alpha,f'_\beta)=(f''_\alpha,0)$, so
$f_\alpha(m)=c\,f'_\alpha(m)$. Hence $c$ is a polynomial constant of $G_2$.

\emph{Case $q\ge1$.} Since $a'>2(p+q)+1\ge3$, write $a'=b+2$. \ref{pf:thmK:L:towerB} gives $G_2$ and $G_1$ of type $(a,b)$
forming a $\beta$-step configuration $(G_2,G_1,H,H')$, with triples; take matching maps as before and reduce to $G_2$ by
\ref{pf:thmK:L:polre}. Let $C,(e_1,e_2,s)$ be given by \cref{pf:thmK:keyB}, let $c_1\ne0$ be a polynomial constant of $G_1$
($\mathcal B(p+q)$), and $\kappa=C\,\thmKst(e_1,e_2,s;\mathfrak b_H,\mathfrak b_{H'};-k_B)\,\thmKLS{k_B+1}{\nu}$, which is
non-zero by \ref{pf:thmK:L:facne} as $k_B+1\ge2$. Put $c=c_1/\kappa\ne0$. In the situation of \ref{pf:thmK:L:endB},
\eqref{pf:thmK:eq:keyB} gives $\kappa f_\beta(m)=c_1f'_\beta(m)$, i.e.\ $f_\beta(m)=c\,f'_\beta(m)$; so
\ref{pf:thmK:L:endB} shows that $c$ is a polynomial constant of $G_2$.
\end{proof}

\begin{lemma}[Non-vanishing on the big side]\label{pf:thmK:nzV}
For a triple with $q\ge1$ satisfying $\mathcal B(p+q)$ there is $\Phi\in\thmKPD(M)$ with $\thmKI(\Phi)\ne0$.
\formallabel{PositivityTheory.ThKG.Leaves.nzV}
\end{lemma}

\begin{proof}
As in the last case above, \ref{pf:thmK:L:towerB} gives a re-choice $G_2$ and $G_1$ forming a $\beta$-step configuration,
with triples, and \cref{pf:thmK:parts}(ii) gives matching maps on both. By \cref{pf:thmK:keyB}(2) there is
$\Phi\in\thmKPD(M_2)$ such that every splitting $(f_\alpha,f_\beta)$ of $\thmKI(\Phi)$ has $f_\beta(m)\ne0$ for some Levi
element $m$. If $\thmKI(\Phi)=0$ then $(0,0)\in\thmKSpl(\thmKI(\Phi))$, a contradiction. So $\thmKI(\Phi)\ne0$, and
\cref{pf:thmK:nzmove} gives the assertion.
\end{proof}

\begin{theorem}[Non-degeneracy of all matched data]\label{pf:thmK:claim}
$\mathcal B(n)$ holds for every $n$; hence all matched data are non-degenerate.
\formallabel{PositivityTheory.ThKG.Leaves.claim_of}\formallabel{PositivityTheory.ThKG.Leaves.below}
\formallabel{PositivityTheory.ThKG.Leaves.claim}
\end{theorem}

\begin{proof}
We first show: matched data satisfying $\mathcal B(p+q)$ are non-degenerate. \Cref{pf:thmK:constof} gives a non-zero
polynomial constant $c$. If $q=0$, \ref{pf:thmK:L:nonzeroan} gives $\Phi'\in\thmKPD(M')\cap\thmKSo(M')$ and $x$ with
$\thmKI'(\Phi')(x)\ne0$; \ref{pf:thmK:L:matchex}(b) gives $\Phi\approx\Phi'$; then $\Phi\in\thmKPD(M)$ and
$\thmKI(\Phi)(x)=c\,\thmKI'(\Phi')(x)\ne0$. If $q\ge1$, \cref{pf:thmK:nzV} gives $\Phi\in\thmKPD(M)$ with $\thmKI(\Phi)\ne0$;
\ref{pf:thmK:L:matchex}(a) gives $\Phi'\in\thmKSo(M')$ with $\Phi\approx\Phi'$; then $\Phi'\in\thmKPD(M')$, and
$\thmKI'(\Phi')\ne0$ since otherwise $\thmKI(\Phi)=c\,\thmKI'(\Phi')=0$.

Now $\mathcal B(0)$ is vacuous, and if $\mathcal B(n)$ holds then matched data with $p+q<n+1$ satisfy $\mathcal B(p+q)$ by
monotonicity, hence are non-degenerate; so $\mathcal B(n+1)$ holds. Given matched data, apply $\mathcal B(p+q+1)$.
\end{proof}

\subsubsection{Constants and steps}

\begin{lemma}\label{pf:thmK:polconst}
For matched data, a number is a constant if and only if it is a polynomial constant.
\formallabel{PositivityTheory.ThKG.Leaves.polConst_of}\formallabel{PositivityTheory.ThKG.Leaves.isConst_of}
\end{lemma}

\begin{proof}
If $c$ is a constant and $\Phi\approx\Phi'$ with $\Phi'\in\thmKSo(M')$, then $\Phi\sim\Phi'$ by \ref{pf:thmK:L:matchex}(c),
so $\thmKI(\Phi)=c\,\thmKI'(\Phi')$. Conversely let $c$ be a polynomial constant and $\Phi\sim\Phi'$, with
$\Phi=\sum_if_i\otimes g_i$, $\Phi'=\sum_jf'_j\otimes g'_j$ and approximating sequences $s_{k,i}\to f_i$, $s'_{k,j}\to f'_j$ as
in \cref{pf:thmK:def:match}. Fix $h$. Since $\thmKI$ is linear on test functions and $f\mapsto\thmKI(f\otimes g_i)(h)$ is
continuous (\ref{pf:thmK:L:weilconv}), $\thmKI(\sum_is_{k,i}\otimes g_i)(h)\to\thmKI(\Phi)(h)$. By \ref{pf:thmK:L:socont} (for
the Weil pair of $V'$), $\thmKI'(\sum_js'_{k,j}\otimes g'_j)(h)\to\thmKI'(\Phi')(h)$. Since the $k$-th approximants satisfy
$\Phi_k\approx\Phi'_k$ with $\Phi'_k\in\thmKSo(M')$, $\thmKI(\Phi_k)(h)=c\,\thmKI'(\Phi'_k)(h)$, and uniqueness of limits
gives $\thmKI(\Phi)(h)=c\,\thmKI'(\Phi')(h)$.
\end{proof}

\begin{proposition}\label{pf:thmK:kconst}
All matched data satisfy the one-constant condition.
\formallabel{PositivityTheory.ThKG.Leaves.kConst}
\end{proposition}

\begin{proof}
By \cref{pf:thmK:claim} there are a non-zero polynomial constant $c$ and $\Phi'\in\thmKPD(M')\cap\thmKSo(M')$, $x$ with
$\thmKI'(\Phi')(x)\ne0$; \ref{pf:thmK:L:matchex}(b) gives $\Phi\approx\Phi'$. If $c'$ is a constant, it is a polynomial
constant (\cref{pf:thmK:polconst}), so $c'\thmKI'(\Phi')(x)=\thmKI(\Phi)(x)=c\,\thmKI'(\Phi')(x)$ and $c'=c$. By
\cref{pf:thmK:polconst}, $c$ is a constant. Hence there is exactly one constant, it is non-zero, and $\thmKI'(\Phi')\ne0$
with $\Phi'\in\thmKSo(M')$.
\end{proof}

\begin{proposition}\label{pf:thmK:ksteps}
The $\beta$-step relation holds with the signs of \cref{pf:thmK:keyB}, and the $\alpha$-step relation holds with the signs
of \cref{pf:thmK:keyA}, for all $F,E,\sigma$.
\formallabel{PositivityTheory.ThKG.Leaves.kBeta}\formallabel{PositivityTheory.ThKG.Leaves.kAlpha}
\end{proposition}

\begin{proof}
\emph{The $\beta$-step relation.} Its hypotheses on $(G,G_1,H,H')$ say exactly that $(G,G_1,H,H')$ is a $\beta$-step
configuration. Take $C$ from \cref{pf:thmK:keyB} (for $R,R_1$); by \cref{pf:thmK:claim}, $\mathcal B(p+q)$ holds. Let
$c,c_1$ be constants of $G,G_1$. \Cref{pf:thmK:keyB}(2) gives $\Phi\in\thmKPD(M)$; \ref{pf:thmK:L:matchex}(a) gives
$\Phi'\in\thmKSo(M')$ with $\Phi\approx\Phi'$, and \ref{pf:thmK:L:splitex} a splitting $(f'_\alpha,f'_\beta)$ of
$\thmKI'(\Phi')$. As $c$ is a polynomial constant, $\thmKI(\Phi)=c\,\thmKI'(\Phi')$, so $(cf'_\alpha,cf'_\beta)\in
\thmKSpl(\thmKI(\Phi))$, and \cref{pf:thmK:keyB}(2) gives $m$ with $cf'_\beta(m)\ne0$. By \eqref{pf:thmK:eq:keyB} (with the
polynomial constant $c_1$), $C\kappa_B\,c\,f'_\beta(m)=c_1f'_\beta(m)$, where $\kappa_B=\thmKst(e_1,e_2,s;\mathfrak
b_H,\mathfrak b_{H'};-k_B)\thmKLS{k_B+1}{\nu}$. Cancelling $f'_\beta(m)$ gives $C\,\thmKst(\dots;-k_B)\,c\,\thmKLS{k_B+1}{\nu}
=c_1$. By \ref{pf:thmK:L:facne}, $C\,\thmKst(\dots;-k_B)\ne0$, and $u=(C\,\thmKst(\dots;-k_B))^{-1}$ satisfies both equations
\eqref{pf:thmK:eq:stepbeta}.

\emph{The $\alpha$-step relation.} Its hypotheses say that $(G,G_1,H,H')$ is an $\alpha$-step configuration. Take $C$ from
\cref{pf:thmK:keyA}; let $c,c_1$ be constants and $\kappa=C\,\thmKst(e_1,e_2,s;\mathfrak b_H,\mathfrak
b_{H'};1-k_A)\,\thmKLS{k_A}{\nu}\,c_1$. \Cref{pf:thmK:keyA}(2) gives $\Phi\approx\Phi'$, $f'_\alpha$ with
$(f'_\alpha,0)\in\thmKSpl(\thmKI'(\Phi'))$ and the property (1), and $m$ with $f'_\alpha(m)\ne0$. Since $q(G)=0$,
\ref{pf:thmK:L:anall} gives $\thmKI(\Phi)=c\,\thmKI'(\Phi')$, so $(cf'_\alpha,0)\in\thmKSpl(\thmKI(\Phi))$, and (1) gives
$cf'_\alpha(m)=\kappa f'_\alpha(m)$; hence $c=\kappa$, which is \eqref{pf:thmK:eq:stepalpha}.
\end{proof}

\begin{theorem}[Theorem K from the leaves]\label{pf:thmK:agree}
The leaves, \textup{(Dec)} and \textup{(Coi)} imply $\sfname{ThetaIntegralsAgree}$ and $\sfname{LocalMatching}$.
\formallabel{PositivityTheory.ThKG.agree_of_leaves}
\end{theorem}

\begin{proof}
Combine \cref{pf:thmK:parts} with \cref{pf:thmK:kconst,pf:thmK:ksteps} and the leaves \ref{pf:thmK:L:kreal},
\ref{pf:thmK:L:kfin}.
\end{proof}


% ===========================================================================
\subsection{Proofs of the leaves}\label{sec:pf:thmK:leafproofs}

In this subsection we prove the twenty-six leaves and the statements \textup{(Dec)}, \textup{(Coi)}. Most leaves are
deduced from a small number of further propositions of the development; when such a proposition is not proved here, we
state what it asserts and cite its formal proof. In \cref{sec:pf:thmK:final} we record which hypotheses are used where.

\subsubsection{Reduction theory and analytic estimates}\label{sec:pf:thmK:est}
The analytic leaves rest on reduction theory for unitary groups, in the following form.

\begin{littheorem}[Reduction theory; \cite{Borel1969,MoeglinWaldspurger1995}]\label{pf:thmK:siegel}
Let $E/F$ be CM, $J\in M_n(E)$ a non-degenerate $\pm$hermitian matrix and $(X,Y)$ a Witt frame of $J$. There are a
compact subgroup $K$ and a compact subset $\Omega$ of the minimal parabolic subgroup, and $c>0$, such that
$\Uni(J)(\A)=P\,A\,K$ (Iwasawa decomposition), the Siegel set $\mathfrak S=\Omega A_cK$ is closed and $\Uni(J)(\A)=\Gamma_J
\mathfrak S$, $\mathfrak S$ has finite Haar measure with the usual integration bound over $A_c$, heights on $\mathfrak S$
are bounded by powers of heights of their $\Gamma_J$-translates, and smooth cuspidal functions of moderate growth decay
rapidly on $\mathfrak S$.
\formal{GlobalTheory.Published.SiegelSetUnitary}
\end{littheorem}

In the development this statement is \lref{PositivityTheory.ThKR.siegel_of_core}, which deduces it from a core
statement on Siegel sets of the adelic unitary group (\lref{PositivityTheory.Prw.siegelCore_holds}).

\begin{proposition}\label{pf:thmK:analytic}
\textup{(a)} For a Weil pair and a test function $\Phi$, the function $(g,h)\mapsto\theta_\omega(\Phi)(g,h)$ is
continuous, left $\Gamma_X\times\Gamma_W$-invariant and bounded by $C(\thmKht(g)\thmKht(h))^k$, and
$f\mapsto\theta_\omega(f\otimes\varphi)(g,h)$ is continuous on Schwartz functions. \textup{(b)} The leaves
\ref{pf:thmK:L:weilconv}, \ref{pf:thmK:L:socont} and the statement \textup{(Dec)} hold. \textup{(c)} The leaves
\ref{pf:thmK:L:evzero}, \ref{pf:thmK:L:evcone}, \ref{pf:thmK:L:stab}, \ref{pf:thmK:L:levi} and \ref{pf:thmK:L:poldense} hold.
\end{proposition}

\begin{proof}
(a) is \lref{PositivityTheory.ThKF0.thetaBasic}, which proves these properties of theta functions. (b) The leaves
\ref{pf:thmK:L:weilconv}, \ref{pf:thmK:L:socont} and \textup{(Dec)} are \lref{PositivityTheory.ThKEst.weilConv},
\lref{PositivityTheory.ThKEst.soCont} and \lref{PositivityTheory.ThKEst.kDecay}, each deduced from (a) and
\cref{pf:thmK:siegel}: Weil's convergence argument \cite{Weil1965} on Siegel sets in the sufficient form ``$d+2\min(r,s)<2n$
or $\min(r,s)=0$'' ($\min(r,s)$ bounds the Witt index of $X$), and, for functions in $\thmKSo$, bounds for lattice sums
restricted to matrices of full rank, which give rapid decay in $g$. (c) \Lref{PositivityTheory.ThKF0.evaluation_at_zero_and_cone}
constructs the evaluation maps of \ref{pf:thmK:L:evzero} and \ref{pf:thmK:L:evcone} and shows that those along $W$ map
polynomial data onto polynomial data; \lref{PositivityTheory.ThKG.Tow.stabEx} shows that the stabiliser of the lattice
and of the Cartan involutions is compact and that Haar measures exist (\ref{pf:thmK:L:stab});
\lref{PositivityTheory.ThKG.Tow.leviEx} constructs the Levi elements $m(1,h_1)$ (\ref{pf:thmK:L:levi}); and
\lref{GlobalTheory.ThK.Dbl.polDense} proves the density of polynomials times Gaussians in the Schwartz space
(\ref{pf:thmK:L:poldense}).
\end{proof}

\subsubsection{Constant terms}\label{sec:pf:thmK:ct}
The constant-term leaves are deduced from three general formulas for the constant term of a theta integral along a
plane $H'$ of $W$. In the classical computation, the theta series is expanded according to the pairing of the variable
with the isotropic vector $x_{H'}$; the terms where it vanishes give an \emph{evaluation} of the theta integral of the
smaller pair $(X,W_1)$, and the other terms, after unfolding over the stabiliser of an isotropic line of $X$, give a
\emph{cone integral} \eqref{pf:thmK:eq:cone}; the two contributions have different laws under the Levi torus.

\begin{proposition}[General constant-term formulas]\label{pf:thmK:gen}
Let $(X,W)$ have dimensions $(n,d)$, $H'$ be adapted to $(W,W_1)$, $M_1$ a good model for $(X,W_1)$ and
$e_0\in\thmKEv(M,M_1;1,\iota_{H'};I_0(H'))$.
\begin{enumerate}
\item If $\min(r,s)=0$, there is $C\ne0$ such that for all Weil pairs $(M,\omega)$ of signature $(r,s;r',s')$ and
$(M_1,\omega_1)$ of signature $(r,s;r_1,s_1)$ with the same characters $(\chi_X,\chi_W)$ and every test function $\Phi$
there is $A_0$ with $\thmKCT_{H'}(\thmKI(\Phi))=A_0$, $A_0$ of law $(\chi_X,n/2)$, and $A_0(mh)=C\,\thmKI_{M_1,\omega_1}
(e_0(\omega(1,h)\Phi))(h_1)$ for all Levi elements $m=m(1,h_1)$ and all $h$.
\item Let moreover $H$ be adapted to $(X,X_1)$, $M_0$ a good model for $(X_1,W_1)$,
$e\in\thmKEv(M,M_0;\iota_H,\iota_{H'};I_A(H,H'))$, $z\in\N$ with $2n=d+2z$ and $\min(r,s)+2\le z$, and $K$ a compact
subgroup with $\Uni(X)(\A)=P_HK$, where $P_H$ is the stabiliser of the line $\thmKAE x_H$, with Haar measures $\mu_K,\mu$.
There are $C_0,C_1\ne0$ such that for all Weil pairs on $M$ (signature $(r+1,s+1;r'+1,s'+1)$), $M_1$ (signature
$(r+1,s+1;r',s')$) and $M_0$ (signature $(r,s;r',s')$) with the same characters, and all test functions $\Phi$ whose
$\omega(K,1)$-translates span a finite-dimensional space, $\thmKCT_{H'}(\thmKI(\Phi))=A_0+A_1$ with $A_0$ of law
$(\chi_X,n/2)$, $A_1$ of law $(\chi_X/\chi_W,(d-n)/2)$, $A_1(mh)=C_1\thmKI_{M_0,\omega_0}(\thmKcone(\omega(1,h)\Phi))(h_1)$
for the cone integral with $u=u_A$, character $\chi_W^{-1}$ and exponent $z$, and $A_0(mh)=C_0\thmKI_{M_1,\omega_1}(e_0(
\omega(1,h)\Phi))(h_1)$.
\item The same with $I_B(H,H')$, $d+2=2n+2z$, $z\ge2$, $\Phi\in\thmKSo(M)$, the cone integral with $u=u_B$ and character
$\chi_W$, and the formula for $A_0(mh)$ asserted only when $e_0(\omega(1,h)\Phi)\in\thmKSo(M_1)$.
\end{enumerate}
These are \lref{PositivityTheory.ThKCT.genAniso_of}, \lref{PositivityTheory.ThKCT.genBig_of} and
\lref{PositivityTheory.ThKCT.genSmall_of}, applied to the hypotheses collected in the proof of \cref{pf:thmK:closesph}.
\formal{PositivityTheory.ThKCT.GenAniso}\formal{PositivityTheory.ThKCT.GenBig}\formal{PositivityTheory.ThKCT.GenSmall}
\end{proposition}

\begin{lemma}[The evaluation part]\label{pf:thmK:evalpart}
Let $(M,\omega)$ be a Weil pair for $(X,W)$ of signature $(r,s;r'+1,s'+1)$ with $n\ge3$, $H'$ adapted to $(W,W_1)$, $(M_1,
\omega_1)$ a Weil pair for $(X,W_1)$ of signature $(r,s;r',s')$ with the same characters, and
$e_0\in\thmKEv(M,M_1;1,\iota_{H'};I_0(H'))$. Assume either $2n=d+2z$ with $\min(r,s)+1\le z$, or $d+2=2n+2z$ with $z\ge2$.
There is $C_0\ne0$, independent of the Weil pairs, such that for every $\Phi\in\thmKPD(M)$ (in the second case,
$\Phi\in\thmKPD(M)\cap\thmKSo(M)$) the constant term of $\thmKI(\Phi)$ along $H'$ is $A_0+A_1$ with $A_0$ of law
$(\chi_X,n/2)$, $A_1$ of law $(\chi_X/\chi_W,(d-n)/2)$, and $A_0(mh)=C_0\thmKI_{M_1,\omega_1}(e_0(\omega(1,h)\Phi))(h_1)$ for
all Levi elements $m=m(1,h_1)$ and all $h$ (in the second case, whenever $e_0(\omega(1,h)\Phi)\in\thmKSo(M_1)$).
\formallabel{PositivityTheory.ThKCT.theta_kernel_general_identity}\formallabel{PositivityTheory.ThKCT.anyBig}
\formallabel{PositivityTheory.ThKCT.anySmall}
\end{lemma}

\begin{proof}
If $\min(r,s)=0$ this is \cref{pf:thmK:gen}(1), with $A_1=0$. Otherwise $r,s\ge1$; if there are no Weil pairs as in the
statement, take $C_0=1$. Otherwise, re-choosing the Cartan involutions of $X$ (which changes neither the model nor the
theta integrals), we find a plane $H$ adapted to $(X,X_1)$ at $(0,1)$ with $X_1$ of signature $(r-1,s-1)$
(\lref{PositivityTheory.ThKG.Tow.planeHerm}: after re-choosing the Cartan data, a hermitian space of signature $(r,s)$
with $r,s\ge1$ contains such a plane), a good model $M_0$ for $(X_1,W_1)$ carrying Weil representations for the given
characters (\lref{PositivityTheory.ThKG.Tow.pairEx_of}: good models exist and carry Weil representations), evaluation
maps of types $I_A$ and $I_B$ (\ref{pf:thmK:L:evcone}), and the stabiliser $K$ of $X$ (\ref{pf:thmK:L:stab}), which satisfies
$\Uni(X)(\A)=P_HK$ (\lref{PositivityTheory.ThKCT.iwaEx_of}: the Iwasawa decomposition with respect to the stabiliser of
an isotropic line). Elements of $\thmKPD(M)$ are test functions whose $\omega(K,1)$-translates span a finite-dimensional
space (\lref{PositivityTheory.ThKCT.kFinData_of}). Hence \cref{pf:thmK:gen}(2), resp.\ (3), applies with
$(r,s)$ replaced by $(r-1,s-1)$, and its constant $C_0$ does the job.
\end{proof}

\begin{lemma}[Uniqueness of splittings]\label{pf:thmK:splituniq}
Let $a'<a$, let $H'$ be adapted to $(W,W_1)$ with $W$ skew-hermitian for a good base datum, and let
$(f_\alpha,f_\beta),(f'_\alpha,f'_\beta)\in\thmKSpl_{H'}(f)$. Then $f_\alpha=f'_\alpha$ and $f_\beta=f'_\beta$.
\formallabel{PositivityTheory.ThKCT.split2_unique}\formallabel{PositivityTheory.ThKCT.split_uniq}
\end{lemma}

\begin{proof}
By \lref{PositivityTheory.ThKCT.unipEx_of}, $N(H')$ carries a Haar measure with a fundamental domain for
$N(H')\cap\Gamma_W$ of volume one, and there is a Levi element $m=m(t,1)$ with $|t|_{\thmKAE}>1$. Evaluating the
definition of $\thmKCT_{H'}$ with this measure and domain gives $f_\alpha+f_\beta=f'_\alpha+f'_\beta$, so
$D:=f_\alpha-f'_\alpha=f'_\beta-f_\beta$. By the laws, $D(mh)=\alpha(t)|t|^{a'/2}D(h)$ and
$D(mh)=\beta(t)|t|^{a/2}D(h)$ for all $h$. If $D(h)\ne0$ for some $h$, then $\alpha(t)|t|^{a'/2}=\beta(t)|t|^{a/2}$;
taking absolute values (the characters are unitary) gives $|t|^{a'/2}=|t|^{a/2}$, which is impossible since $|t|>1$ and
$a'<a$. Hence $D=0$.
\end{proof}

\begin{proposition}[Constant-term leaves]\label{pf:thmK:ctleaves}
The leaves \ref{pf:thmK:L:splitex}--\ref{pf:thmK:L:ctsc} hold.
\formallabel{PositivityTheory.ThKCT.splitEx_of}\formallabel{PositivityTheory.ThKCT.ctBigZero_of}
\formallabel{PositivityTheory.ThKCT.ctSmallZero_of}\formallabel{PositivityTheory.ThKCT.ctBigCone_of}
\formallabel{PositivityTheory.ThKCT.ctSmallCone_of}\formallabel{PositivityTheory.ThKCT.ct_all}
\end{proposition}

\begin{proof}
Let $(B,G,X,\omega,\omega')$ be a triple. For the Weil pair of $V$ ($n=a$, $d=a+a'$, characters $(\beta,\beta/\alpha)$) one
has $\chi_X/\chi_W=\alpha$ and $(d-n)/2=a'/2$; for the Weil pair of $V'$ ($n=a'$, characters $(\alpha,\alpha/\beta)$) one has
$\chi_X/\chi_W=\beta$ and $(d-n)/2=a/2$. The numerical hypotheses of \cref{pf:thmK:evalpart,pf:thmK:gen} follow from the
inequalities of \cref{pf:thmK:def:geo}: an adapted plane lowers both signature indices of $W$ by one, so they are both
positive (\lref{PositivityTheory.ThKCT.low_of}), whence $p+q\ge1$ and $a>a'>2(p+q)+1\ge3$; for $V$,
$z=\lfloor(a-a')/2\rfloor>2(p+q)+1\ge\min(a-p,p)+1$; for $V'$, $z=\lfloor(a-a')/2\rfloor+1\ge2$; and $W_1$ has the signature
required of the second Weil pair. Models of $(V,W_1)$, $(V',W_1)$ with Weil
representations exist by \lref{PositivityTheory.ThKG.Tow.pairEx_of}, and evaluation maps along $W$ by
\ref{pf:thmK:L:evzero}. In each case below the constant is fixed before the triple, and if no triple exists the statement
is vacuous with constant $1$.

\ref{pf:thmK:L:splitex}: for $\Phi\in\thmKPD(M)$, \cref{pf:thmK:evalpart} for the Weil pair of $V$ shows that $(A_1,A_0)$ is a
splitting of $\thmKI(\Phi)$; for $\Phi'\in\thmKPD(M')\cap\thmKSo(M')$, applied to the Weil pair of $V'$, that $(A_0,A_1)$ is a
splitting of $\thmKI'(\Phi')$. Uniqueness is \cref{pf:thmK:splituniq}.

\ref{pf:thmK:L:ctbz} ($\beta$-step configuration, Weil pair of $V$ with $W_1=W(G_1)$ and $M_1=M(G_1)$): by uniqueness the
$\beta$-part of a splitting of $\thmKI(\Phi)$ is $A_0$, whence the formula.
\ref{pf:thmK:L:ctsz} ($\alpha$-step configuration, $q=0$, so $V'$ is definite): \cref{pf:thmK:gen}(1) for the Weil pair of $V'$
gives $(A_0,0)\in\thmKSpl(\thmKI'(\Phi'))$ and the formula for $A_0$.
\ref{pf:thmK:L:ctbc} ($\alpha$-step configuration, $H$ adapted to $(V,V(G_1))$): \cref{pf:thmK:gen}(2) for the Weil pair of
$V$, with $\chi_W^{-1}=\alpha/\beta=\nu$ and $z=k_A$, together with uniqueness, gives
$f_\alpha(mh)=A_1(mh)=C_1\thmKI_1(\thmKcone_A(\omega(1,h)\Phi))(h_1)$.
\ref{pf:thmK:L:ctsc} ($\beta$-step configuration, $H$ adapted to $(V',V'(G_1))$): \cref{pf:thmK:gen}(3) for the Weil pair of
$V'$, with $\chi_W=\nu$ and $z=k_B+1$, gives $f'_\beta(mh)=A_1(mh)=C_1\thmKI'_1(\thmKcone_B(\omega'(1,h)\Phi'))(h_1)$.
\end{proof}

\begin{remark}\label{pf:thmK:rem:ctcone}
The development contains a second proof of the statement \ref{pf:thmK:L:ctbc}, from a different collection of hypotheses,
namely \lref{GlobalTheory.ThK.ctBigCone} (which proves the same cone formula for the $\alpha$-part of the constant term
of $\thmKI(\Phi)$). It is not used in the proof of Theorem~K, which uses \cref{pf:thmK:ctleaves}.
\end{remark}

\subsubsection{Towers}\label{sec:pf:thmK:towers}

\begin{proposition}[Towers]\label{pf:thmK:towers}
The leaves \ref{pf:thmK:L:towerA} and \ref{pf:thmK:L:towerB} hold.
\formallabel{PositivityTheory.ThKG.Tow.towerA_of}\formallabel{PositivityTheory.ThKG.Tow.towerB_of}
\end{proposition}

\begin{proof}
We use, besides \cref{pf:thmK:weilex}, the existence of adapted planes (\lref{PositivityTheory.ThKG.Tow.planeEx}: for
good $B$, a hermitian or skew-hermitian space $Y$ of dimension $n_1+2$, $n_1\ge1$, of signature $(r,s)$ with $r,s\ge1$, and
positions $i_0<i_1<n_1+2$, there are a re-choice $Y[Y']$ of the same signature, a space $Y_1$ of signature $(r-1,s-1)$ and a
plane adapted to $(Y[Y'],Y_1)$ at $(i_0,i_1)$) and of good models (\lref{PositivityTheory.ThKG.Tow.modelEx}: every pair
of a hermitian and a skew-hermitian space has a good model).

\ref{pf:thmK:L:towerB}: let the triple be given with $G$ of type $(a,b+2)$ and $q\ge1$. Since $a'=b+2>2(p+q)+1$, $b\ge1$.
Apply the existence of planes to $V'\in\mathrm{Herm}(q,b+2-q)$ at $(0,b+1)$ and to $W\in\mathrm{Skew}(a+b+2-(p+q),p+q)$ at
$(a,a+b+1)$; this gives $V'_2,V'_1,H$ and $W_2,W_1,H'$. The re-choice $G_2=G[V,V'_2,W_2]$ is still good and carries the triple
(the models and $\omega,\omega'$ are unchanged). Put $G_1=(p,q-1,V,V'_1,W_1,M_1,M'_1)$ with good models $M_1,M'_1$; it is good:
$a+b$ is even, $a+b<a+b+2<m_\alpha-m_\beta$, the inequalities $2(p+q-1)+1<\lfloor(a-b)/2\rfloor$ and $2(p+q-1)+1<b$ follow from
those for $G$, and the signatures are those of $V$, $V'_1\in\mathrm{Herm}(q-1,b-(q-1))$ and
$W_1\in\mathrm{Skew}(a+b-(p+q-1),p+q-1)$. The character datum is good for $(B,a,b)$ because
$\alpha|_{\A^\times}=\eta^{b+2}=\eta^b$. \Cref{pf:thmK:weilex} gives Weil representations $\omega_1$ (characters
$(\beta,\beta/\alpha)$) and $\omega'_1$ (characters $(\alpha,\alpha/\beta)$), and $(G_2,G_1,H,H')$ is a $\beta$-step
configuration by construction. \ref{pf:thmK:L:towerA} is proved in the same way, removing the plane from
$V\in\mathrm{Herm}(b+2-p,p)$ and from $W$ at $(0,b+1)$, with $G_1=(p-1,0,V_1,V',W_1,M_1,M'_1)$ and
$\beta|_{\A^\times}=\eta^b$.
\end{proof}

\subsubsection{Non-vanishing of the step factors}\label{sec:pf:thmK:facne}

\begin{lemma}[Gauss sums]\label{pf:thmK:gauss}
Let $v$ be a finite place, $\psi_v$ unramified ($\psi_v(x\OO_v)=1$ iff $x\in\OO_v$) and $\nu_v$ a continuous character of
$F_v^\times$ which is non-trivial on $\OO_v^\times$. Then $|g(\psi_v,\nu_v)|^2=\Nm v^{c}$, where $c=c(\nu_v)\ge1$; in
particular $g(\psi_v,\nu_v)\ne0$.
\formallabel{PositivityTheory.ThK.GaussNe.gauss_ne_zero}
\end{lemma}

\begin{proof}
This is the classical computation. Let $q=\Nm v$, $\varpi$ a uniformiser and $\psi_c(y)=\psi_v(y\varpi^{-c})$, a function on
$\OO_v/\varpi^c$. For $t\in\OO_v$, $y\mapsto\psi_c(ty)$ is a character of $\OO_v/\varpi^c$ which is trivial iff
$t\varpi^{-c}\in\OO_v$, so $\sum_{y\in\OO_v/\varpi^c}\psi_c(ty)=q^c[t\in\varpi^c\OO_v]$ and similarly
$\sum_{y\in\varpi\OO_v/\varpi^c}\psi_c(ty)=q^{c-1}[t\in\varpi^{c-1}\OO_v]$. Writing $x=uy$,
\begin{align*}
|g|^2&=\sum_{x,y}\nu_v(xy^{-1})^{-1}\psi_c(x-y)=\sum_{u}\nu_v(u)^{-1}\sum_{y\in(\OO_v/\varpi^c)^\times}\psi_c((u-1)y)\\
&=q^c-q^{c-1}\sum_{u\equiv1\bmod\varpi^{c-1}}\nu_v(u)^{-1},
\end{align*}
the sums being over $(\OO_v/\varpi^c)^\times$. By minimality of the conductor, $\nu_v$ is non-trivial on
$(1+\varpi^{c-1}\OO_v)/(1+\varpi^c\OO_v)$ (on $(\OO_v/\varpi)^\times$ if $c=1$), so the last sum vanishes.
\end{proof}

\begin{proposition}\label{pf:thmK:facne}
The leaf \ref{pf:thmK:L:facne} holds.
\formallabel{PositivityTheory.ThK.Close.facNe}
\end{proposition}

\begin{proof}
Let a triple be given. The ideal values $\nu(\mathfrak b)^{e_1}\nu(\mathfrak c)^{e_2}$ and the numbers $\nu_v(s)$ are values of
characters, hence non-zero. For $v\in P\cap T$, $\nu$ is ramified at the place of $E$ determined by $\tau_v$, so $\nu_v$ is
non-trivial on $\OO_v^\times$ (\lref{PositivityTheory.ThK.Close.exists_unit_chrW_ne}: a character ramified at a place is
non-trivial on some unit there); it is continuous because $\alpha,\beta$ are, and $\psi$ is unramified at $v\in P$ since $B$
is good. By \cref{pf:thmK:gauss}, $g(\psi_v,\nu_v)\ne0$, hence $\eps(\psi_v,\nu_v,z)\ne0$ by \eqref{pf:thmK:eq:gauss}. So the
step factor \eqref{pf:thmK:eq:stepfac} is non-zero. For an integer $k\ge2$ the Euler product $\thmKLS{k}{\nu}$ converges
absolutely, with non-zero factors, since $|\nu|=1$ and $k>1$; an absolutely convergent product of non-zero factors is
non-zero (\lref{GlobalTheory.LFun.converges_charFam}: the Euler product of a unitary character converges for $\Re s>1$).
\end{proof}

\subsubsection{Local matching}\label{sec:pf:thmK:local}

\begin{proposition}[Local matching]\label{pf:thmK:localmatching}
Assume that matching maps exist at the places of $T$, at the split places outside $T$ and at the inert places (in the sense
of \cref{pf:thmK:def:maps}(2) at a single place). Then \ref{pf:thmK:L:kfin} holds. Moreover \ref{pf:thmK:L:kreal} and
\textup{(Coi)} hold.
\formallabel{PositivityTheory.ThK.kFin_of}
\end{proposition}

\begin{proof}
\ref{pf:thmK:L:kfin}: for every finite $v$, choose a space $\Pi_v$ and maps $p_v,p'_v$ with the local properties of
\cref{pf:thmK:def:maps}(2) at $v$ (including the sphericity at $v\notin T$ and the zeta condition at $v\in P\cap T$),
according as $v\in T$, $v\notin T$ split ($E\to F_v$ exists) or $v\notin T$ inert; together they form finite maps, which
are matching. The local existence statements are \lref{PositivityTheory.BernsteinRamification.Close.finRam} (places of
$T$), \lref{PositivityTheory.BernsteinRamification.Close.finSplit} (split places, by Bernstein--Zelevinsky type arguments)
and \lref{PositivityTheory.Sph.finInert_of} (inert places; it combines the facts that the vacuum generates the
coinvariants at inert places, \lref{PositivityTheory.Sph.vacInert}, deduced from Howe's unramified computation
\lref{PositivityTheory.Sph.howeInert}, cf.\ \cite{Howe1979}; that the spherical functionals of the two sides agree,
\lref{PositivityTheory.Sph.sphFun_of_modSt}, which uses \cref{pf:thmK:modst}; and that local Weil representations exist,
\lref{PositivityTheory.ThKH.Locrep.locRepEx}).

\ref{pf:thmK:L:kreal} is \lref{PositivityTheory.ThKR.kReal'}, which produces matching real maps from the existence of the
infinitesimal oscillator representation on Schwartz space (\lref{PositivityTheory.ThKR.localLift_holds}); the archimedean
quotient is an irreducible quotient of the module generated by the Fock space, identified through the joint harmonics of the
Fock model. \textup{(Coi)} is \lref{PositivityTheory.BernsteinRamification.Close.kCoinv}, which deduces the coinvariant
condition from a dictionary statement at split places (linear models in which the local actions of $\Uni(X)(F_v)\cong
\GL_n(F_v)$ and $\Uni(W)(F_v)$ become explicit); the dictionary statement is \lref{PositivityTheory.KCo.dictSt_of}, from
\lref{PositivityTheory.ThKG.StdDict.dictX} (existence of the dictionary for the standard models) and
\lref{PositivityTheory.KCo.section_statement} (a section statement for linear models).
\end{proof}

\begin{proposition}[Partners]\label{pf:thmK:matchex}
\textup{(Coi)} implies the leaf \ref{pf:thmK:L:matchex}.
\formallabel{PositivityTheory.ThKClose.matchEx_of}
\formallabel{PositivityTheory.ThKClose.Partner.partner_of_big}
\formallabel{PositivityTheory.ThKClose.Partner.partner_of_small}
\formallabel{PositivityTheory.ThKClose.Partner.match_of_polMatch}
\end{proposition}

\begin{proof}
Fix matched data. (a) Let $\Phi=\sum_{i\le m}\thmKpure(f_i,\varphi_i)\in\thmKPD(M)$. Let $T'$ be the union of $T$, $\{v_1\}$
and the finitely many places where some $\varphi_{i,v}\ne\thmKvac_v$. Since $p'(\thmKPol(M'))=\Pi_\infty$ there are
$f'_i\in\thmKPol(M')$ with $p'(f'_i)=p(f_i)$; since $p'_v(\thmKS)=\Pi_v$ there are $\chi_{i,v}\in\thmKS$ with
$p'_v(\chi_{i,v})=p_v(\varphi_{i,v})$; and since $v_1$ is good for $X$, the coinvariant condition (ii) gives
$\zeta_i\in\thmKSo_{v_1}(\tau_1)$ with $p'_{v_1}(\zeta_i)=p_{v_1}(\varphi_{i,v_1})$. Put $\varphi'_{i,v}=\chi_{i,v}$ for
$v\in T'\smallsetminus\{v_1\}$, $\varphi'_{i,v_1}=\zeta_i$, and $\varphi'_{i,v}=\thmKvac'_v$ for $v\notin T'$. Then
$\Phi'=\sum_i\thmKpure(f'_i,\varphi'_i)\in\thmKSo(M')$, and \eqref{pf:thmK:eq:polmatch} holds termwise with $T'$.
(b) The same construction without the condition at $v_1$, exchanging the roles of $p,p'$ (using $p(\thmKPol(M))=\Pi_\infty$
and $p_v(\thmKS)=\Pi_v$). (c) If $\Phi\approx\Phi'$ then $\Phi\in\thmKPD(M)$, $\Phi'\in\thmKPD(M')$; writing each pure tensor
as $f\otimes g$ with $g$ the product of its finite components, the constant sequences show $\Phi\sim\Phi'$ when
$\Phi'\in\thmKSo(M')$.
\end{proof}

\begin{lemma}[The theta integral kills the augmentation module]\label{pf:thmK:killaug}
Let $(M,\omega)$ be a Weil pair for $(X,W)$ such that the theta integrals of all test functions converge, and assume
$\thmKI(\omega(g,1)\Phi)=\thmKI(\Phi)$ for all test functions $\Phi$ and $g\in\Uni(X)(\A)$. Then $\thmKI$ extends to a linear
map $J$ on all functions, and for every finite $v$, Schwartz $f$, family $\varphi$ with $\varphi_u\in\thmKS$ and
$\varphi_u=\thmKvac_u$ for almost all $u$, and $z\in\thmKaug_v(\omega)$, one has $J(\thmKpure(f,\varphi[v\mapsto z]))=0$.
\formallabel{PositivityTheory.ThKF.exists_Ilin}\formallabel{PositivityTheory.ThKF.kill_aug}
\end{lemma}

\begin{proof}
$\thmKI$ is linear on the space of test functions (the theta series are summable and the integrals converge), so it extends
linearly. Write $\thmKpure(f,\varphi[v\mapsto y])=y\otimes_v\rho$ with $\rho$ invariant under translations supported at $v$;
then $y\mapsto J(y\otimes_v\rho)$ is linear, and it suffices to treat a generator $z=y_1-y$ of $\thmKaug_v(\omega)$, with
$\omega(g,1)\colon y\thmKacts{v}y_1$, $g\in\Uni(X)(F_v)$, $y,y_1\in\thmKS$. Then $y\otimes_v\rho$ and $y_1\otimes_v\rho$ are
test functions and $\omega(g,1)(y\otimes_v\rho)=y_1\otimes_v\rho$, so $J(y_1\otimes_v\rho)=\thmKI(\omega(g,1)(y\otimes_v
\rho))=\thmKI(y\otimes_v\rho)$.
\end{proof}

The hypothesis $\thmKI(\omega(g,1)\Phi)=\thmKI(\Phi)$ holds for every Weil pair and test function $\Phi$
(\lref{PositivityTheory.ThKF.iLeft}: by left invariance of the theta function and of the Haar measure); more generally
$\thmKI(\omega(g,h)\Phi)(x)=\thmKI(\Phi)(xh)$ (\lref{PositivityTheory.ThKF.I_move}).

\begin{proposition}[Definite small side]\label{pf:thmK:anall}
The leaf \ref{pf:thmK:L:anall} follows from \ref{pf:thmK:L:weilconv}, \textup{(Coi)} and the left invariance of theta
integrals.
\formallabel{PositivityTheory.ThKF.anAll_of}
\end{proposition}

\begin{proof}
Let matched data with $q=0$, a polynomial constant $c$ and $\Phi\approx\Phi'$ be given, with decompositions
$\Phi=\sum_i\thmKpure(f_i,\varphi_i)$, $\Phi'=\sum_j\thmKpure(f'_j,\varphi'_j)$ and a set $T'$ as in
\eqref{pf:thmK:eq:polmatch}. Since $q=0$, the Weil pair of $V'$ has $\min(q,a'-q)=0$, so by \ref{pf:thmK:L:weilconv} the theta
integrals of all test functions converge and \cref{pf:thmK:killaug} applies to the Weil pair of $V'$. As $v_1$ is good for
$X$, the coinvariant condition (ii) gives $y_j\in\thmKSo_{v_1}(\tau_1)$ with $p'_{v_1}(y_j)=p'_{v_1}(\varphi'_{j,v_1})$. Let
$\varphi''_j=\varphi'_j[v_1\mapsto y_j]$ and $\Phi''=\sum_j\thmKpure(f'_j,\varphi''_j)\in\thmKSo(M')$. Enlarging $T'$ to
$T'\cup\{v_1\}$ (the added tensor factors are the images of the vacua, which agree on both sides by the sphericity
$p_v(\thmKvac_v)=p'_v(\thmKvac'_v)$ for $v\notin T$), we get $\Phi\approx\Phi''$, since the $v_1$-components have the same
images. Hence $\thmKI(\Phi)=c\,\thmKI'(\Phi'')$. Finally $y_j-\varphi'_{j,v_1}\in\ker p'_{v_1}\cap\thmKS=\thmKaug_{v_1}(\omega')$
by the coinvariant condition (i), and $\thmKpure(f'_j,\varphi''_j)=\thmKpure(f'_j,\varphi'_j)+\thmKpure(f'_j,\varphi'_j[v_1
\mapsto y_j-\varphi'_{j,v_1}])$, so by \cref{pf:thmK:killaug} and linearity $\thmKI'(\Phi'')=\thmKI'(\Phi')$.
\end{proof}

\subsubsection{Change of Cartan data}\label{sec:pf:thmK:polre}

\begin{proposition}\label{pf:thmK:polre}
The leaf \ref{pf:thmK:L:polre} holds.
\formallabel{PositivityTheory.HoweDualityUnitaryAlt.polRe_of2T}\formallabel{GlobalTheory.CompanionTheory.ThetaK.polRe_of2T}
\end{proposition}

\begin{proof}
Let matched data on $G$ (with $R,R_f$) and on a re-choice $G_2=G[V_2,V'_2,W_2]$ (with $R_2,R_{f,2}$, the same $\omega,\omega'$)
be given, and let $c_2\ne0$ be a polynomial constant of $G_2$. There is $h\in\Uni(W)(F_\infty)$ conjugating the Cartan
involutions of $W_2$ into those of $W$ (\lref{PositivityTheory.ThKF.cartan_conjugacy_holds}: Cartan involutions of a
skew-hermitian space of fixed signature are conjugate), and then there are $g\in\Uni(V)(F_\infty)$, $g'\in\Uni(V')(F_\infty)$
such that the archimedean operators $T$, $T'$ induced by $\omega(g,h)$, $\omega'(g',h)$ map the Fock spaces of $G_2$ onto
those of $G$ (\lref{PositivityTheory.ThKF.cartanMove_of}: archimedean operators of the Weil representation transport the
Fock space of the conjugated Cartan data; the operators exist by \lref{PositivityTheory.ThKF.A.archOp}). Comparing the
matching maps: there are a linear map $e$ of the archimedean quotients and $c\ne0$ with $e\circ p\circ T=p_2$ and $e\circ
p'\circ T'=c\,p'_2$ on the Fock spaces of $G_2$ (\lref{PositivityTheory.ThKF.archimedean_matched_maps_compare}: two systems
of matching real maps for the same triple differ by a linear map and a scalar on the small side; it rests on the joint
harmonic description of the Fock spaces, \lref{PositivityTheory.ThKR.Join.joinTri_of}, and archimedean uniqueness); and
for every finite $v$ there are a linear map $e_v$ and $c_v\ne0$ with $e_v\circ p_v=p_{2,v}$ and $e_v\circ p'_v=c_vp'_{2,v}$ on
$\thmKS$ (\lref{GlobalTheory.CompanionTheory.ThetaK.cmp_finT}: the same comparison at a finite place, using the maximality
of the local theta quotients, \lref{PositivityTheory.Sph.maxFinT_of}); for $v\notin T$, $c_v=1$ by comparing the images
of the vacua. Put $\Lambda=c\prod_{v\in T}c_v$.

If $\Phi\approx\Phi'$ on $G$ with $\Phi'\in\thmKSo(M')$, write the Fock components of $\Phi,\Phi'$ as $Tu_i$, $T'u'_j$ with
$u_i,u'_j$ in the Fock spaces of $G_2$, and let $\Phi_2=\sum_i\thmKpure(u_i,\varphi_i)$, $\Phi'_2=\sum_j\thmKpure(u'_j,
\varphi'_j)$. Then $\Phi=\omega(g,h)\Phi_2$ and $\Phi'=\omega'(g',h)\Phi'_2$; $\Phi'_2\in\thmKSo$ because $\thmKSo$ is stable
under $\omega'$ (coinvariant condition); and applying $e\otimes\bigotimes_ve_v$ to \eqref{pf:thmK:eq:polmatch} shows
$\Phi_2\approx\Lambda\Phi'_2$ on $G_2$. Since $\thmKI(\omega(g,h)\Psi)(x)=\thmKI(\Psi)(xh)$ for test functions $\Psi$, we get
$\thmKI(\Phi)(x)=\thmKI(\Phi_2)(xh)=c_2\Lambda\,\thmKI'(\Phi'_2)(xh)=c_2\Lambda\,\thmKI'(\Phi')(x)$. So $c_2\Lambda\ne0$ is a
polynomial constant of $G$.
\end{proof}


\subsubsection{Non-vanishing of theta integrals}\label{sec:pf:thmK:nonzero}
We use three auxiliary statements about Weil pairs and triples. Polynomial data are test functions
(\lref{GlobalTheory.ThK.Dbl.polSch}). If $\min(r,s)=0$ and $\thmKI$ vanishes on $\thmKPD(M)$, then it vanishes on all test
functions (\lref{GlobalTheory.ThK.Dbl.schToPol_of}; by the continuity in \ref{pf:thmK:L:weilconv} and the density
\ref{pf:thmK:L:poldense}, as in \cref{pf:thmK:nzmove}). For a triple with $q=0$, every $\Phi'\in\thmKPD(M')$ has some
$\Psi\in\thmKPD(M')\cap\thmKSo(M')$ with $\thmKI'(\Psi)=\thmKI'(\Phi')$ (\lref{GlobalTheory.ThK.Dbl.toSo_of}, which replaces
the component at $v_1$ by an element of $\thmKSo_{v_1}(\tau_1)$ with the same image in the coinvariants, as in the proof of
\cref{pf:thmK:anall}).

\begin{lemma}[Identity theorem]\label{pf:thmK:identity}
Let $U\subseteq\C$ be open with discrete complement, let $f$ be holomorphic on $U\cap\{\Re z>0\}$ and suppose that there is
$t$ with $z\in U$ and $f(z)=0$ whenever $\Re z>t$. Then $f(s)=0$ for all $s\in U$ with $\Re s>0$.
\formallabel{GlobalTheory.ThK.Dbl.eq_zero_of_right}
\end{lemma}

\begin{proof}
The complement of $U$ is closed and discrete, hence countable. The map $w\mapsto e^{\Re w}+i\Im w$ is a homeomorphism of
$\C$ onto the half plane $\{\Re z>0\}$; the preimage of the complement of $U$ is countable, and the complement of a countable
subset of $\C$ is connected. Hence $U\cap\{\Re z>0\}$ is connected, and the identity theorem applies.
\end{proof}

\begin{proposition}[The doubling argument]\label{pf:thmK:nonzerobot}
For a triple with $p=q=0$ there are $\Phi\in\thmKPD(M)$ with $\thmKI(\Phi)\ne0$ and $\Phi'\in\thmKPD(M')$ with
$\thmKI'(\Phi')\ne0$.
\formallabel{GlobalTheory.ThK.Dbl.nonZeroBottom_of_doubling_integral}\formallabel{GlobalTheory.ThK.Dbl.pair_ne_zero}
\end{proposition}

The proof uses the following statements about a Weil pair $(M,\omega)$ for $(X,W)$ of signature $(r,s;r',s')$ with
characters $(\chi_X,\chi_W)$, $\min(r,s)=\min(r',s')=0$, $1\le n<d$ and $d$ even. Let $G_2(X)$ be the doubled group with its
stabiliser $K$ and height function $a$ (\cref{sec:pf:thmK:sw}), $\iota\colon\Uni(X)\times\Uni(X)\to G_2(X)$ the doubling
embedding, $\nu$ a Haar measure on $\A^{n\times d}$, and for $\Phi\in\thmKPD(M)$ and $z\in\C$ let
\begin{equation}\label{pf:thmK:eq:zint}
Z_\Phi(z)=\int_{\Uni(X)(\A)}a(\iota(g,1))^{z-(d-n)/2}\,\langle\omega(g,1)\Phi,\Phi\rangle\,d\mu(g),\qquad
\langle\Phi_1,\Phi_2\rangle=\int\Phi_1\overline{\Phi_2}\,d\nu .
\end{equation}
\begin{enumerate}[label=(D\arabic*)]
\item\label{pf:thmK:D:norm} (\lref{GlobalTheory.ThK.Dbl.normZeta_of}, from the Siegel--Weil formula for the doubled pair
and the meromorphic continuation.) There are an open set $U\ni(d-n)/2$ whose complement is discrete, a holomorphic
function $Z$ on $U$ and $t_0$ such that for $\Re z>t_0$, $z\in U$, the integral \eqref{pf:thmK:eq:zint} converges and equals
$Z(z)$; and if $\thmKI(\Phi)=0$ then $Z((d-n)/2)=0$.
\item\label{pf:thmK:D:euler} (\lref{GlobalTheory.ThK.Dbl.eulerProd}, \lref{GlobalTheory.ThK.Dbl.locMeas},
\lref{GlobalTheory.ThK.Dbl.locUnr_of}.) For a finite set $T$ of places such that $\chi_W$ is unramified outside $T$, an
integrable function on $\Uni(X)(\A)$ with spherical behaviour at all $v\notin T$ integrates to the limit of the products of its
local integrals at $v\notin T$ times its integral over the elements integral outside $T$; and for a spherical function $f$ at
$v\notin T$, $\int_{\Uni(X)(F_v)}f(\iota(u,1))\,du\cdot d_v(z)=\prod_{i=0}^{n-1}L_v(z+\tfrac n2-i,\chi_W)$, where
$d_v(z)=\prod_{i=0}^{n-1}(1-\epsilon_v^{\,n-1-i}\Nm v^{-(2z+i+1)})^{-1}$, with $\epsilon_v=1$ if $v$ splits in $E$ and $-1$
otherwise (the unramified doubling integral; at inert places it is computed in \lref{GlobalTheory.ThK.Dbl.Inert.main}).
\item\label{pf:thmK:D:dprod} (\lref{GlobalTheory.ThK.Dbl.dProd}.) $\prod_{v\notin T}d_v(z)$ converges to a non-vanishing
holomorphic function $d_T(z)$ on $\Re z>0$.
\item\label{pf:thmK:D:bad} (\lref{GlobalTheory.ThK.Dbl.realInv_of}, \lref{GlobalTheory.ThK.Dbl.badFactor}.) If $M$ has a
non-trivial real quotient, there is a non-zero $K_\infty(X)$-invariant $f_0\in\thmKPol(M)$; and then, for $T\subseteq P$ with
$\chi_X,\chi_W$ unramified outside $T$, there is $\Phi\in\thmKPD(M)$ whose integrand in \eqref{pf:thmK:eq:zint} is spherical at
all $v\notin T$, whose integral $H(z)$ over the elements integral outside $T$ is entire, and with $H((d-n)/2)\ne0$.
\item\label{pf:thmK:D:lint} (\lref{GlobalTheory.ThK.Dbl.lInt2}.) For a triple and $\chi\in\{\nu,\nu^{-1}\}$ there is an
entire function $\Lambda$ with $\Lambda(s)=L^{\emptyset}(s,\chi)$ for $\Re s>1$ and $\Lambda(z)\ne0$ for all integers $z$ with
$2z\ge-(a+a')$.
\end{enumerate}

\begin{proof}[Proof of \cref{pf:thmK:nonzerobot}]
We first show, for a Weil pair as above, a finite set $T\subseteq P$ with $\chi_X,\chi_W$ unramified outside $T$ and $\chi_W$
ramified at all places above $T$, a non-trivial real quotient of $M$, and an entire $\Lambda$ with $\Lambda=L^\emptyset(\cdot,
\chi_W)$ on $\Re s>1$ and $\Lambda(z)\ne0$ for integers $2z\ge-d$: there is $\Phi\in\thmKPD(M)$ with $\thmKI(\Phi)\ne0$.
Take $\Phi$ and $H$ from \ref{pf:thmK:D:bad}. For $\Re z>\max(t_0,n/2)$, \ref{pf:thmK:D:euler} and \ref{pf:thmK:D:dprod} give
\begin{equation}\label{pf:thmK:eq:dbl}
Z(z)\,d_T(z)=\prod_{i=0}^{n-1}\Lambda\bigl(z+\tfrac n2-i\bigr)\cdot H(z):
\end{equation}
the local integrals at $v\notin T$ multiplied by $d_v(z)$ give $\prod_iL_v(z+n/2-i,\chi_W)$; the product over $v\notin T$ of
these local factors is $\prod_iL^\emptyset(z+n/2-i,\chi_W)$, because $\chi_W$ is ramified above $T$ (so its local factors there
are $1$); and $\Re(z+n/2-i)>1$ for $i<n$. Both sides of \eqref{pf:thmK:eq:dbl} are holomorphic on $U\cap\{\Re z>0\}$, so by
\cref{pf:thmK:identity} the identity holds at $z_0=(d-n)/2$, which lies in $U$ and has $\Re z_0>0$. Suppose $\thmKI(\Phi)=0$;
then $Z(z_0)=0$ by \ref{pf:thmK:D:norm}, so $\prod_i\Lambda(z_0+n/2-i)\,H(z_0)=0$. But with $d=2k$, $z_0+n/2-i=k-i$ is an
integer with $2(k-i)\ge-d$ (as $i<n<d$), so $\Lambda(k-i)\ne0$; and $H(z_0)\ne0$: a contradiction.

Now let a triple with $p=q=0$ be given. By \ref{pf:thmK:L:kreal} there are matching real maps, whose quotients for $M$ and $M'$
are non-trivial. $\nu=\alpha/\beta$ is a Hecke character with $\nu|_{\A^\times}=\eta^{a+a'}=1$ ($a+a'$ even), unramified
outside $T$ and ramified above $T$, and $T\subseteq P$. Apply the above to the Weil pair of $V$ ($n=a<d=a+a'$,
$\chi_W=\beta/\alpha=\nu^{-1}$, $\Lambda$ from \ref{pf:thmK:D:lint} for $\nu^{-1}$) and to the Weil pair of $V'$ ($n=a'<d$,
$\chi_W=\nu$, $\Lambda$ for $\nu$); $V,V',W$ are definite since $p=q=0$.
\end{proof}

\begin{proposition}\label{pf:thmK:nonzeroan}
The leaf \ref{pf:thmK:L:nonzeroan} holds.
\formallabel{GlobalTheory.ThK.Dbl.nonZeroAn_of}\formallabel{GlobalTheory.ThK.Dbl.big_of}\formallabel{GlobalTheory.ThK.Dbl.small_of}
\end{proposition}

\begin{proof}
\emph{The case $p=0$.} We show by induction on $q$: for every triple with $p=0$ there is $\Phi\in\thmKPD(M)$ with
$\thmKI(\Phi)\ne0$. For $q=0$ this is \cref{pf:thmK:nonzerobot}. Let $q\ge1$ and write $a'=b+2$. \ref{pf:thmK:L:towerB} gives a
re-choice $G_2$ carrying the triple and $G_1$ of type $(a,b)$ forming a $\beta$-step configuration $(G_2,G_1,H,H')$, with a
triple $(\omega_1,\omega'_1)$ on $G_1$; $p(G_1)=0$, $q(G_1)=q-1$. By induction there is $\Phi_1\in\thmKPD(M_1)$ with
$\thmKI_1(\Phi_1)\ne0$. \ref{pf:thmK:L:evzero} gives $e\in\thmKEv(M,M_1;1,\iota_{H'};I_0)$ (for the Weil pair of $V$ in $G_2$)
mapping $\thmKPD(M_2)$ onto $\thmKPD(M_1)$, so $\Phi_1=e\Phi$ with $\Phi\in\thmKPD(M_2)$; let $C\ne0$ be the constant of
\ref{pf:thmK:L:ctbz}. Suppose $\thmKI$ vanished on $\thmKPD(M)$; since $V$ is definite, it then vanishes on all test functions,
in particular $\thmKI(\Phi)=0$. Then $(0,0)\in\thmKSpl(\thmKI(\Phi))$, and \ref{pf:thmK:L:ctbz} (with $h=1$ and a Levi element
$m=m(1,g)$, which exists by \lref{GlobalTheory.ThK.Dbl.leviEx}) gives $0=C\,\thmKI_1(\Phi_1)(g)$ for every $g$, a
contradiction.

\emph{The case $q=0$.} By induction on $p$: for every triple with $q=0$ there is $\Phi'\in\thmKPD(M')$ with
$\thmKI'(\Phi')\ne0$. For $p=0$ this is \cref{pf:thmK:nonzerobot}. For $p\ge1$, write $a=b+2$; \ref{pf:thmK:L:towerA} gives
$G_2$ and $G_1$ of type $(b,a')$ forming an $\alpha$-step configuration, with $p(G_1)=p-1$, $q(G_1)=0$, and by induction
$\Phi'_1\in\thmKPD(M'_1)$ with $\thmKI'_1(\Phi'_1)\ne0$; write $\Phi'_1=e'\Phi'$ by \ref{pf:thmK:L:evzero}. If $\thmKI'$ vanished
on $\thmKPD(M')$, then ($V'$ definite) $\thmKI'(\Phi')=0$; by \ref{pf:thmK:L:ctsz} there is $f'_\alpha$ with
$(f'_\alpha,0)\in\thmKSpl(0)$ and $f'_\alpha(m)=C\thmKI'_1(\Phi'_1)(g)$ for $m=m(1,g)$; by the uniqueness in
\ref{pf:thmK:L:splitex} (applied to $0=\thmKI(0)$, $0\in\thmKPD(M_2)$, for which $(0,0)$ is a splitting) $f'_\alpha=0$, a
contradiction. Finally, the replacement statement recalled above turns $\Phi'$ into an element of $\thmKPD(M')\cap\thmKSo(M')$
with the same theta integral.
\end{proof}

\subsubsection{The definite case}\label{sec:pf:thmK:polzero}
Let matched data with $p=q=0$ be given. Then $V$, $V'$ and $W$ are definite, and their automorphic quotients are compact:
for $Y$ hermitian or skew-hermitian with $\min(r,s)=0$, $\Uni(Y)(\A)$ is locally compact and second countable, $\Gamma_Y$ is
countable, there is a compact set meeting every orbit of $\Gamma_Y$, fundamental domains have finite non-zero volume and Haar
measures are right invariant (\lref{GlobalTheory.ThK.Dbl.cptAll_of}). In particular continuous left
$\Gamma_Y$-invariant functions are bounded, rapidly decreasing, and vanish if their Petersson norm vanishes. Fix a Haar
measure $\mu$ on $\Uni(W)(\A)$ and a fundamental domain $D$. Let $\mathcal Y=\{(\Phi,\Phi'):\Phi'\in\thmKSo(M'),\
\Phi\approx\Phi'\}$; for $y=(\Phi,\Phi')$ and $s,t\in\C$ put $\phi_{s,t}(y)=s\,\thmKI(\Phi)+t\,\thmKI'(\Phi')$, a continuous
left $\Gamma_W$-invariant function, and $\langle\phi_1,\phi_2\rangle=\int_D\phi_1\overline{\phi_2}\,d\mu$. For a function
$\phi$ on $\Uni(W)(\A)$ and a test function $\Psi$, let $\vartheta'(\phi,\Psi)(g)=\langle\phi,\theta_{\omega'}(\Psi)(g,
\cdot)\rangle$, a function on $\Uni(V')(\A)$. We use three statements, valid for matched data with $p=q=0$:
\begin{enumerate}[label=(P\arabic*)]
\item\label{pf:thmK:P:herm} (\lref{GlobalTheory.ThK.Dbl.hermLine_of'}, an irreducibility argument for the
$\Uni(W)$-stable space $\mathcal Y$.) For all $s,t,s',t'$ there is $\lambda$ with $\langle\phi_{s,t}(y),\phi_{s',t'}(y')
\rangle=\lambda\,\langle\thmKI'(\Phi'),\thmKI'(\Phi'_2)\rangle$ for all $y=(\Phi,\Phi')$, $y'=(\Phi_2,\Phi'_2)$ in $\mathcal Y$.
\item\label{pf:thmK:P:coef} (\lref{GlobalTheory.ThK.Dbl.coefLine_of}, using the Siegel--Weil formula for the doubled
pair, derived from \cref{pf:thmK:sw:thm,pf:thmK:eis}.) If $\langle\phi_{s,t}(y),\phi_{s,t}(y')\rangle=\lambda\langle
\thmKI'(\Phi'),\thmKI'(\Phi'_2)\rangle$ for all $y,y'\in\mathcal Y$, then for all $y\in\mathcal Y$ and $\Psi\in\thmKPD(M')$,
$\|\vartheta'(\phi_{s,t}(y),\Psi)\|^2=\lambda\,\|\vartheta'(\thmKI'(\Phi'),\Psi)\|^2$, the norms being Petersson norms on
$\Uni(V')$ over the domain of $M'$.
\item\label{pf:thmK:P:const} (\lref{GlobalTheory.ThK.Dbl.thConst_of'}, using that Hecke $L$-functions are entire and
non-zero at $s=1$, \cref{pf:thmK:hecke}.) For $y\in\mathcal Y$ and test functions $\Psi$, the functions
$g\mapsto\vartheta'(\phi_{s,t}(y),\Psi)(g)$ (and the analogous ones on $\Uni(V)$) are constant.
\end{enumerate}

\begin{littheorem}[Hecke; \cite{Hecke1920,Tate1967thesis}]\label{pf:thmK:hecke}
For a number field $K$ and a unitary Hecke character $\chi$ of $K$: if $\chi$ is not of the form $|\cdot|^{it}$, the
Euler product $L(s,\chi)$ over the finite places ($\Re s>1$) extends to an entire function; and if $\chi\ne1$, it extends
analytically to a neighbourhood of $s=1$ with a non-zero value at $s=1$.
\formal{GlobalTheory.Published.HeckeLEntireAndNonzeroAtOne}
\formallabel{GlobalTheory.Published.HeckeLEntireAndNonzeroAtOne_holds}
\end{littheorem}

This published statement enters the proof of Theorem~K only through \ref{pf:thmK:P:const} and the orthogonality
statement \ref{pf:thmK:E:orthiso} of the end steps.

\begin{lemma}[Fubini on compact quotients]\label{pf:thmK:fubini}
For a continuous left $\Gamma_W$-invariant function $\phi$ on $\Uni(W)(\A)$ and a test function $\Psi$,
$\int_{\mathrm{dom}'}\vartheta'(\phi,\Psi)\,d\mu'=\langle\phi,\thmKI'(\Psi)\rangle$, where $(\mu',\mathrm{dom}')$ are the
measure and domain of $M'$.
\formallabel{GlobalTheory.ThK.Dbl.int_th}
\end{lemma}

\begin{proof}
The function $(g,h)\mapsto\phi(h)\overline{\theta_{\omega'}(\Psi)(g,h)}$ is continuous and invariant under $\Gamma_{V'}\times
\Gamma_W$ (\cref{pf:thmK:analytic}(a)), hence bounded, and both domains have finite measure; so Fubini's theorem applies, and
the inner integral over $\mathrm{dom}'$ is $\overline{\thmKI'(\Psi)(h)}$.
\end{proof}

\begin{proposition}\label{pf:thmK:polzeroprop}
The leaf \ref{pf:thmK:L:polzero} holds.
\formallabel{PositivityTheory.ThKR.polynomial_vanishing_of_compactness}\formallabel{GlobalTheory.ThK.Dbl.polZero_of}
\end{proposition}

\begin{proof}
Let matched data with $p=q=0$ be given. Elements of $\mathcal Y$ are test functions, so the functions $\phi_{s,t}(y)$ are
continuous and $\Gamma_W$-invariant, and all Petersson products below converge.

By \ref{pf:thmK:P:herm} with $(s,t,s',t')=(1,0,0,1)$ there is $c$ with $\langle\thmKI(\Phi),\thmKI'(\Phi'_2)\rangle=c\,
\langle\thmKI'(\Phi'),\thmKI'(\Phi'_2)\rangle$ for $y,y'\in\mathcal Y$; in particular
\begin{equation}\label{pf:thmK:eq:orth}
\langle\phi_{1,-c}(y),\thmKI'(\Phi')\rangle=0\qquad(y=(\Phi,\Phi')\in\mathcal Y).
\end{equation}
By \ref{pf:thmK:P:herm} with $(1,-c,1,-c)$ there is $\lambda$ with $\langle\phi_{1,-c}(y),\phi_{1,-c}(y')\rangle=\lambda\langle
\thmKI'(\Phi'),\thmKI'(\Phi'_2)\rangle$ for $y,y'\in\mathcal Y$.

By \ref{pf:thmK:L:nonzeroan} there are $\Phi_0\in\thmKPD(M)$ with $\thmKI(\Phi_0)\ne0$ and $\Phi'\in\thmKPD(M')\cap\thmKSo(M')$
with $\thmKI'(\Phi')\ne0$; by \ref{pf:thmK:L:matchex}(b) there is $\Phi$ with $y=(\Phi,\Phi')\in\mathcal Y$. By
\ref{pf:thmK:P:const} and \cref{pf:thmK:fubini}, $\vartheta'(\phi_{s,t}(y),\Phi')$ is the constant
$\langle\phi_{s,t}(y),\thmKI'(\Phi')\rangle/\mu'(\mathrm{dom}')$. Apply \ref{pf:thmK:P:coef} with $(s,t)=(1,-c)$, $\lambda$, $y$
and $\Psi=\Phi'$: by \eqref{pf:thmK:eq:orth} the left side is $0$, and the right side is $\lambda\,|x|^2/\mu'(\mathrm{dom}')$
with $x=\langle\thmKI'(\Phi'),\thmKI'(\Phi')\rangle$. Since $\thmKI'(\Phi')$ is continuous, invariant and non-zero, $x\ne0$;
hence $\lambda=0$.

Now for every $z=(\Psi,\Psi')\in\mathcal Y$, $\|\phi_{1,-c}(z)\|^2=\lambda\|\thmKI'(\Psi')\|^2=0$, so $\phi_{1,-c}(z)=0$, i.e.\
$\thmKI(\Psi)=c\,\thmKI'(\Psi')$: $c$ is a polynomial constant. If $c=0$, take $\Psi'\in\thmKSo(M')$ with $\Phi_0\approx\Psi'$
(\ref{pf:thmK:L:matchex}(a)); then $\thmKI(\Phi_0)=c\,\thmKI'(\Psi')=0$, a contradiction. So $c\ne0$.
\end{proof}

\subsubsection{Local factors of the steps}\label{sec:pf:thmK:loc}

\begin{proposition}\label{pf:thmK:locAB}
The leaves \ref{pf:thmK:L:locA} and \ref{pf:thmK:L:locB} hold.
\formallabel{PositivityTheory.ThKG.Loc.locA_of}\formallabel{PositivityTheory.ThKG.Loc.locB_of}
\end{proposition}

\begin{proof}
We prove \ref{pf:thmK:L:locA}; the proof of \ref{pf:thmK:L:locB} is the same with the roles of $V,V'$ and of the two
characters exchanged, using the second part of \lref{PositivityTheory.ThKG.Loc.Cone.coneFac_all}, the local computation
\lref{PositivityTheory.ThKG.Std.stdB} (the analogue of (b) below for $\beta$-steps) and the bump functions
\lref{PositivityTheory.ThKG.Loc.Bump.bumpB_fin}; in addition, the local cone integral at $v_1$ maps
$\thmKSo_{v_1}(\tau_1)$ into $\thmKSo_{v_1}(\tau_1)$ and is compatible with the augmentation modules at $v_1$
(\lref{PositivityTheory.ThKG.Loc.FL.fLocB_of}), which gives $\thmKcone_B\Phi'\in\thmKSo(M'_1)$. The proof combines four
groups of statements.
\begin{enumerate}[label=(\alph*)]
\item \emph{Factorisation of the cone integral} (\cref{pf:thmK:conefac}): there are signs $e_1,e_2$, $\kappa\ne0$, a
continuous linear map $L$ of Schwartz spaces mapping $\thmKPol(M)$ into $\thmKPol(M_1)$, the local cone integrals $\Lambda_v$
(preserving $\thmKS$) and numbers $u_v\ne0$ ($v\notin T$) with $\Lambda_v(\thmKvac_v)=u_v\thmKvac^{M_1}_v$, such that
$\thmKcone_A$ is linear and \eqref{pf:thmK:eq:conefac} holds for every finite $T'\supseteq T$ and every pure tensor
$\thmKpure(f,\varphi)$ with $\varphi_v=\thmKvac_v$ outside $T'$.
\item \emph{Local computation at the places of $T$} (\lref{PositivityTheory.ThKG.Std.stdA}): there is $s\in\{\pm1\}$
such that for $v\in T$ there are $\varphi,\varphi'\in\thmKS$ with $p_v(\varphi)=p'_v(\varphi')$, $p'_{1,v}(e'_v\varphi')\ne0$
and $p_{1,v}(\Lambda_v\varphi)=\nu_v(s)\,\eps(\psi_v,\nu_v,1-k_A)\,p'_{1,v}(e'_v\varphi')$.
\item \emph{Induced local maps} (\lref{PositivityTheory.ThKG.Loc.lams}): the global maps $\thmKcone_A$ and $e'$ induce
linear maps $\lambda_\infty,\lambda'_\infty\colon\Pi_\infty(G)\to\Pi_\infty(G_1)$ and $\lambda_v,\lambda'_v\colon\Pi_v(G)\to
\Pi_v(G_1)$ with $\lambda'_\infty\circ p'=p'_1\circ e'_\infty$, $\lambda_\infty\circ p=p_1\circ L$,
$\lambda'_v\circ p'_v=p'_{1,v}\circ e'_v$, $\lambda_v\circ p_v=p_{1,v}\circ\Lambda_v$, all equivariant for the parabolic
subgroup of $H'$ with the character $(\alpha,a'/2)$ of the Levi torus. Well-definedness uses that a theta integral of a pure
tensor vanishes as soon as one local component is killed by the matching map
(\lref{PositivityTheory.Sph.theta_kills_cusp_forms_spherical}, a consequence of square integrability, archimedean
uniqueness and the spherical statements) together with an injectivity statement for the smaller configuration; equivariance
comes from the constant-term formulas \ref{pf:thmK:L:ctsz}, \ref{pf:thmK:L:ctbc} and the invariance of the components of the
constant terms (\lref{PositivityTheory.ThKCT.splitInv2_of}). The maps are non-zero by $\mathcal B(p+q)$ (the theta
integrals of $G_1$ do not vanish identically) and by bump functions at the archimedean place
(\lref{PositivityTheory.ThKG.Loc.Bump.bumpA_fin}).
\item \emph{Multiplicity one}: two non-zero equivariant maps between the archimedean quotients are proportional
(\lref{PositivityTheory.ThKG.CayleyTransformDarboux.archimedean_multiplicity_one_of_lie}); so are two such maps at a place
$v\in T$, the first being non-zero (hypothesis \ref{pf:thmK:H:ram} of \cref{pf:thmK:closesph}); and at $v\notin T$ two such
maps are proportional with the ratio of their values on the image of the vacuum
(\lref{PositivityTheory.ThKG.Loc.unramifiedProportional_holds}).
\end{enumerate}
By (d) there is $c\ne0$ with $\lambda'_\infty=c\lambda_\infty$; for $v\in T$ (using (b)) $\lambda_v=\nu_v(s)\eps(\psi_v,\nu_v,
1-k_A)\,\lambda'_v$; and for $v\notin T$ (using $\Lambda_v(\thmKvac_v)=u_v\thmKvac_v$ and sphericity) $\lambda_v=u_v\lambda'_v$.
The constant $\kappa c$ does not depend on the triples (the maps $\lambda_\infty$ attached to two triples coincide, and the
proportionality constant is unique; \lref{PositivityTheory.ThKG.Loc.one_const} is this elementary uniqueness argument).
Let $\Phi\approx\Phi'$ with decompositions over $T'$. Apply $\lambda'_\infty\otimes\bigotimes_{v\in T'}\lambda'_v$ to
\eqref{pf:thmK:eq:polmatch}; replacing $\lambda'$ by $\lambda$ on the left side produces the factor
$c^{-1}\prod_{v\in T}(\nu_v(s)\eps(\psi_v,\nu_v,1-k_A))^{-1}\prod_{v\in T'\smallsetminus T}u_v^{-1}$. By
\eqref{pf:thmK:eq:conefac} and the definition \eqref{pf:thmK:eq:stepfac} of the step factor (note $T\subseteq P$), the
resulting identity of tensors says that $\thmKcone_A\Phi\approx C\,\thmKst(e_1,e_2,s;\mathfrak b_H,\mathfrak b_{H'};1-k_A)\,
\thmKLS{k_A}{\nu}\,e'\Phi'$ on $G_1$, with $C=\kappa c$ (\lref{PositivityTheory.ThKG.Loc.tens_step} is this tensor
computation). If there is no triple on $G$ the statement holds with $C=1$.
\end{proof}

\subsubsection{End steps}\label{sec:pf:thmK:end}\label{sec:pf:thmK:endsteps}
Fix a triple. Let $\mathcal D$ be the set of functions $\thmKI(\Phi)-\thmKI'(\Phi')$ with $\Phi\in\thmKPD(M)$ and
$\Phi'\in\thmKPD(M')\cap\thmKSo(M')$. A function on $\Uni(W)(\A)$ is \emph{automorphic} if it is continuous and left
$\Gamma_W$-invariant; it is \emph{cuspidal} if for every non-zero totally isotropic subspace $U\subseteq E^{a+a'}$ its integral
over a fundamental domain of $\Gamma_W\cap N_U$ in the unipotent radical $N_U$ of the stabiliser of $U$ (for any Haar measure)
vanishes after any right translation. Two functions $g,f$ are \emph{orthogonal}, $g\perp f$, if $\int_{\mathcal F}g\bar f\,
d\mu=0$ for every Haar measure $\mu$ on $\Uni(W)(\A)$ and every fundamental domain $\mathcal F$ of $\Gamma_W$. A function $f$
is \emph{pure of law $(\chi,z)$} if for some plane $H'$ in $W$, $\thmKCT_{H'}(f)$ exists and has law $(\chi,z)$. Let $K_W$ be
the set of $k\in\Uni(W)(\A)$ whose archimedean components commute with the Cartan involutions of $W$. For a plane $H'$ adapted
to $(W,W_1)$, a function $T$ on $\Uni(W)(\A)$ and $k\in K_W$, a function $f_1$ on $\Uni(W_1)(\A)$ is a \emph{low} of $T$ at
$k$ if $f_1(h_1)=T(m(1,h_1)k)$ for all Levi elements $m(1,h_1)$. Finally, for $G$ of type $(b+2,a')$ and $G_1$ of type
$(b,a')$, $(G,G_1,H,H')$ is a \emph{plane-removal configuration} if $V'(G_1)=V'(G)$, $q(G_1)=q(G)$, $p(G_1)+1=p(G)$, $H$ is
adapted to $(V(G),V(G_1))$ at $(0,b+1)$ and $H'$ to $(W(G),W(G_1))$ at $(0,b+1)$. The end steps rest on the following fifteen
statements, quantified over all triples (and over the indicated further data); \lref{PositivityTheory.ThKEnd.Leaves}
collects them.
\begin{enumerate}[label=(E\arabic*)]
\item\label{pf:thmK:E:iwa} For a plane $H'$ adapted to $(W,W_1)$, every $x\in\Uni(W)(\A)$ is $u\,m(t,1)\,m(1,h_1)\,k$ with
$u\in N(H')$ and $k\in K_W$.
\item\label{pf:thmK:E:ct} Constant terms along $H'$ are linear on automorphic functions; and if $\thmKCT_{H'}(f)=T_1+T_2$ for
an automorphic $f$ with $T_1,T_2$ of laws with distinct exponents, then $T_1,T_2$ are left $N(H')$-invariant.
\item\label{pf:thmK:E:linemove} If $f$ is automorphic and pure of law $(\chi,z)$, then for every plane $H'$ in $W$,
$\thmKCT_{H'}(f)$ exists and has law $(\chi,z)$.
\item\label{pf:thmK:E:aut} Theta integrals of elements of $\thmKPD(M)$ and of $\thmKPD(M')\cap\thmKSo(M')$ are automorphic.
\item\label{pf:thmK:E:stab} For matched data, $k\in K_W$ and $\Phi\approx\Phi'$ with $\Phi'\in\thmKSo(M')$:
$\omega'(1,k)\Phi'\in\thmKSo(M')$ and $\omega(1,k)\Phi\approx\omega'(1,k)\Phi'$.
\item\label{pf:thmK:E:sigre} If a re-choice $G[V_2,V'_2,W_2]$ (with the same $\omega,\omega'$) is again part of a triple, there
is $h_0\in\Uni(W)(\A)$ such that $f\mapsto f(\cdot\,h_0)$ maps $\mathcal D$ into the set $\mathcal D$ of the re-choice.
\item\label{pf:thmK:E:tower2} If $G$ has type $(b+2,a')$ and $p\ge1$, there are a re-choice of $G$ which is part of a triple
and a plane-removal configuration $(G[V_2,V'_2,W_2],G_1,H,H')$ with a triple on $G_1$.
\item\label{pf:thmK:E:alphazero} If $p=0<q$, every $f\in\mathcal D$ with a splitting $(f_\alpha,f_\beta)$ along a plane
adapted to $(W,W_1)$, $W_1$ skew-hermitian, has $f_\alpha=0$.
\item\label{pf:thmK:E:beta} If $q=0<p$, every $f\in\mathcal D$ with a splitting $(f_\alpha,f_\beta)$ along a plane adapted to
$(W,W_1)$ has $f_\beta=0$.
\item\label{pf:thmK:E:alphalow} In a plane-removal configuration with triples on $G$ and $G_1$ and $q\ge1$: for $f\in
\mathcal D$ with a splitting $(f_\alpha,f_\beta)$ along $H'$ and $k\in K_W$ there is $f_1\in\mathcal D(G_1)$ which is a low of
$f_\alpha$ at $k$.
\item\label{pf:thmK:E:twostep} In the situation of \ref{pf:thmK:E:alphalow}, if $(f_\alpha,0)$ is a splitting of $f$, then
every $f_1\in\mathcal D(G_1)$ which is a low of $f_\alpha$ at some $k\in K_W$ is pure of law $(\alpha,a'/2)$.
\item\label{pf:thmK:E:cusp} If $f\in\mathcal D$ has $\thmKCT_{H'}(f)=0$ for a plane $H'$, then $f$ is cuspidal.
\item\label{pf:thmK:E:orthiso} If $f\in\mathcal D$ is cuspidal, then $\thmKI(\Phi)\perp f$ for $\Phi\in\thmKPD(M)$ if $p\ge1$,
and $\thmKI'(\Phi')\perp f$ for $\Phi'\in\thmKPD(M')\cap\thmKSo(M')$ if $q\ge1$.
\item\label{pf:thmK:E:orthan} If $f\in\mathcal D$ is cuspidal, then $\thmKI(\Phi)\perp f$ for $\Phi\in\thmKPD(M)$ if $p=0<q$,
and $\thmKI'(\Phi')\perp f$ for $\Phi'\in\thmKPD(M')\cap\thmKSo(M')$ if $q=0<p$.
\item\label{pf:thmK:E:perp} If $f\in\mathcal D$ is orthogonal to all $\thmKI(\Phi)$, $\Phi\in\thmKPD(M)$, and to all
$\thmKI'(\Phi')$, $\Phi'\in\thmKPD(M')\cap\thmKSo(M')$, then $f=0$.
\end{enumerate}

We first derive the end steps from these statements.

\begin{lemma}\label{pf:thmK:zeroline}
Assume \ref{pf:thmK:E:cusp}--\ref{pf:thmK:E:perp}. If $p+q\ge1$ and $f\in\mathcal D$ has $\thmKCT_{H'}(f)=0$ for a plane
$H'$, then $f=0$.
\formallabel{PositivityTheory.ThKEnd.zero_of_line}
\end{lemma}

\begin{proof}
By \ref{pf:thmK:E:cusp} $f$ is cuspidal. If $p\ge1$, $\thmKI(\Phi)\perp f$ by \ref{pf:thmK:E:orthiso}; if $p=0$, then $q\ge1$
and $\thmKI(\Phi)\perp f$ by \ref{pf:thmK:E:orthan}. Similarly $\thmKI'(\Phi')\perp f$. By \ref{pf:thmK:E:perp}, $f=0$.
\end{proof}

We also use two elementary consequences of \ref{pf:thmK:E:iwa}--\ref{pf:thmK:E:aut}. (Z) If $T$ has some law $(\chi,z)$, is
left $N(H')$-invariant and $T(m(1,h_1)k)=0$ for all Levi elements $m(1,h_1)$ and all $k\in K_W$, then $T=0$: by
\ref{pf:thmK:E:iwa}, $T(u\,m(t,1)m(1,h_1)k)=\chi(t)|t|^zT(m(1,h_1)k)=0$. (S) If $f\in\mathcal D$ is pure of law
$(\alpha,a'/2)$ and $H'$ is adapted to $(W,W_1)$, then $(c,0)$ is a splitting of $f$ along $H'$ with
$c=\thmKCT_{H'}(f)$: $f$ is automorphic by \ref{pf:thmK:E:aut} and \ref{pf:thmK:E:ct}, so by \ref{pf:thmK:E:linemove} its
constant term along $H'$ exists and has law $(\alpha,a'/2)$, while $0$ has every law.

\begin{lemma}[Pure differences vanish]\label{pf:thmK:sumb}
Assume \ref{pf:thmK:E:iwa}--\ref{pf:thmK:E:perp} and the leaf \ref{pf:thmK:L:towerB}. For every triple with $q\ge1$, every
$f\in\mathcal D$ which is pure of law $(\alpha,a'/2)$ vanishes.
\formallabel{PositivityTheory.ThKEnd.sumB}
\end{lemma}

\begin{proof}
By induction on $p$. Right translation by $h_0$ and passing to a re-choice preserve purity of law $(\alpha,a'/2)$ (the law
of a constant term along $H'$ is preserved by right translations, and a re-choice has the same constant terms), and $f=0$ iff
$f(\cdot\,h_0)=0$; so by \ref{pf:thmK:E:sigre} we may replace $G$ by any re-choice which is part of a triple.

$p=0$. Write $a'=b+2$. By \ref{pf:thmK:L:towerB} there are a re-choice which is part of a triple and a $\beta$-step
configuration $(G,G_1,H,H')$ with $H'$ adapted to $(W,W_1)$, $W_1$ skew-hermitian. By (S), $f$ has the splitting $(c,0)$ along
$H'$, and \ref{pf:thmK:E:alphazero} gives $c=0$. So $\thmKCT_{H'}(f)=0$ and $f=0$ by \cref{pf:thmK:zeroline}.

$p\ge1$. Write $a=b+2$. By \ref{pf:thmK:E:tower2} we may assume that there is a plane-removal configuration $(G,G_1,H,H')$
with a triple on $G_1$; then $p(G_1)=p-1$ and $q(G_1)=q\ge1$. By (S), $f$ has the splitting $(c,0)$ along $H'$. By
\ref{pf:thmK:E:alphalow}, for every $k\in K_W$ there is a low $f_1\in\mathcal D(G_1)$ of $c$ at $k$; by
\ref{pf:thmK:E:twostep} it is pure of law $(\alpha,a'/2)$, so $f_1=0$ by induction. Hence $c(m(1,h_1)k)=0$ for all $h_1,k$.
By \ref{pf:thmK:E:ct} (the exponents $a'/2$ and $a/2$ are distinct) $c$ is left $N(H')$-invariant, so $c=0$ by (Z), and
$f=0$ by \cref{pf:thmK:zeroline}.
\end{proof}

\begin{proposition}\label{pf:thmK:endAB}
Assume \ref{pf:thmK:E:iwa}--\ref{pf:thmK:E:perp} and the leaves \ref{pf:thmK:L:splitex}, \ref{pf:thmK:L:matchex},
\ref{pf:thmK:L:towerB}. Then the leaves \ref{pf:thmK:L:endA} and \ref{pf:thmK:L:endB} hold.
\formallabel{PositivityTheory.ThKEnd.endA_of}\formallabel{PositivityTheory.ThKEnd.endB_of}
\formallabel{PositivityTheory.ThKG.endA'_of}\formallabel{PositivityTheory.ThKG.endB'_of}
\end{proposition}

\begin{proof}
\ref{pf:thmK:L:endA}. Let matched data with $q=0<p$, a plane $H'$ adapted to $(W,W_1)$ and $c$ be as in
\ref{pf:thmK:L:endA} (the hypothesis $\mathcal B(p+q)$ is not needed), and let $\Phi\approx\Phi'$ with $\Phi'\in\thmKSo(M')$.
By \ref{pf:thmK:L:matchex}(c), $\Phi\in\thmKPD(M)$ and $\Phi'\in\thmKPD(M')$, so $f=\thmKI(\Phi)-c\,\thmKI'(\Phi')=
\thmKI(\Phi)-\thmKI'(c\Phi')\in\mathcal D$. By \ref{pf:thmK:L:splitex} there are splittings $(f_\alpha,f_\beta)$ of
$\thmKI(\Phi)$ and $(f'_\alpha,f'_\beta)$ of $\thmKI'(\Phi')$, and by \ref{pf:thmK:E:aut} and the linearity in
\ref{pf:thmK:E:ct}, $f$ has the splitting $(T_1,T_2)=(f_\alpha-cf'_\alpha,f_\beta-cf'_\beta)$. By \ref{pf:thmK:E:ct} ($a'\ne
a$) $T_1$ is left $N(H')$-invariant. For $k\in K_W$, $\omega(1,k)\Phi\approx\omega'(1,k)\Phi'$ with second member in
$\thmKSo(M')$ \ref{pf:thmK:E:stab}; their theta integrals are the right translates by $k$, whose splittings are the right
translates of the splittings above; so the hypothesis of \ref{pf:thmK:L:endA} gives $T_1(m(1,h_1)k)=0$ for all $h_1$ and
$k$, and $T_1=0$ by (Z). By \ref{pf:thmK:E:beta}, $T_2=0$. Thus $\thmKCT_{H'}(f)=0$, and \cref{pf:thmK:zeroline} gives
$f=0$, i.e.\ $\thmKI(\Phi)=c\,\thmKI'(\Phi')$: $c$ is a polynomial constant.

\ref{pf:thmK:L:endB}. With $q\ge1$ the same argument, applied to the second components, gives $T_2=0$. So
$\thmKCT_{H'}(f)=T_1$ has law $(\alpha,a'/2)$, $f$ is pure of law $(\alpha,a'/2)$, and $f=0$ by \cref{pf:thmK:sumb}.
\end{proof}

\begin{theorem}[End steps]\label{pf:thmK:endsteps}
Assume \ref{pf:thmK:H:dict}, \ref{pf:thmK:H:finram}, \ref{pf:thmK:H:finsplit}, \ref{pf:thmK:H:fininert},
\ref{pf:thmK:H:lie2}, \ref{pf:thmK:H:archlie} and \ref{pf:thmK:H:siegel}--\ref{pf:thmK:H:weilex} of
\cref{pf:thmK:closesph}. Then \ref{pf:thmK:E:iwa}--\ref{pf:thmK:E:perp} hold, and the leaves \ref{pf:thmK:L:endA} and
\ref{pf:thmK:L:endB} hold.
\formallabel{PositivityTheory.Sph.Tr.endAB_sph}
\end{theorem}

\begin{proof}
The fifteen statements are proved as follows; most of the auxiliary leaves used below are among those established in
\cref{sec:pf:thmK:leafproofs} from the same hypotheses.
\begin{itemize}
\item \ref{pf:thmK:E:iwa} is \lref{PositivityTheory.ThKG.Tow.iwa}: an Iwasawa decomposition at the infinite places with
respect to the parabolic subgroup of $H'$, obtained from frames adapted to $H'$ and to the Cartan involutions.
\item \ref{pf:thmK:E:ct} is \lref{PositivityTheory.ThKEnd.Grp.ctBasic}. Linearity holds because constant terms are
integrals over the compact quotient $N(H')/(N(H')\cap\Gamma_W)$ of functions which are continuous on it. For the invariance,
fix $u\in N(H')$ and put $\phi_i(h)=T_i(uh)-T_i(h)$; then $\phi_1=-\phi_2$ since $\thmKCT_{H'}(f)$ is left
$N(H')$-invariant, and for $m=m(t,1)$ the laws give $\phi_i(mh)=\chi_i(t)|t|^{z_i}\phi_i'(h)$, where $\phi'_i$ is defined
with $m^{-1}um$ in place of $u$. Choosing $t$ with $|t|^{z_1}\ne|t|^{z_2}$ forces $\phi'_1=0$, and $m^{-1}um$ runs over
$N(H')$ with $u$.
\item \ref{pf:thmK:E:linemove}, \ref{pf:thmK:E:aut}, \ref{pf:thmK:E:stab}, \ref{pf:thmK:E:sigre} are
\lref{PositivityTheory.ThKEnd.Grp.lineMove}, \lref{PositivityTheory.ThKEnd.Grp.iAut_of} (continuity and
invariance of theta integrals, \cref{pf:thmK:analytic}), \lref{PositivityTheory.ThKEnd.Grp.matchStab_of} (the
archimedean operators of $K_W$ preserve the Fock spaces and the matching, and $\thmKSo(M')$ is stable) and
\lref{PositivityTheory.ThKEnd.Grp.sigRe_of} (Cartan involutions of the same signature are conjugate under
$\Uni(W)(F_\infty)$, and conjugating the Cartan data moves the Fock spaces, \lref{PositivityTheory.ThKF.cartanMove_of}).
\item \ref{pf:thmK:E:tower2} is \lref{PositivityTheory.ThKG.Tow.towerCodimensionTwo_of}, proved like the towers
\cref{pf:thmK:towers} from the existence of adapted planes, of good models and of Weil representations
(\cref{pf:thmK:weilex}).
\item \ref{pf:thmK:E:alphazero}--\ref{pf:thmK:E:twostep} are \lref{PositivityTheory.ThKEnd.Comp.alphaZero_of},
\lref{PositivityTheory.ThKEnd.Comp.betaZero_of}, \lref{PositivityTheory.ThKEnd.Comp.alphaLow_of} and
\lref{PositivityTheory.ThKEnd.Two.twoStep_of'}. The first three compute the components of the constant terms of
elements of $\mathcal D$ with the general constant-term formulas (\cref{pf:thmK:gen}), the factorisation of the cone integral
(\cref{pf:thmK:conefac}) and a vanishing statement for real theta lifts
(\lref{PositivityTheory.ThKEnd.Comp.realVanish_of}); the fourth combines the moderate growth of elements of $\mathcal D$
with the existence of splittings along a second plane (\lref{PositivityTheory.ThKEnd.Two.two_kill}).
\item \ref{pf:thmK:E:cusp} is \lref{PositivityTheory.ThKEnd.Cusp.cuspOfLine_of_keys}.
\item \ref{pf:thmK:E:orthiso} is \lref{PositivityTheory.ThKEnd.Orth.orthIso_fin}; it uses the rapid decay of cuspidal
elements of $\mathcal D$ (\lref{PositivityTheory.ThKEst.sigRapid}) and a constancy statement for theta lifts
(\lref{PositivityTheory.ThKEnd.Orth.sigThConst_fin}, which uses the non-vanishing of Hecke $L$-functions at $s=1$,
\cref{pf:thmK:hecke}).
\item \ref{pf:thmK:E:orthan} is \lref{GlobalTheory.ThK.Dbl.orthAn_of}; it uses the doubling method with the
Siegel--Weil formula and the meromorphic continuation (\cref{pf:thmK:sw:thm,pf:thmK:eis}), the rapid decay of cuspidal
elements of $\mathcal D$, and the vanishing of the constant terms of the theta lifts of cuspidal elements
(\lref{GlobalTheory.ThK.Dbl.cusp_lift}).
\item \ref{pf:thmK:E:perp}: Let $f=\thmKI(\Phi_0)-\thmKI'(\Phi'_0)\in\mathcal D$ be orthogonal to all theta integrals as
stated. All these theta integrals are square integrable over $\mathcal F$
(\lref{PositivityTheory.ThKEst.Sq.sqInt_of_leaves}), hence
$\int_{\mathcal F}|f|^2\,d\mu=\int_{\mathcal F}\thmKI(\Phi_0)\bar f\,d\mu-\int_{\mathcal F}\thmKI'(\Phi'_0)\bar f\,d\mu=0$.
As $f$ is automorphic \ref{pf:thmK:E:aut}, it vanishes (as in the proof of \cref{pf:thmK:kills}).
\end{itemize}
The leaves \ref{pf:thmK:L:splitex}, \ref{pf:thmK:L:matchex} and \ref{pf:thmK:L:towerB} hold by
\cref{pf:thmK:ctleaves,pf:thmK:matchex,pf:thmK:towers} (the coinvariant condition entering \ref{pf:thmK:L:matchex} follows from
\ref{pf:thmK:H:dict} as in \cref{pf:thmK:closesph}). \Cref{pf:thmK:endAB} gives the leaves \ref{pf:thmK:L:endA} and
\ref{pf:thmK:L:endB}.
\end{proof}

% ===========================================================================
\subsection{Proofs of the deferred results}\label{sec:pf:thmK:deferred}
This subsection proves the results which were used above through their formal statements: the analytic statements of the
doubling method and of the cone integrals, the local multiplicity-one and matching statements, the vanishing of theta
integrals of killed tensors, the existence of Weil representations, and the end steps. The references
\lref{GlobalTheory.ThK.Dbl.lInt2}, \lref{PositivityTheory.ThKG.Loc.Cone.ideleFub}, \dots\ used in the previous
subsections refer to the results proved here.

\subsubsection{Partial $L$-functions at integers}\label{sec:pf:thmK:lint}

\begin{littheorem}[Hecke; \cite{Hecke1920,Tate1967thesis}]\label{pf:thmK:heckefe}
Let $K$ be a totally complex number field and $\mu$ a unitary Hecke character of $K$ which is not of the form
$|\cdot|^{i\tau}$, $\tau\in\R$, with $\mu_w(z)=(z/|z|)^{m_w}$ on $K_w^\times\cong\C^\times$ for the infinite places $w$. Let
$L(s,\mu)$ be the Euler product over the finite places ($\Re s>1$), $D=|d_K|\Nm\mathfrak f_\mu$, $\Gamma_\C(s)=
2(2\pi)^{-s}\Gamma(s)$ and $\Gamma_\mu(s)=\prod_w\Gamma_\C(s+|m_w|/2)$. Then $\Lambda(s,\mu)=D^{s/2}\Gamma_\mu(s)L(s,\mu)$
extends to an entire function; $\Lambda(s,\mu)=W(\mu)\Lambda(1-s,\mu^{-1})$ with $W(\mu)W(\mu^{-1})=1$; and an entire
function which agrees with $L(s,\mu)$ on $\Re s>1$ has no zero on $\Re s\ge1$.
\formal{GroundTheory.Global.HeckeCharacter.differentiable_completedLFunction}
\formal{GroundTheory.Global.HeckeCharacter.completedLFunction_eq_rootNumber_mul_of_not}
\formal{GroundTheory.Global.HeckeCharacter.continuation_ne_zero_of_one_le_re}
\end{littheorem}

\begin{lemma}[Partial $L$-functions at integers]\label{pf:thmK:lint}
For a triple and $\chi\in\{\nu,\nu^{-1}\}$ there is an entire function $\Lambda$ with $\Lambda(s)=L^{\emptyset}(s,\chi)$ for
$\Re s>1$ and $\Lambda(z)\ne0$ for every integer $z$ with $2z\ge-(a+a')$.
\formallabel{GlobalTheory.ThK.Dbl.lInt2}
\end{lemma}

\begin{proof}
Since $B$ is good, $E$ is a CM field, hence totally complex, and the embeddings $\tau_w$ ($w$ an infinite place of $F$) form
a CM type: every embedding $E\to\C$ is $\tau_w$ or $\overline{\tau_w}$ for exactly one $w$. The character $\chi$ is a unitary
Hecke character of $E$ whose infinity type is $\pm(m_\alpha-m_\beta)$, computed through the $\tau_w$. Let $\mu$ be $\chi$
viewed as a Hecke character in the sense of \cref{pf:thmK:heckefe}; identifying the two notions (same local components,
same Euler factors) is a matter of definitions. Each infinite place of $E$ is induced by exactly one $\tau_w$ up to
conjugation, so $\mu$ has an infinity type $(m_{w'})$ with $|m_{w'}|=|m_\alpha(w)-m_\beta(w)|>a+a'$ for the place $w'$
induced by $\tau_w$, by the conditions on a good configuration. In particular $\mu$ is not of the form $|\cdot|^{i\tau}$
(whose infinity type is $0$), and $L(s,\mu)=L^\emptyset(s,\chi)$ for $\Re s>1$ (both are the products of the same local
factors).

Let $G(s)=D^{-s/2}\Gamma_\mu(s)^{-1}\Lambda(s,\mu)$. As $1/\Gamma$ is entire, $G$ is entire, and $G(s)=L(s,\mu)$ for
$\Re s>1$. Let $z$ be an integer with $2z\ge-(a+a')$; then $2z+|m_{w'}|>0$ for all $w'$. If $z\ge1$, then $G(z)\ne0$ by the
last assertion of \cref{pf:thmK:heckefe}. If $z\le0$, then $\Gamma_\mu(z)=\prod_{w'}2(2\pi)^{-z-|m_{w'}|/2}
\Gamma(z+|m_{w'}|/2)$ is finite and non-zero because $z+|m_{w'}|/2>0$; $W(\mu)\ne0$; and, writing $G^-$ for the function
attached in the same way to $\mu^{-1}$ (which is not of the form $|\cdot|^{i\tau}$ either), $G^-(1-z)\ne0$ since
$1-z\ge1$, so $\Lambda(1-z,\mu^{-1})\ne0$. Hence $G(z)=D^{-z/2}\Gamma_\mu(z)^{-1}W(\mu)\Lambda(1-z,\mu^{-1})\ne0$. Take
$\Lambda=G$.
\end{proof}

\subsubsection{An idelic Fubini theorem}\label{sec:pf:thmK:idelefub}
For a finite place $v$ of $F$ put $E_v=E\otimes_FF_v=\prod_{w\mid v}E_w$, and $E_\infty=E\otimes_F F_\infty$; we view
$E_v^\times$ and $E_\infty^\times$ as subgroups of $\thmKAE^\times$ (ideles supported at $v$, resp.\ at the infinite
places). A \emph{system of local measures} consists of Haar measures $\mu_\infty$ on $E_\infty^\times$ and $\mu_v$ on
$E_v^\times$ with $\mu_v(\OO_{E,v}^\times)=1$ (and the analogous measures on the local components of a compact subgroup,
used in \cref{sec:pf:thmK:conefac}). For a non-zero fractional ideal $\mathfrak a$ of $E$ let $\mathfrak a_v=\mathfrak
a\OO_{E,v}$, and for a character $\nu$ of $\thmKAE^\times$ unramified at all $w$ with $\ord_w\mathfrak a\ne0$ put
$\nu(\mathfrak a)=\prod_w\nu_w(\varpi_w)^{\ord_w\mathfrak a}$. Let $|\cdot|$ be the idelic absolute value of $E$.

\begin{lemma}[Idelic Fubini theorem]\label{pf:thmK:idelefub}
Let $B$ be good, $(\mu_\infty,\mu_v)$ a system of local measures and $\mu$ a Haar measure on $\thmKAE^\times$. Let
$\mathfrak a\ne0$ be a fractional ideal of $E$ and $z\ge2$ an integer. There is $\kappa\ne0$ such that for every continuous
character $\nu$ of $\thmKAE^\times$ and all finite sets $S_0\subseteq S$ of finite places such that $\nu$ is unramified at
the places above $v\notin S$, and $\nu$ is ramified and $\ord_w\mathfrak a=0$ at every place $w$ above $S$:
\begin{enumerate}
\item for $v\notin S$ the integral $u_v=\int_{E_v^\times}\one_{\mathfrak a_v}(l)\,\nu(l)|l|^z\,d\mu_v(l)$ converges
absolutely and is non-zero;
\item for every finite $T\supseteq S$ and all functions $g_\infty$ on $E_\infty^\times$ and $g_v$ on $E_v^\times$
($v\in T$) such that $g_\infty\nu|\cdot|^z$ and the $g_v\nu|\cdot|^z$ are integrable, the function
$\Psi(l)=g_\infty(l_\infty)\prod_{v\in T}g_v(l_v)\prod_{v\notin T}\one_{\mathfrak a_v}(l_v)\cdot\nu(l)|l|^z$ is
$\mu$-integrable and
\begin{equation}\label{pf:thmK:eq:idelefub}
\Bigl(\prod_{v\in T\smallsetminus S}u_v\Bigr)\int\Psi\,d\mu=\kappa\,\nu(\mathfrak a)\,L^{S_0}(z,\nu)\int g_\infty\nu
|\cdot|^z\,d\mu_\infty\prod_{v\in T}\int g_v\nu|\cdot|^z\,d\mu_v .
\end{equation}
\end{enumerate}
\formallabel{PositivityTheory.ThKG.Loc.Cone.ideleFub}
\end{lemma}

\begin{proof}
\emph{Haar measures.} For a finite set $T$ of finite places, $l\mapsto(l_\infty,(l_v)_{v\in T},l^T)$ is an isomorphism of
topological groups from $\thmKAE^\times$ onto $E_\infty^\times\times\prod_{v\in T}E_v^\times\times\thmKAE^{T,\times}$,
where $\thmKAE^{T,\times}$ is the restricted product of the $E_v^\times$, $v\notin T$, with respect to the unit groups. Let
$\mu^T$ be the Haar measure on $\thmKAE^{T,\times}$ giving $\prod_{v\notin T}\OO_{E,v}^\times$ measure $1$. By uniqueness
of Haar measures there is $c>0$ with $\mu=c\,(\mu_\infty\times\mu^\emptyset)$, and $\mu^\emptyset=\prod_{v\in T}\mu_v\times
\mu^T$ for every $T$ (both sides are Haar measures and give $\prod_v\OO_{E,v}^\times$ measure $1$). Since $\nu$ and $|\cdot|$
are multiplicative, $\Psi$ is a product of functions of $l_\infty$, of the $l_v$ ($v\in T$) and of $l^T$, and Fubini's
theorem gives
\begin{equation}\label{pf:thmK:eq:fubsplit}
\int\Psi\,d\mu=c\int g_\infty\nu|\cdot|^z\,d\mu_\infty\prod_{v\in T}\int g_v\nu|\cdot|^z\,d\mu_v\cdot
\int_{\thmKAE^{T,\times}}\prod_{v\notin T}\one_{\mathfrak a_v}(l_v)\,\nu(l)|l|^z\,d\mu^T(l),
\end{equation}
as soon as the last integral converges absolutely.

\emph{Local integrals.} Let $v\notin S$. Then $E_v^\times=\prod_{w\mid v}E_w^\times$, $\mu_v$ is the product of the Haar
measures giving $\OO_w^\times$ measure $1$, and $E_w^\times$ is the disjoint union of the shells $\varpi_w^k\OO_w^\times$,
$k\in\Z$, on which $\nu_w$ is constant ($\nu$ is unramified at $w$) and $|\cdot|$ equals $\Nm w^{-k}$. Hence
\[
u_v=\prod_{w\mid v}\sum_{k\ge\ord_w\mathfrak a}\bigl(\nu_w(\varpi_w)\Nm w^{-z}\bigr)^k
=\prod_{w\mid v}\bigl(\nu_w(\varpi_w)\Nm w^{-z}\bigr)^{\ord_w\mathfrak a}L_w(z,\nu),
\]
a product of absolutely convergent geometric series ($|\nu_w(\varpi_w)\Nm w^{-z}|<1$) with non-zero sums. This proves (1).

\emph{The tail.} Let $T_0\supseteq T$ contain all places below the prime divisors of $\mathfrak a$. On
$\thmKAE^{T_0,\times}$ the integrand is $\prod_{v\notin T_0}\one_{\OO_{E,v}}(l_v)\nu(l)|l|^z$. Its absolute value is the
increasing limit, over finite sets $T_0\subseteq T_1$, of the functions $\prod_{v\in T_1\smallsetminus T_0}
\one_{\OO_{E,v}}(l_v)|l_v|^z\cdot\prod_{v\notin T_1}\one_{\OO_{E,v}^\times}(l_v)$, whose integrals are
$\prod_{w}(1-\Nm w^{-z})^{-1}$ over the $w$ above $T_1\smallsetminus T_0$; these are bounded by $\zeta_E(z)<\infty$ as
$z\ge2$. By monotone and dominated convergence the integral converges absolutely and equals $\prod_{v\notin T_0}u_v$ (an
absolutely convergent product). By Fubini over the finitely many places of $T_0\smallsetminus T$ the integral over
$\thmKAE^{T,\times}$ in \eqref{pf:thmK:eq:fubsplit} converges absolutely and equals $\prod_{v\notin T}u_v$.

\emph{Conclusion.} Multiplying by $\prod_{v\in T\smallsetminus S}u_v$ and using the formula for $u_v$,
\[
\Bigl(\prod_{v\in T\smallsetminus S}u_v\Bigr)\prod_{v\notin T}u_v=\prod_{v\notin S}u_v
=\prod_{w}\Nm w^{-z\ord_w\mathfrak a}\cdot\nu(\mathfrak a)\cdot L^S(z,\nu),
\]
where we used $\ord_w\mathfrak a=0$ above $S$. Since $\nu$ is ramified at every place above $S$, its local factors there
are $1$ and $L^S(z,\nu)=L^{S_0}(z,\nu)$. With $\kappa=c\prod_w\Nm w^{-z\ord_w\mathfrak a}\ne0$, which depends only on
$\mu$, the system of measures, $\mathfrak a$ and $z$, \eqref{pf:thmK:eq:fubsplit} gives \eqref{pf:thmK:eq:idelefub}.
\end{proof}

\subsubsection{Theta integrals of killed pure tensors}\label{sec:pf:thmK:kills}
For a Weil pair, a finite place $v$ and $x\in\thmKS(F_v^{n\times d})$, let $\Phi[v\mapsto x]$ denote the function obtained
from a pure tensor $\Phi=\thmKpure(f,\varphi)$ by replacing its component at $v$ by $x$; if $\omega(g,1)\colon x\thmKacts{v}
x_1$ for $g\in\Uni(X)(F_v)$, then $\omega(g,1)\Phi[v\mapsto x]=\Phi[v\mapsto x_1]$
(\lref{PositivityTheory.ThKF.intersection_update_fsHerm_neg_one}). A subspace $N\subseteq\thmKS(F_v^{n\times d})$ is a
\emph{theta subspace at $v$} if it contains $\thmKaug_v(\omega)$ and is stable under the local action of $\Uni(W)(F_v)$.
Similarly, a subspace $N$ of $\thmKPol(M)$ is an \emph{archimedean theta subspace} if it contains the derivatives
$\frac{d}{dt}\omega(\gamma(t),1)f|_{t=0}$ ($f\in\thmKPol(M)$, $\gamma$ a one-parameter subgroup of $X$) and the differences
$\omega(k,1)f-f$ ($k\in K_\infty(X)$), and is stable under the derivatives along the one-parameter subgroups of $W$ and
under $K_\infty(W)$; the span of these derivatives and differences is denoted $\thmKaug_\infty$.

\begin{proposition}[Killed pure tensors]\label{pf:thmK:kills}
Let matched data be given, and assume the conclusion of \ref{pf:thmK:H:so} and the statements \ref{pf:thmK:H:sph},
\ref{pf:thmK:H:lie2}, \ref{pf:thmK:H:archlie} of \cref{pf:thmK:closesph}, and the coinvariant condition.
\begin{enumerate}
\item Let $f\in\thmKPol(M)$ and $\varphi=(\varphi_v)$ with $\varphi_v\in\thmKS$, $\varphi_v=\thmKvac_v$ for almost all
$v$. If $p(f)=0$ or $p_v(\varphi_v)=0$ for some $v$, then $\thmKI(\thmKpure(f,\varphi))=0$.
\item Let $f\in\thmKPol(M')$ and $\varphi$ as in (1) for $M'$ with $\varphi_{v_1}$ in the span of $\thmKSo_{v_1}(\tau_1)$.
If $p'(f)=0$ or $p'_v(\varphi_v)=0$ for some $v$, then $\thmKI'(\thmKpure(f,\varphi))=0$.
\end{enumerate}
\formallabel{PositivityTheory.Sph.theta_kills_cusp_forms_spherical}
\end{proposition}

\begin{proof}
We treat (1); the differences for (2) are indicated at the end. Fix a fundamental domain $\mathcal F$ of $\Gamma_W$ in
$\Uni(W)(\A)$ and a Haar measure $\mu_W$, which is right invariant as $\Uni(W)(\A)$ is unimodular
(\lref{PositivityTheory.ThKF.uniW}). For every test function $\Phi$, $\thmKI(\Phi)$ is continuous and left
$\Gamma_W$-invariant, and for $\Phi\in\thmKPD(M)$ it is square integrable over $\mathcal F$; for (2) the same holds for
$\Phi\in\thmKSo(M')$ (\lref{PositivityTheory.ThKF.clsV}, \lref{PositivityTheory.ThKF.clsV'}: these follow from
\ref{pf:thmK:L:weilconv}, \ref{pf:thmK:L:socont}, the decay condition, the coinvariant condition and the square
integrability of theta integrals, \lref{PositivityTheory.ThKEst.Sq.sqInt_of_leaves}).

\emph{The positivity argument.} Let $Q$ be a space of local components at a place (below, $Q=\thmKS(F_v^{n\times d})$ or
$Q=\thmKPol(M)$) and $u\colon Q\to\{\text{functions on }\Uni(W)(\A)\}$ a linear map such that each $u(x)$ is continuous,
left $\Gamma_W$-invariant and square integrable on $\mathcal F$. Then
\[
\beta(x,y)=\int_{\mathcal F}\overline{u(x)(h)}\,u(y)(h)\,d\mu_W(h)
\]
is a positive semi-definite hermitian form on $Q$, and $\beta(x,x)=0$ implies $u(x)=0$: a continuous left
$\Gamma_W$-invariant function vanishing almost everywhere on $\mathcal F$ vanishes almost everywhere on
$\Gamma_W\mathcal F=\Uni(W)(\A)$ ($\Gamma_W$ is countable), hence everywhere, since Haar measure charges every non-empty
open set. If $A$ is a linear map of $Q$ and $h_0\in\Uni(W)(\A)$ with $u(Ax)(h)=u(x)(hh_0)$ for all $x,h$, then
$\beta(Ax,Ay)=\beta(x,y)$, because $\mathcal Fh_0$ is again a fundamental domain and $\mu_W$ is right invariant. If $D$ is a
linear map of $Q$ and $\delta$ a one-parameter subgroup with $\frac{d}{dt}u(x)(h\delta(t))|_{t=0}=u(Dx)(h)$ for all $x,h$,
then $\beta(Dx,y)+\beta(x,Dy)=0$: the function $t\mapsto\int_{\mathcal F}\overline{u(x)(h\delta(t))}u(y)(h\delta(t))\,
d\mu_W(h)$ is constant by right invariance, and its derivative at $0$ may be computed under the integral sign
(\lref{PositivityTheory.ThKF.L2.skew}, which justifies this from the continuity and square integrability of the
functions involved).

\emph{Finite places.} Suppose $p_v(\varphi_v)=0$, let $Q=\thmKS(F_v^{n\times d})$ and $u(x)=\thmKI(\thmKpure(f,\varphi)
[v\mapsto x])$. By the above, $u$ has the required properties. Let $\omega_v$ be a local representation of
$\Uni(V)(F_v)\times\Uni(W)(F_v)$ on $Q$ compatible with $\omega$ (\lref{PositivityTheory.ThKH.Locrep.locRepEx}). If
$\omega(g,1)\colon x\thmKacts{v}x_1$, then $u(x_1)=\thmKI(\omega(g,1)\Phi[v\mapsto x])=u(x)$ by the left invariance of
the theta integral (\lref{PositivityTheory.ThKF.iLeft}); so $u$ vanishes on $\thmKaug_v(\omega)$, and $\beta(x,y)=0$
when $x$ or $y$ lies in $\thmKaug_v(\omega)$. For $h\in\Uni(W)(F_v)$, $u(\omega_v(1,h)x)(g)=u(x)(gh)$, hence $\beta$ is
invariant under $\omega_v(1,h)$. The subspace $N=\ker p_v\cap Q$ is a theta subspace at $v$ (as $p_v$ kills the
augmentation module and is equivariant) and $N\ne Q$ (as $p_v(Q)=\Pi_v\ne0$). By \ref{pf:thmK:H:sph} the local Weil
representation at $v$ is irreducible or spherical.

If it is irreducible, $N=\thmKaug_v(\omega)$, so $\varphi_v\in\thmKaug_v(\omega)$ and $u(\varphi_v)=0$.

If it is spherical, there is $s\in Q$ such that (S1) $Q=\thmKaug_v(\omega)+\operatorname{span}\{\omega_v(1,h)s:h\in
\Uni(W)(F_v)\}$, and (S2) for every $y\in Q$ there are $k_0,\dots,k_m\in\Uni(W)(F_v)$, $w_0,\dots,w_m\in\C$ and $c\in\C$
with $w_i\,\omega_v(1,k_i)s=s$ and $\sum_iw_i\,\omega_v(1,k_i)y-cs\in\thmKaug_v(\omega)$. Then $s\ne0$ (otherwise
$Q=\thmKaug_v(\omega)\subseteq N$ by (S1)) and $s\notin N$ (otherwise $Q\subseteq N$ by (S1), as $N$ is a theta
subspace). Let $y\in N$ and choose $k_i,w_i,c$ as in (S2). Since $N$ is stable and contains the augmentation module,
$cs\in N$, so $c=0$, and $z=\sum_iw_i\,\omega_v(1,k_i)y\in\thmKaug_v(\omega)$. As $s=w_i\,\omega_v(1,k_i)s$, invariance
gives $\beta(s,\omega_v(1,k_i)y)=\overline{w_i}\,\beta(\omega_v(1,k_i)s,\omega_v(1,k_i)y)=\overline{w_i}\,\beta(s,y)$, hence
\[
0=\beta(s,z)=\Bigl(\sum_i|w_i|^2\Bigr)\beta(s,y).
\]
Some $w_i$ is non-zero (as $s\ne0$), so $\beta(s,y)=0$. Since $N$ is stable, $\beta(\omega_v(1,h)s,y)=\beta(s,
\omega_v(1,h^{-1})y)=0$ for all $h$, and by (S1) and the vanishing of $\beta$ on the augmentation module, $\beta(x,y)=0$
for all $x\in Q$ and $y\in N$. In particular $\beta(\varphi_v,\varphi_v)=0$, so $\thmKI(\thmKpure(f,\varphi))=u(\varphi_v)
=0$.

\emph{The archimedean place.} Suppose $p(f)=0$, let $Q=\thmKPol(M)$ and $u(x)=\thmKI(\thmKpure(x,\varphi))$. For
$k\in K_\infty(W)$, $\omega(1,k)$ preserves $Q$ and $u(\omega(1,k)x)(h)=u(x)(hk)$; for a one-parameter subgroup $\delta$ of
$W$ the derivative $D_\delta$ of $t\mapsto\omega(1,\delta(t))$ preserves $Q$ and $\frac{d}{dt}u(x)(h\delta(t))|_{t=0}=
u(D_\delta x)(h)$ (\lref{PositivityTheory.ThKF.A.archOp} and \ref{pf:thmK:H:archlie}). Similarly $u$ vanishes on
$\thmKaug_\infty$, by the left invariance of theta integrals and its derivative. Hence $\beta$ vanishes on
$\thmKaug_\infty$, is $K_\infty(W)$-invariant and satisfies $\beta(D_\delta x,y)+\beta(x,D_\delta y)=0$; consequently the
radical $\ker\beta$ is an archimedean theta subspace. The space $N=\ker p\cap Q$ is a maximal proper archimedean theta
subspace (\lref{PositivityTheory.ThKF.maxInfP}: this follows from the matching properties of $p$ and the irreducibility
of $\Pi_\infty$). We use two further facts. (R) For a form $\beta$ as above and a subspace $P\supseteq\ker\beta$ stable
under all $D_\delta$, either $P\subseteq\ker\beta$ or $P=Q$ (\lref{PositivityTheory.ThKF.Unit.radMax_of}, a
consequence of the explicit action of the Lie algebra on the Fock space, \ref{pf:thmK:H:lie2}). (J) The sum of two proper
archimedean theta subspaces is proper (\lref{PositivityTheory.ThKR.Join.joinTri_of}). If $\ker\beta=Q$ there is
nothing to prove. Otherwise apply (R) to $P=N+\ker\beta$: $P=Q$ is impossible by (J), so $N\subseteq\ker\beta$. Thus
$f\in\ker\beta$, $\beta(f,f)=0$ and $\thmKI(\thmKpure(f,\varphi))=u(f)=0$.

\emph{The case (2).} The same arguments apply to $M'$ and $\omega'$, with $u$ built from functions in $\thmKSo(M')$:
replacing the component at a place $v\ne v_1$, or the archimedean component, of a pure tensor whose $v_1$-component lies in
the span of $\thmKSo_{v_1}(\tau_1)$ gives again an element of $\thmKSo(M')$. It remains to treat $p'_{v_1}(\varphi_{v_1})=0$.
The place $v_1$ is good, so by the coinvariant condition $\varphi_{v_1}\in\thmKaug_{v_1}(\omega')$, and by the conclusion of
\ref{pf:thmK:H:so}, $\varphi_{v_1}$ lies in the span of the differences $y-x$ with $x$ in the span of
$\thmKSo_{v_1}(\tau_1)$ and $\omega'(g,1)\colon x\thmKacts{v_1}y$. For such $x,y$ the functions $\Phi[v_1\mapsto x]$ and
$\Phi[v_1\mapsto y]=\omega'(g,1)\Phi[v_1\mapsto x]$ lie in $\thmKSo(M')$ (which is stable under $\omega'$), and their theta
integrals agree by left invariance. By linearity $\thmKI'(\Phi[v_1\mapsto\varphi_{v_1}])=0$.
\end{proof}

\subsubsection{Factorisation of the cone integral}\label{sec:pf:thmK:conefac}
Let $K$ be the stabiliser of $X$. Its elements supported at a finite place $v$, resp.\ at the infinite places, form compact
groups $K_v$, $K_\infty$; a \emph{system of local measures} for $K$ (and $\thmKAE^\times$) consists of Haar measures
$\mu_{K,v}$ of total mass $1$ on the $K_v$, a Haar measure $\mu_{K,\infty}$ on $K_\infty$, and measures $\mu_\infty,\mu_v$ on
$E_\infty^\times,E_v^\times$ as in \cref{sec:pf:thmK:idelefub}. For an evaluation map $e$ with local components
$(e_\infty,(e_v))$, a map $u$ as in \eqref{pf:thmK:eq:cone}, a character $\nu$ and $z\in\N$, the \emph{local cone integrals}
are
\begin{align*}
\Lambda_v(\varphi)&=\int_{K_v}\int_{E_v^\times}e_v\bigl(\rho_\psi(u(l))\omega(k,1)\varphi\bigr)\nu(l)|l|^z\,d\mu_v(l)\,
d\mu_{K,v}(k),\\
L(f)&=\int_{K_\infty}\int_{E_\infty^\times}e_\infty\bigl(\rho_\psi(u(l))\omega(k,1)f\bigr)\nu(l)|l|^z\,d\mu_\infty(l)\,
d\mu_{K,\infty}(k),
\end{align*}
for $\varphi\in\thmKS(F_v^{n\times d})$ and archimedean Schwartz functions $f$, where the operators act through the local
actions at $v$, resp.\ at the infinite places.

\begin{proposition}[Factorisation of the cone integral]\label{pf:thmK:conefac}
There are signs $e_1,e_2$ with the following property. Let $(G,G_1,H,H')$ be an $\alpha$-step configuration,
$e\in\thmKEv(M,M_1;\iota_H,\iota_{H'};I_A(H,H'))$ with local components $(e_\infty,(e_v))$, $K$ the stabiliser of $V$ with a
Haar measure $\mu_K$, $\mu$ a Haar measure on $\thmKAE^\times$ and a system of local measures. There is $\kappa\ne0$ such
that for every pair of triples on $G$ and $G_1$ (with the same character datum), with $\nu=\alpha/\beta$, $z=k_A$ and
$u=u_A$:
\begin{enumerate}
\item $\thmKcone_A$ is linear on test functions;
\item $L$ is a continuous linear map of Schwartz spaces with $L(\thmKPol(M))\subseteq\thmKPol(M_1)$; the local cone
integrals $\Lambda_v$ converge and map $\thmKS$ into $\thmKS$; and for $v\notin T$ there is $u_v\ne0$ with
$\Lambda_v(\thmKvac_v)=u_v\thmKvac^{M_1}_v$;
\item for every finite $T'\supseteq T$ and every pure tensor $\thmKpure(f,\varphi)$ with $\varphi_v\in\thmKS$ for all $v$
and $\varphi_v=\thmKvac_v$ for $v\notin T'$,
\begin{multline}\label{pf:thmK:eq:conefac}
\Bigl(\prod_{v\in T'\smallsetminus T}u_v\Bigr)\thmKcone_A(\thmKpure(f,\varphi))\\
=\kappa\,\nu(\mathfrak b_H)^{e_1}\nu(\mathfrak b_{H'})^{e_2}\,\thmKLS{k_A}{\nu}\;\thmKpure\bigl(Lf,(\Lambda_v\varphi_v)_{v\in
T'},(\thmKvac^{M_1}_v)_{v\notin T'}\bigr).
\end{multline}
\end{enumerate}
The same holds, with other signs, for $\beta$-step configurations ($e\in\thmKEv(M',M'_1;\iota_H,\iota_{H'};I_B(H,H'))$,
the cone integral $\thmKcone_B$ on the $V'$-side, $K$ the stabiliser of $V'$, $u=u_B$, $z=k_B+1$, and $L^{S_0}(k_B+1,\nu)$
in \eqref{pf:thmK:eq:conefac}), and for the plane-removal configurations of \cref{sec:pf:thmK:endsteps}.
\formallabel{PositivityTheory.ThKG.Loc.Cone.coneFac_all}
\end{proposition}

\begin{proof}
We treat $\alpha$-steps; the other two cases are proved in the same way (for $\beta$-steps with $u_B$ and the ideal
identity below for $u_B$; the plane-removal case contains the $\alpha$-step case as a special case). If there is no pair of
triples, take $\kappa=1$. Otherwise the following ingredients are available.
\begin{enumerate}[label=(\roman*)]
\item\emph{Ideal identity} (\lref{PositivityTheory.ThKG.Loc.Cone.Id.idealId'}). There are signs $e_1,e_2$, independent
of all data, and a non-zero fractional ideal $\mathfrak a$ of $E$ such that: for every finite $v$ and $l\in E_v^\times$,
the matrix $u(l)$ is integral at $v$ (lies in the lattice generated by the $a\,xy^{\mathsf T}$, $x\in L_V$, $y\in L_W$) iff
$l\in\mathfrak a_v$; $\ord_w\mathfrak a=0$ for the places $w$ above $P$; and $\nu(\mathfrak a)=\nu(\mathfrak b_H)^{e_1}
\nu(\mathfrak b_{H'})^{e_2}$ for every continuous character $\nu$ trivial on $\A^\times$ and unramified outside $P$. The
ideal $\mathfrak a$ is computed from the coordinates of $x_H$ and $y_{H'}$ in the lattices; the signs come from the
duality between $L_X$ and its dual lattice.
\item\emph{Pure tensors} (\lref{PositivityTheory.ThKG.Loc.Cone.pureFac_of}). For a finite set $T'$ outside of which
the splitting characters are unramified, a pure tensor $\Phi=\thmKpure(f,\varphi)$ with vacua outside $T'$, $k\in K$ and
$l\in\thmKAE^\times$,
\begin{multline*}
e\bigl(\rho_\psi(u(l))\omega(k,1)\Phi\bigr)=\one\bigl[\,u(l_v)\text{ integral for all }v\notin T'\bigr]\\
\cdot\thmKpure\Bigl(e_\infty\bigl(\rho_\psi(u(l_\infty))\omega(k_\infty,1)f\bigr),\bigl(e_v(\rho_\psi(u(l_v))\omega(k_v,1)
\varphi_v)\bigr)_{v\in T'},(\thmKvac^{M_1}_v)_{v\notin T'}\Bigr),
\end{multline*}
where $k_\infty,k_v$ and $l_\infty,l_v$ are the components. This rests on the factorisation of the global operators into
local actions, the invariance of the vacuum under $K_v$ at places where the splitting characters are unramified
(\lref{PositivityTheory.ThKG.Loc.Cone.vacFixX_of}), and the formula $e_v(\rho_\psi(u(l))\thmKvac_v)=\thmKvac^{M_1}_v$
or $0$ according as $u(l)$ is integral at $v$ or not.
\item\emph{Local integrals} (\lref{PositivityTheory.ThKG.Loc.Cone.Fin.finCone2_of},
\lref{PositivityTheory.ThKG.Loc.Bump.archCone_of}). For $z\ge2$ the local cone integrals $\Lambda_v$ converge
absolutely and map $\thmKS$ into $\thmKS$, and the integrands are integrable over $K_v\times E_v^\times$; for $z\ge1$ the
archimedean cone integral $L$ converges absolutely, is a continuous linear map and maps the Fock space of $M$ into that of
$M_1$.
\item\emph{Fubini on $K$} (\lref{PositivityTheory.ThKG.Loc.Cone.kFub_of}). There is $\kappa_2\ne0$ such that for
integrable $h_\infty$ on $K_\infty$ and $h_v$ on $K_v$ ($v\in T'$) the function $k\mapsto h_\infty(k_\infty)\prod_{v\in T'}
h_v(k_v)$ is $\mu_K$-integrable with integral $\kappa_2\int h_\infty\,d\mu_{K,\infty}\prod_{v\in T'}\int h_v\,d\mu_{K,v}$; this
is the decomposition of Haar measure on $K=K_\infty\times\prod_vK_v$.
\item\emph{Spanning} (\lref{PositivityTheory.ThKG.Loc.Cone.Fin.schSpan}). Test functions are linear combinations of
pure tensors $\thmKpure(f,\varphi)$ with $\varphi_v=\thmKvac_v$ outside a finite set; this follows from the fact that the
locally constant compactly supported functions on $\A_f^{n\times d}$ are spanned by characteristic functions of products of
cosets.
\end{enumerate}
Note $k_A\ge2$, by the inequalities defining a good configuration.

Let $\kappa_1$ be the constant of \cref{pf:thmK:idelefub} for $\mathfrak a$, $z=k_A$ and $\mu$, and $\kappa=\kappa_1\kappa_2$. Fix
a pair of triples, $T'\supseteq T$, a pure tensor $\Phi=\thmKpure(f,\varphi)$ as in (3) and $y$. By (i) the indicator in (ii)
equals $\prod_{v\notin T'}\one_{\mathfrak a_v}(l_v)$. Hence, by (ii), the integrand of \eqref{pf:thmK:eq:cone} at $(k,l)$ is
\[
\Bigl(\prod_{v\notin T'}\thmKvac^{M_1}_v(y_v)\Bigr)\,g_\infty(k_\infty,l_\infty)\prod_{v\in T'}g_v(k_v,l_v)\cdot
\prod_{v\notin T'}\one_{\mathfrak a_v}(l_v)\,\nu(l)|l|^{k_A},
\]
with $g_\infty(k,l)=e_\infty(\rho_\psi(u(l))\omega(k,1)f)(y_\infty)$ and $g_v(k,l)=e_v(\rho_\psi(u(l))\omega(k,1)
\varphi_v)(y_v)$. Apply \cref{pf:thmK:idelefub} with $S=T$, $S_0$ the corresponding set of the character datum and $T'$: its
hypotheses hold since $\nu$ is unramified outside $T$, ramified above $T$, and $\ord_w\mathfrak a=0$ above $T\subseteq P$; the
integrability hypotheses are (iii). For each $k$ this gives the integral over $l$, multiplied by $\prod_{v\in T'\smallsetminus
T}u_v$, as $\kappa_1\nu(\mathfrak a)L^{S_0}(k_A,\nu)$ times the product of the local integrals over $l_\infty$ and the $l_v$
($v\in T'$). Integrating over $k$ with (iv) gives
\[
\Bigl(\prod_{v\in T'\smallsetminus T}u_v\Bigr)\thmKcone_A(\Phi)(y)=\kappa\,\nu(\mathfrak a)\,L^{S_0}(k_A,\nu)\,(Lf)(y_\infty)
\prod_{v\in T'}(\Lambda_v\varphi_v)(y_v)\prod_{v\notin T'}\thmKvac^{M_1}_v(y_v),
\]
which is \eqref{pf:thmK:eq:conefac} by (i), since $\nu=\alpha/\beta$ is trivial on $\A^\times$ (its restriction is
$\eta^{a'}\eta^{-(b+2)}$ with $a'+b$ even) and unramified outside $T\subseteq P$. Here $u_v=\int_{E_v^\times}
\one_{\mathfrak a_v}\nu|\cdot|^{k_A}d\mu_v\ne0$ by \cref{pf:thmK:idelefub}(1). For $v\notin T$ the vacuum is fixed by $K_v$
and $e_v(\rho_\psi(u(l))\thmKvac_v)=\one_{\mathfrak a_v}(l)\thmKvac^{M_1}_v$, so, as $\mu_{K,v}(K_v)=1$, $\Lambda_v(\thmKvac_v)
=u_v\thmKvac^{M_1}_v$. This proves (2) and (3).

For (1), the integrability statements above show that the double integral \eqref{pf:thmK:eq:cone} converges absolutely for
pure tensors with vacua outside a finite set; the set of test functions for which it converges absolutely is a vector
space on which $\thmKcone_A$ is linear, and by (v) it contains all test functions.
\end{proof}

\subsubsection{Bump functions for $\beta$-steps}\label{sec:pf:thmK:bumpb}

\begin{lemma}[Archimedean bump functions]\label{pf:thmK:bumpb}
Let $(G,G_1,H,H')$ be a $\beta$-step configuration, $e'\in\thmKEv(M',M'_1;\iota_H,\iota_{H'};I_B(H,H'))$ with archimedean
component $e'_\infty$, $K$ the stabiliser of $V'$ with a system of local measures, and let triples on $G$ and $G_1$ be
given. Let $L$ be the archimedean cone integral of $\thmKcone_B$ (with $u=u_B$, $\nu$ and $z=k_B+1$). If $l$ is a non-zero
continuous linear form on the archimedean Schwartz space of $M'_1$ such that $l(f_1)=l(f)$ whenever
$\omega'_1(g,1)\colon f\thmKacts{\infty}f_1$ with $g\in\Uni(V'_1)(F_\infty)$, then $l(Lf)\ne0$ for some $f$.
\formallabel{PositivityTheory.ThKG.Loc.Bump.bumpB_fin}
\end{lemma}

\begin{proof}
Suppose $l(Lf)=0$ for all $f$. Identify $E_\infty=E\otimes_FF_\infty$ with $\C^I$, $I$ the set of infinite places of $F$,
through the $\tau_w$. The character $\nu$ has infinity type $\kappa=(\kappa_w)$ with $\kappa_w=|m_\alpha(w)-m_\beta(w)|$ up
to sign, so $2k_B<\kappa_w$ for all $w$; write $r(c)=\prod_w(\bar c_w/|c_w|)^{\kappa_w}|c_w|^{2k_B}$ (the density of
$\nu(l)|l|^{k_B+1}d^\times l$ with respect to Lebesgue measure $dc$ on $\C^I$, up to conjugation of the variable). For an
archimedean Schwartz function $f$ put
\[
a_f(k,c)=l\Bigl(e'_\infty\bigl(\rho_\psi(\bar c\,y_Hx_{H'}^{\mathsf T})\,\omega'(k,1)f\bigr)\Bigr)\qquad(k\in K_\infty,\ c\in
\C^I).
\]
By the description of the archimedean operators (\lref{PositivityTheory.ThKF.A.archOp}) and the Heisenberg calculus,
$a_f$ is continuous, rapidly decreasing in $c$ uniformly in $k$, and $l(Lf)=c_0\int_{K_\infty}\int_{\C^I}a_f(k,c)r(c)\,dc\,dk$
with a constant $c_0\ne0$ (\lref{PositivityTheory.ThKG.Loc.Bump.archSlice_of}). Hence
\begin{equation}\label{pf:thmK:eq:bumpzero}
\int_{K_\infty}\int_{\C^I}a_f(k,c)\,r(c)\,dc\,dk=0\qquad\text{for every }f.
\end{equation}
\emph{Translations.} The form $\langle x\,x_{H'}^{\mathsf T},x'\,x_{H'}^{\mathsf T}\rangle$ vanishes, and $e'$ is invariant
under the Heisenberg translations by $I_B(H,H')$. Using this, replacing $f$ by the Heisenberg translate
$\rho_\psi(Y(\eta)x_{H'}^{\mathsf T})f$ by a vector $Y(\eta)$ built from $\eta\in\C^I$ changes $a_f(k,c)$ into
$a_f(k,c+\beta(k,\eta))$, where $\beta(k,\eta)$ is a continuous function, linear in $\eta$ and bounded by a multiple of
$|\eta|$, with $\langle\beta(k,\eta),t\rangle=\langle\eta,\pi(k,t)\rangle$ for a continuous map $\pi$ recording the
coordinate of $k^{-1}x_H$ along $x_H$. So \eqref{pf:thmK:eq:bumpzero} holds with $a_f(k,c+\beta(k,\eta))$ for all $\eta$.

\emph{Scaling.} At the points $(k,t)$ with $\pi(k,t)=\pi(1,1)$, the element $k$ acts on the line of $x_H$ through a Levi
element $m(t,\cdot)$, and the mixed-model formula (\cref{pf:thmK:mixx} for $V'$, applied to $(G,G_1)$, and the invariance
of $l$) gives $a_f(k,c)=\chi(t)|t|^{\rho}a_f(1,tc)$ for the splitting character $\chi=\nu$ and an exponent $\rho$; hence the
Fourier transform $\widehat{a_f(k,\cdot)}$ satisfies
\[
\widehat{a_f(k,\cdot)}(t)\,r(t)^{-1}|t|^{-1}\,\widehat{a_f(1,\cdot)}(1)^{-1}=|t|^{\rho-k_B-2}>0 .
\]
\emph{Choice of $f$.} Since $l\ne0$ and the image of $e'_\infty$ is dense, there is $f_0$ with $a_{f_0}(1,0)\ne0$; then the
Fourier transform of $c\mapsto a_{f_0}(1,c)$ does not vanish identically, and after a Heisenberg translation of $f_0$ we
obtain $f$ with $\widehat{a_f(1,\cdot)}(1)\ne0$.

\emph{Conclusion.} The following analytic statement (\lref{PositivityTheory.ThKG.Loc.Bump.coreB}) gives a
contradiction. Let $a$ be a continuous function on $K_\infty\times\C^I$ which decreases rapidly in the second variable,
$\beta,\pi$ as above, $2e<\kappa_w$ for all $w$, $r=r_{\kappa,e}$, and $c\in\C$; if $\int\int a(k,x+\beta(k,\eta))r(x)\,dx\,
dk=0$ for all $\eta$, then $\Re\bigl(c\,\widehat{a(k,\cdot)}(t)r(t)^{-1}|t|^{-1}\bigr)>0$ cannot hold at all points of
$\pi^{-1}(x_0)$ for a point $x_0$ of the image of $\pi$. Its proof combines a local functional equation
$\int\widehat H(t)r(t)^{-1}|t|^{-1}dt=Z\int H(x)r(x)\,dx$ ($Z\ne0$; \lref{PositivityTheory.ThKG.Loc.Bump.tate}),
Fourier inversion in $\eta$ and a bump function concentrated near $x_0$. We apply it with $a=a_f$, $e=k_B$ and
$c=\widehat{a_f(1,\cdot)}(1)^{-1}$.
\end{proof}

\subsubsection{The local cone integral at $v_1$}\label{sec:pf:thmK:flocb}

\begin{lemma}[The cone integral at $v_1$]\label{pf:thmK:flocb}
In the situation of \cref{pf:thmK:bumpb}, let $\Lambda_{v_1}$ be the local cone integral of $\thmKcone_B$ at $v=v_1$.
\begin{enumerate}
\item If $g\in\Uni(V')(F_{v_1})$, $\omega'(g,1)\colon\varphi\thmKacts{v_1}\varphi_1$ with $\varphi,\varphi_1\in\thmKS$, and
$l$ is a linear form vanishing on $\thmKaug_{v_1}(\omega'_1)$, then $l(\Lambda_{v_1}\varphi_1)=l(\Lambda_{v_1}\varphi)$.
\item $\Lambda_{v_1}$ maps $\thmKSo_{v_1}(\tau_1)$ for $M'$ into the span of $\thmKSo_{v_1}(\tau_1)$ for $M'_1$.
\end{enumerate}
\formallabel{PositivityTheory.ThKG.Loc.FL.fLocB_of}
\end{lemma}

\begin{proof}
Write $v=v_1$, let $\omega_v$ be a local representation for $\omega'$ at $v$
(\lref{PositivityTheory.ThKH.Locrep.locRepEx}) and
\[
A(\varphi)=\int_{E_v^\times}e'_v\bigl(\rho_\psi(u_B(s))\varphi\bigr)\,\nu(s)|s|^{z}\,d\mu_v(s),\qquad z=k_B+1\ge2,
\]
so that $\Lambda_v(\varphi)=\int_{K_v}A(\omega_v(k,1)\varphi)\,d\mu_{K,v}(k)$. For $\varphi\in\thmKS$ the integrand of $A$
takes finitely many values in $\thmKS$, on measurable sets, and $A(\varphi)\in\thmKS$
(\lref{PositivityTheory.ThKG.Loc.FL.tate'}); and $k\mapsto\omega_v(k,1)\varphi$ is locally constant. Hence linear
forms commute with both integrals.

(1) Let $f(x)=l(A(\omega_v(x,1)\varphi))$ for $x\in\Uni(V')(F_v)$; it is locally constant. For $u$ in the unipotent radical of
the stabiliser $P_H$ of the line $E x_H$, a Levi element $m=m(t,g_1)$ at $v$ and $x$, the mixed-model formula and the
invariance of $l$ under $\thmKaug_v(\omega'_1)$ give
\[
l\bigl(e'_v(\rho_\psi(u_B(s))\omega_v(um,1)\psi)\bigr)=\nu(t)\,|t|^{(a+b+2)/2}\,l\bigl(e'_v(\rho_\psi(u_B(ts))\psi)\bigr)
\]
(\lref{PositivityTheory.ThKG.Loc.FL.point}); substituting $s\mapsto t^{-1}s$ in $A$ and using $(a+b+2)/2=(b+1)+z$,
\[
f(umx)=\nu(t)|t|^{(a+b+2)/2}\,\nu(t)^{-1}|t|^{-z}f(x)=|t|^{b+1}f(x).
\]
So $f$ lies in the representation of $\Uni(V')(F_v)$ induced from the modulus character of $P_H(F_v)$ (note $\dim V'_1=b$).
On such functions integration over $K_v$ is invariant under right translations by $\Uni(V')(F_v)$
(\lref{PositivityTheory.ThKG.Loc.KInv.kInv}; this uses the Iwasawa decomposition $\Uni(V')(F_v)=P_H(F_v)K_v$ and the
decomposition of Haar measure along it). Since $\omega_v(k,1)\varphi_1=\omega_v(kg,1)\varphi$, we get
$l(\Lambda_v\varphi_1)=\int_{K_v}f(kg)\,dk=\int_{K_v}f(k)\,dk=l(\Lambda_v\varphi)$.

(2) Let $\varphi\in\thmKSo_v(\tau_1)$, $k\in K_v$, $s\in E_v^\times$, and let $\Phi=\thmKpure(f_0,\varphi[v])$ be a pure tensor
with vacua outside $v$ and $e'_\infty(f_0)\ne0$. Then $\Phi\in\thmKSo(M')$; $\omega'(k,1)\Phi\in\thmKSo(M')$ by the
stability of $\thmKSo(M')$ (\lref{PositivityTheory.ThKCT.soStab_of}); and $e'(\rho_\psi(u_B(s))\omega'(k,1)\Phi)\in
\thmKSo(M'_1)$, since evaluation maps of type $I_B$ map the Heisenberg translates of elements of $\thmKSo(M')$ by
$u_B(s)$ into $\thmKSo(M'_1)$ (\lref{PositivityTheory.ThKCT.soEv_of'}). This function is the pure tensor with archimedean
component $e'_\infty(f_0)\ne0$, the vacua outside $v$ and the component $e'_v(\rho_\psi(u_B(s))\omega_v(k,1)\varphi)$ at $v$;
hence that component lies in the span of $\thmKSo_v(\tau_1)$ for $M'_1$. As the integrals defining $\Lambda_v\varphi$ are
finite linear combinations of such components, (2) follows.
\end{proof}

\subsubsection{Mixed models}\label{sec:pf:thmK:mixx}

\begin{proposition}[Mixed models]\label{pf:thmK:mixx}
Let $(M,\omega)$ be a Weil pair for $(X,W)$ with characters $(\chi_X,\chi_W)$, $H$ a plane adapted to $(X,X_1)$ and
$(M_1,\omega_1)$ a Weil pair for $(X_1,W)$ with the same characters. Then there is an evaluation map
$e\in\thmKEv(M,M_1;\iota_H,1;\{x_Hz^{\mathsf T}\})$, and every such $e$ satisfies
\[
e\bigl(\omega(um,h)\Phi\bigr)=\chi_W(t)\,|t|^{d/2}\,\omega_1(g_1,h)\,e(\Phi)
\]
for all test functions $\Phi$, $u$ in the unipotent radical of the stabiliser of $Ex_H$, Levi elements $m=m(t,g_1)$ and
$h\in\Uni(W)(\A)$.
\formallabel{PositivityTheory.Mix.mixX_closed}
\end{proposition}

\begin{proof}
Write $\omega(um,h)=\omega(u,1)\omega(m,1)\omega(1,h)$. The existence of $e$ comes from a mixed model of the Weil
representation adapted to the decomposition of $X$ given by $x_H$, $y_H$ and $\iota_H(X_1)$
(\lref{GlobalTheory.ThKMixXEx.exists_isEv_mixedX_m}). The three identities
\[
e(\omega(1,h)\Phi)=\omega_1(1,h)e(\Phi),\quad e(\omega(m,1)\Phi)=\chi_W(t)|t|^{d/2}\omega_1(g_1,1)e(\Phi),\quad
e(\omega(u,1)\Phi)=e(\Phi)
\]
are \lref{PositivityTheory.Mix.theta_tower_hermitian_first_holds},
\lref{PositivityTheory.Mix.theta_tower_hermitian_second_holds} and \lref{GlobalTheory.MixXUnip.unip_trivial_m}.
They are derived from the defining properties of Weil representations and of evaluation maps: the intertwining of the
Heisenberg operators, Siegel's rule (which produces the factor $\chi_W(t)|t|^{d/2}$ for the Levi element), and the
invariance of $e$ under the Heisenberg translations by the matrices $x_Hz^{\mathsf T}$. Applying the three identities in
turn (all operators preserve test functions) gives the formula.
\end{proof}

\subsubsection{The modulus identity and spherical functions at inert places}\label{sec:pf:thmK:inert}

\begin{proposition}[Modulus identity at inert places]\label{pf:thmK:modst}
Let $B$ be good, $W\in\mathrm{Skew}_B(r,s)$ of dimension $d$, $v$ a finite place inert in $E$, $C_y\in M_{d\times
m}(\thmKAE)$ a matrix which is Lagrangian and adapted at $v$ with stabiliser $P$, $K$ the stabiliser of the local lattice
at $v$, $p\in P$ and $t$ with $pC_y=C_yt$. Then $[P\cap K:(P\cap K)\cap p^{-1}(P\cap K)p]=|\Nm\det t_v|_v^{m}\,[P\cap
K:(P\cap K)\cap p(P\cap K)p^{-1}]$.
\formal{PositivityTheory.Sph.ModSt}
\formallabel{PositivityTheory.Sph.modSt2}
\end{proposition}

\begin{proof}
Let $Y$ be the column span of $C_y$ at $v$, a Lagrangian subspace of $W_v$, so $d=2m$. $P$ is a closed subgroup of the
locally compact totally disconnected group $\Uni(W)(F_v)$ and $U=P\cap K$ is a compact open subgroup of $P$. Let $\lambda$
be a left Haar measure on $P$ and $\Delta$ its modular function, $\lambda(Ap)=\Delta(p)\lambda(A)$. Then
\[
\lambda(U\cap p^{-1}Up)=\lambda\bigl(p^{-1}(pUp^{-1}\cap U)p\bigr)=\Delta(p)\,\lambda(pUp^{-1}\cap U),
\]
and dividing $\lambda(U)$ by both sides gives $[U:U\cap p^{-1}Up]=\Delta(p)^{-1}[U:U\cap pUp^{-1}]$. It remains to show
$\Delta(p)=|\Nm\det t|^{-m}$. Choose a complementary Lagrangian $Y'$; then $P=M\ltimes N$, where $M\cong\GL(Y)$ consists of
the elements $m(t)$ acting by $t$ on $Y$ and preserving $Y'$, and $N$ is the unipotent radical, the group of elements
$n(X)$ acting trivially on $Y$ and on $W/Y$, with $X$ in an $F_v$-vector space $\mathfrak n$ of hermitian-type $m\times m$
matrices of $F_v$-dimension $m^2$ and $m(t)n(X)m(t)^{-1}=n(tX\sigma(t)^{\mathsf T})$. The group $N$ is a union of compact open
subgroups and $M\cong\GL_m(E_v)$ is unimodular, so $\Delta$ is trivial on $N$ and $\Delta(m(t))=|\det(\mathrm{Ad}\,m(t)|
_{\mathfrak n})|_v^{-1}$ (as $\lambda$ is the product of Haar measures of $M$ and $N$ in the coordinates $p=mn$). For
$t=\mathrm{diag}(t_1,\dots,t_m)$ the map $X\mapsto tX\sigma(t)^{\mathsf T}$ multiplies the diagonal entries (in $F_v$) by
$\Nm t_i$ and the pairs of off-diagonal entries (in $E_v$) by $t_i\sigma(t_j)$, of $F_v$-determinant $\Nm t_i\Nm t_j$;
so its determinant is $\prod_i\Nm t_i^{\,m}=\Nm(\det t)^m$. Both sides being characters of $\GL_m(E_v)$ which agree on the
diagonal torus and are trivial on the elementary matrices (unipotent), the formula holds for all $t$. As $\Delta(p)$
depends only on the image $t$ of $p$ in $M$, $\Delta(p)=|\Nm\det t|_v^{-m}$. (The formal proof carries out the same
computation by counting cosets in coordinates adapted to $C_y$ and the lattice:
\lref{PositivityTheory.Sph.M2.locMod}.)
\end{proof}

\begin{proposition}[Spherical functions at inert places]\label{pf:thmK:sphfun}
Assume the modulus identity of \cref{pf:thmK:modst}. Let $(B,G,X,\omega,\omega')$ be a triple, $v$ a finite place inert in
$E$, and $\omega_v,\omega'_v$ local representations for $\omega,\omega'$ at $v$ such that the vacua are fixed by the
stabiliser $K_v$ of the lattice of $W$ at $v$. Then there are linear forms $l,l'$ on the local Schwartz spaces, vanishing
on the augmentation modules, $K_v$-invariant and with $l(\thmKvac_v)=l'(\thmKvac'_v)=1$, such that
$l(\omega_v(1,h)\thmKvac_v)=l'(\omega'_v(1,h)\thmKvac'_v)$ for all $h\in\Uni(W)(F_v)$.
\formallabel{PositivityTheory.Sph.sphFun_of_modSt}
\end{proposition}

\begin{proof}
Since $v$ is inert and $v\notin T$, the characters are unramified at $v$. Let $m=(a+a')/2$.
\begin{enumerate}[label=(\roman*)]
\item There is a Lagrangian $C_y$ adapted at $v$ such that $\Uni(W)(F_v)=PK_v$ for its stabiliser $P$
(\lref{PositivityTheory.ThKH.Iwas.iwasawa}); moreover $h^{-1}\in K_vhK_v$ for all $h$
(\lref{PositivityTheory.Sph.cartanSt2}, from the Cartan decomposition), and each $K_vhK_v$ is a finite union of cosets
$pK_v$ (\lref{PositivityTheory.Sph.finSt}).
\item There is a linear form $l_0$ on $\thmKS(F_v^{a\times d})$ with $l_0(\thmKvac_v)\ne0$ which vanishes on the
augmentation module (\lref{PositivityTheory.ThKH.Evl.evVac}, \lref{PositivityTheory.Sph.evInvSt}) and satisfies
Siegel's rule $l_0(\omega_v(1,p)\varphi)=\beta(z)\,|\Nm\det t|^{a/2}\,l_0(\varphi)$ for $p\in P$ with $pC_y=C_yt$, where
$z\in E_v^\times$ is a unit-adjusted lift of $\det t$ (\lref{PositivityTheory.Sph.siegSt},
\lref{PositivityTheory.Sph.unitLiftSt}); similarly $l'_0$ for $\omega'$ with $\alpha(z)|\Nm\det t|^{a'/2}$; and
$\alpha(z)\beta(z)=1$ for these $z$ (\lref{PositivityTheory.Sph.charSt}).
\item Put $F(h)=l_0(\omega_v(1,h)\thmKvac_v)/l_0(\thmKvac_v)$ and $F'$ similarly. Then $F(phk)=c(p)F(h)$ and
$F'(phk)=c'(p)F'(h)$ for $p\in P$, $k\in K_v$, with $c(p)=\beta(z)|\Nm\det t|^{a/2}$, $c'(p)=\alpha(z)|\Nm\det t|^{a'/2}$,
and $F(1)=F'(1)=1$. By (ii) and the modulus identity,
\[
c(p)\,c'(p)=|\Nm\det t|^{m}=\frac{[U:U\cap p^{-1}Up]}{[U:U\cap pUp^{-1}]}\qquad(U=P\cap K_v).
\]
\item The forms $l(\varphi)=\int_{K_v}l_0(\omega_v(1,k)\varphi)\,dk/l_0(\thmKvac_v)$ and $l'$ (normalised Haar measure)
are $K_v$-invariant, vanish on the augmentation modules, satisfy $l(\thmKvac_v)=1$ (the vacuum is $K_v$-fixed), and
$l(\omega_v(1,h)\thmKvac_v)=\int_{K_v}F(kh)\,dk$ (\lref{PositivityTheory.Sph.exists_form},
\lref{PositivityTheory.Sph.avgSt}).
\item Finally, for functions $F,F'$ as in (iii) the averages $\int_{K_v}F(kh)\,dk$ and $\int_{K_v}F'(kh)\,dk$ coincide
(\lref{PositivityTheory.Sph.countSt}). This is the classical relation between the spherical functions attached to a
character $c$ of $P$ and to $c^{-1}\Delta^{-1}$, evaluated at $h$ and $h^{-1}$ \cite{Casselman1980}: decomposing $K_vhK_v$
into cosets $p_iK_v$ and $K_vh^{-1}K_v=K_vhK_v$ into cosets $p'_jK_v$, both averages become finite sums of the values
$c(p_i)$, resp.\ $c'(p'_j)$, weighted by indices of the form $[U:U\cap p^{\pm1}Up^{\mp1}]$, and the relation in (iii)
identifies the two sums.
\end{enumerate}
\end{proof}

\subsubsection{The unramified doubling integral at inert places}\label{sec:pf:thmK:dblinert}

\begin{proposition}[Unramified doubling integral]\label{pf:thmK:dblinert}
Let $B$ be good, $Y\in\mathrm{Herm}_B(r,s)$ of dimension $n$, $\chi$ a Hecke character with $\chi|_{\A^\times}=\eta^0$,
$G_2(Y)$ the doubled group with stabiliser $K$ and doubling embedding $\iota$, $v$ a finite place inert in $E$ at which
$\chi$ is unramified, and $\mu$ the Haar measure on $\Uni(Y)(F_v)$ with $\mu(\Uni(Y)(F_v)\cap K)=1$. For $\Re z>n/2$ and
every $K$-spherical section $f$ at $v$ of the degenerate principal series of $G_2(Y)(F_v)$ with parameter $z+n/2$ and
character $\chi$, normalised by $f(1)=1$, the function $u\mapsto f(\iota(u,1))$ is integrable and
\[
\int_{\Uni(Y)(F_v)}f(\iota(u,1))\,d\mu(u)\cdot\prod_{i=0}^{n-1}\bigl(1-(-1)^{n-1-i}\Nm v^{-(2z+i+1)}\bigr)^{-1}
=\prod_{i=0}^{n-1}L_v\bigl(z+\tfrac n2-i,\chi\bigr).
\]
\formallabel{GlobalTheory.ThK.Dbl.Inert.main}
\end{proposition}

\begin{proof}
The place $w$ of $E$ above $v$ is unique and $E_w/F_v$ is the unramified quadratic extension. Choosing a uniformiser of
$F_v$ (which is also one of $E_w$) and a basis of the lattice of $Y$ at $v$ adapted to the hermitian form, one identifies
$\Uni(Y)(F_v)$ with an unramified unitary group, $K\cap\Uni(Y)(F_v)$ with its standard maximal compact subgroup,
$G_2(Y)(F_v)$ with the corresponding quasi-split unitary group in $2n$ variables and $f$ with the normalised spherical
section in standard coordinates; since $\chi$ is trivial on $\A^\times$ and unramified at the inert place $w$, $\chi_w$ is
trivial, and the factors $L_v(\cdot,\chi)$ are those of the trivial character of $E_w$. These identifications preserve the
measures and the integrals. The statement in these coordinates is the classical computation of the unramified doubling
zeta integral of the trivial representation \cite{GelbartPiatetskiShapiroRallis1987,Li1992}, which the formal proof carries
out by an explicit count of lattices (\lref{GlobalTheory.CompanionTheory.InertZetaTop.unramDoublingInert}): the integral
is decomposed along double cosets of the maximal compact subgroup, the values of $f$ are powers of $\Nm v^{-(z+n/2)}$
determined by elementary divisors (computed through Smith normal forms and the action of a Levi subgroup on cosets), the
volumes are computed by counting cosets in Heisenberg-type groups, and summing the resulting series gives the right-hand side
divided by the normalising factor.
\end{proof}

\subsubsection{Local multiplicity one}\label{sec:pf:thmK:multone}

\begin{proposition}[Archimedean multiplicity one]\label{pf:thmK:archmult}
Assume \ref{pf:thmK:H:lie2}. Let matched data on $G$ and $G_1$ be related by an $\alpha$- or $\beta$-step with Levi
character $(\chi,r)$, and let $H'$ be adapted to $(W,W_1)$. If $\lambda,\lambda'\colon\Pi_\infty(G)\to\Pi_\infty(G_1)$ are
linear maps compatible, through the archimedean matching maps, with the action of the Levi torus of $H'$ by $(\chi,r)$, of
its unipotent radical and of the Levi factor $\Uni(W_1)$, and $\lambda\ne0$, then $\lambda'=c\lambda$ for some $c$.
\formallabel{PositivityTheory.ThKG.CayleyTransformDarboux.archimedean_multiplicity_one_of_lie}
\end{proposition}

\begin{proof}
Choose, at each infinite place, frames of $W$ and $W_1$ adapted to the Cartan involutions and to $H'$ (a basis completing
$x_{H'},y_{H'}$ and the image of $\iota_{H'}$). Through the Fock model, the archimedean matching maps identify
$\Pi_\infty(G)$ and $\Pi_\infty(G_1)$ with modules over the complexified Lie algebras of $\Uni(W)(F_\infty)$ and
$\Uni(W_1)(F_\infty)$, on which these Lie algebras act by explicit operators in the chosen frames (by \ref{pf:thmK:H:lie2}
and the first-order formulas; \lref{PositivityTheory.ThKG.CayleyTransformDarboux.exists_mod}). The equivariance of a map
$\lambda$ under the Levi torus (through the derivative of $(\chi,r)$), the unipotent radical and $\Uni(W_1)$ translates,
by differentiating along one-parameter subgroups, into a list of commutation relations between $\lambda$ and these
operators. Two maps satisfying these relations, the first non-zero, are proportional
(\lref{PositivityTheory.ThKG.CayleyTransformDarboux.par_line}): the relations determine such a map from its value on a
generator of the cyclic module $\Pi_\infty(G)$, and confine that value to a line.
\end{proof}

\begin{proposition}[Multiplicity one at the places of $T$]\label{pf:thmK:rammult}
Let $v$ be a place with an embedding $t\colon E\to F_v$ at which $\chi_W$ is ramified, $(M,\omega)$ a Weil pair for $(X,W)$
with characters $(\chi_X,\chi_W)$, $H'$ adapted to $(W,W_1)$ with $n\le d_1$, and $(\tau,P)$ an irreducible representation of
$\Uni(W_1)(F_v)\cong\GL_{d_1}(F_v)$ of countable dimension. Let $\Lambda,\Lambda'$ be linear maps from
$\thmKS(F_v^{n\times d})$ to $P$ which vanish on $\thmKaug_v(\omega)$ and are equivariant for the parabolic subgroup of $H'$
with the Levi character $(\chi_X,r)$ and the action $\tau$ of $\Uni(W_1)(F_v)$. If $\Lambda\ne0$ on $\thmKS$, then
$\Lambda'=c\Lambda$ on $\thmKS$. In particular hypothesis \ref{pf:thmK:H:ram} of \cref{pf:thmK:closesph} holds.
\formallabel{PositivityTheory.BernsteinRamification.Close.ramifiedMultiplicityOne}
\end{proposition}

\begin{proof}
\emph{Linear models.} At the split place $v$ the dictionary for linear models
(\lref{PositivityTheory.KCo.section_statement}; \lref{PositivityTheory.ThKG.StdDict.dictX}) gives a bijection
$\iota$ of $\thmKS$ onto the locally constant compactly supported functions on $M_{n\times d}(F_v)$ in which
$\Uni(X)(F_v)\cong\GL_n(F_v)$ acts by left translation twisted by $\chi_W$, and an element of $\Uni(W)(F_v)$ corresponding to
$CpC^{-1}$ ($p\in\GL_d(F_v)$, $C$ a frame adapted to $H'$) acts by $\Phi\mapsto\chi_X(\det p)|\det p|^{n/2}\Phi(\cdot\,p)$.
In these coordinates the parabolic subgroup of $H'$ becomes the group of block matrices $p$ preserving the decomposition
$F_v^{d}=F_v^{d_1}\oplus F_v^2$ in the appropriate way, its Levi torus $m(z,1)$ corresponds to $\mathrm{diag}(1,x,y)$ with
$x,y$ given by $z$ (\lref{PositivityTheory.ThKG.LocalRamifiedLeviRestriction.idele_two},
\lref{PositivityTheory.ThKG.LocalRamifiedLeviRestriction.is_levi}), and $\Lambda\circ\iota^{-1}$, $\Lambda'\circ
\iota^{-1}$ become linear maps $L,L'$ which are invariant under $\GL_n(F_v)$, transform under the block-diagonal elements
$\mathrm{diag}(1,x,y)$ with $x,y$ units by a character $\nu$ with $\nu(l)\ne1$ for some unit $l$ (because $\chi_W$ is
ramified), and are $\GL_{d_1}(F_v)$-equivariant.

\emph{Reduction to full rank.} Because $\nu$ is ramified, $L$ and $L'$ are determined by their restrictions $\mu,\mu'$ to
the space $S$ of functions of the form $\psi\otimes\one\otimes\one$, with $\psi$ supported on the matrices of rank $n$ in
$M_{n\times d_1}(F_v)$ and the characteristic functions of the integral vectors in the last two columns; $S$ is stable under
$\GL_{d_1}(F_v)$, and $\mu\ne0$
(\lref{GlobalTheory.CompanionTheory.BernNarrow.local_multiplicity_one_irreducible_GL}).

\emph{Irreducibility and Schur's lemma.} The kernel of $\mu$ contains the augmentation module of $S$ and is
$\GL_{d_1}(F_v)$-stable; using the irreducibility of the representations of $\GL_{d_1}(F_v)$ induced from characters of
maximal parabolic subgroups --- for smooth unitary characters $\chi_1,\chi_2$ of $F_v^\times$ and $a\ne b$ or
$\chi_1\ne\chi_2$, the representation of $\GL_{a+b}(F_v)$ induced from $\chi_1\circ\det\otimes\chi_2\circ\det$ is
irreducible (\lref{PositivityTheory.BernsteinRamification.narrowBern}, deduced from the geometric lemma of
Bernstein--Zelevinsky \cite{BernsteinZelevinsky1977,Bernstein1984}) --- one obtains $\ker\mu\subseteq\ker\mu'$. Thus $\mu'$
factors through $\mu(S)=P$, giving a $\GL_{d_1}(F_v)$-endomorphism $A$ of the irreducible representation $P$. Since $P$ has
countable dimension, $A$ is a scalar: otherwise $A-\lambda$ is injective, hence (its image being a non-zero
subrepresentation) bijective, for every $\lambda\in\C$, and for $0\ne x\in P$ the uncountably many vectors
$(A-\lambda)^{-1}x$ would be linearly independent. Hence $\mu'=c\mu$ and $\Lambda'=c\Lambda$ on $\thmKS$.
\end{proof}

\subsubsection{Local matching at split places}\label{sec:pf:thmK:finsplit}

\begin{proposition}[Matching maps at split places]\label{pf:thmK:finsplit}
For every triple and every finite place $v\notin T$ at which $E/F$ splits there are a non-zero complex vector space $\Pi_v$
and linear maps $p_v,p'_v$ satisfying the conditions of \cref{pf:thmK:def:maps} at $v$. In particular hypothesis
\ref{pf:thmK:H:finsplit} of \cref{pf:thmK:closesph} holds.
\formallabel{PositivityTheory.BernsteinRamification.Close.finSplit}
\end{proposition}

\begin{proof}
At $v$ fix $t\colon E\to F_v$ and linear models (\lref{PositivityTheory.ThK.SplitDict.split_place_dictionary}): $\Uni(W)(F_v)\cong
\GL_{a+a'}(F_v)$, $\Uni(V)(F_v)\cong\GL_a(F_v)$, $\Uni(V')(F_v)\cong\GL_{a'}(F_v)$, and the local Weil representations become
the natural actions on the locally constant compactly supported functions on $M_{a\times(a+a')}(F_v)$, resp.\
$M_{a'\times(a+a')}(F_v)$, twisted by the unramified characters $\alpha_v,\beta_v$ (\lref{PositivityTheory.ThK.SplitAsm.split_character_facts};
the actions are unique, \lref{PositivityTheory.KCo.action_uniqueness}). Let $\Pi_v$ be the representation of
$\GL_{a+a'}(F_v)$ induced from the characters $\beta_v\circ\det\otimes\alpha_v\circ\det$ of the Levi subgroup
$\GL_a\times\GL_{a'}$ of a maximal parabolic subgroup; it is irreducible, since $a\ne a'$
(\lref{PositivityTheory.BernsteinRamification.narrowBern}, as in the proof of \cref{pf:thmK:rammult}). Define $p_v$ and
$p'_v$ by local zeta integrals of Godement--Jacquet type \cite{GodementJacquet1972}: $p_v(\varphi)$ is the function
$g\mapsto\int_{\GL_a(F_v)}(\iota\varphi)\bigl((h,0)g\bigr)\beta_v(\det h)|\det h|^{s}\,dh$ at the appropriate point $s$, and
$p'_v$ similarly with $\GL_{a'}$ and $\alpha_v$ (\lref{GlobalTheory.CompanionTheory.BernBridge.Split.unrMaps_of_lt}).
These integrals converge for $a'<a$ (\lref{PositivityTheory.ThK.SplitLoc.Unr.conv_finite}); they are invariant under
$\GL_a(F_v)$, resp.\ $\GL_{a'}(F_v)$, so they vanish on the augmentation modules; they are equivariant for $\GL_{a+a'}(F_v)$;
they are surjective onto the irreducible $\Pi_v$ as soon as they are non-zero; and on the vacua (characteristic functions of
the integral matrices) both take the value of the normalised spherical vector, by the unramified computation of the zeta
integrals (\lref{PositivityTheory.ThK.SplitLoc.Unr.pUnr'_eq}). A function vanishing on the integrals over all
translates is zero (\lref{PositivityTheory.ThK.SplitDict.induced_vanish_of_vanish_on_integral}), which
gives the compatibility of $p_v$ and $p'_v$ with the action of $\Uni(W)(F_v)$ required in \cref{pf:thmK:def:maps}.
\end{proof}

\subsubsection{Existence of Weil representations}\label{sec:pf:thmK:weilex}

\begin{proposition}[Existence of Weil representations]\label{pf:thmK:weilex}
For good $B$, $X\in\mathrm{Herm}_B(r,s)$, $W\in\mathrm{Skew}_B(r',s')$, a good model $M$ for $(X,W)$ and Hecke characters
$\chi_X,\chi_W$ with $\chi_X|_{\A^\times}=\eta^n$, $\chi_W|_{\A^\times}=\eta^d$, there is a Weil representation for $M$
with splitting characters $(\chi_X,\chi_W)$.
\formal{PositivityTheory.ThKG.Tow.WeilEx}
\formallabel{PositivityTheory.ThKW.weilEx}
\end{proposition}

\begin{proof}
This is the existence of the global Weil representation of $\Uni(X)\times\Uni(W)$ \cite{Weil1964} with the splitting of
Kudla \cite{Kudla1994}, in the axiomatic form of \cref{pf:thmK:def:models}. The formal construction proceeds in four steps.
\begin{enumerate}[label=(\roman*)]
\item\emph{Hermitian reformulation.} Choose $\xi_0\in E^\times$ with $\sigma(\xi_0)=-\xi_0$. Then $J_S=-\xi_0J_W^{-1}$ is
hermitian and non-degenerate, as is $J_T=J_X^{\mathsf T}$, and the lattices of $X$ and $W$ give a self-dual lattice for the
pair $(J_S,J_T)$ (\lref{PositivityTheory.ThKW.latH_of}).
\item\emph{A Weil datum with characters $\eta^{n},\eta^{d}$.} Choose a Hecke character $\eta$ of $E$ with
$\eta|_{\A^\times}=\eta_{E/F}$ (\lref{PositivityTheory.ThKW.exists_eta}). The theta distribution attached to the
lattice is invariant under the rational Heisenberg group (\lref{GlobalTheory.T.theta_function_heisenberg_invariant});
the lattice model then gives a Weil datum for the pair (\lref{PositivityTheory.ThKW.latRep}), and Kudla's splitting
with the characters $\eta^n,\eta^d$ turns it into a representation of the two unitary groups
(\lref{PositivityTheory.ThKW.datumEta}).
\item\emph{Twisting.} Multiplying the operators by $(\chi_X\eta^{-n})(\det g)$ and $(\chi_W\eta^{-d})(\det h)$, characters
trivial on $\A^\times$, changes the splitting characters to $(\chi_X,\chi_W)$
(\lref{GlobalTheory.T.weil_datum_twist_by_hecke_characters_trivial_base}).
\item\emph{Transport to the model.} Transporting the resulting representation to the Schwartz space of $M$ through the
Darboux coordinates of $M$ (which map rational matrices to rational vectors since $M$ is good) gives a representation
$\omega$ which preserves test functions and the admissible linear forms, intertwines the Heisenberg operators, is unitary
and continuous (\lref{GlobalTheory.T.exists_weilTransport_adelicSchwartz_cmExtension}); the extensions to the doubled
groups and Siegel's rule are obtained by applying the same construction to the doubled pairs
(\lref{PositivityTheory.ThKW.S.exists_dbl_target}). These are exactly the conditions of \cref{pf:thmK:def:models}.
\end{enumerate}
\end{proof}


% ===========================================================================
\subsection{Proof of Theorem K}\label{sec:pf:thmK:final}

\begin{theorem}[Reduction to the leaves]\label{pf:thmK:closesph}
Assume the following fourteen statements.
\begin{enumerate}[label=(H\arabic*),leftmargin=*]
\item\label{pf:thmK:H:dict} If the dictionary statement at split places holds (at every finite place $v$ with an embedding
$t\colon E\to F_v$, every Weil pair admits linear models of conjugate type at $v$, and in each of them the local actions of
$\Uni(X)(F_v)$ and $\Uni(W)(F_v)$ are the standard actions of $\GL_n(F_v)\times\GL_d(F_v)$ twisted by two unitary
characters), then \textup{(Coi)} holds.
\item\label{pf:thmK:H:so} If \ref{pf:thmK:L:kfin} holds, then for every triple, an element of the span of
$\thmKSo_{v_1}(\tau_1)$ lying in $\thmKaug_{v_1}(\omega')$ lies in the augmentation module generated inside that span.
\item\label{pf:thmK:H:ram} Local multiplicity one at the places of $T$: for matched data on $G$ and $G_1$ related by a step, a
plane $H'$ and $v\in T$, two linear maps $\Pi_v(G)\to\Pi_v(G_1)$ which are equivariant for the parabolic subgroup of $H'$ with
the character of the step, the first non-zero, are proportional.
\item\label{pf:thmK:H:finram} Matching maps exist at the places of $T$ (with the zeta condition).
\item\label{pf:thmK:H:finsplit} Matching maps exist at the split places outside $T$.
\item\label{pf:thmK:H:sph} For every triple and finite $v$, each of the two local Weil representations at $v$ is either
spherical (it admits a spherical vector in every local model) or irreducible (every $\Uni(W)(F_v)$-stable subspace of $\thmKS$
containing the augmentation module is the augmentation module or $\thmKS$).
\item\label{pf:thmK:H:fininert} Matching maps exist at the inert places outside $T$.
\item\label{pf:thmK:H:lie2} The derivatives of a Weil representation along one-parameter subgroups act on the Fock model
by the standard second-order operators, shifted by half the infinity type of the splitting character.
\item\label{pf:thmK:H:archlie} For a Weil pair and one-parameter subgroups $\gamma$ of $X$ and $\delta$ of $W$, every $f$ in
the Fock space is differentiable along $t\mapsto\omega(\gamma(t),\delta(t))$ with derivative in the Fock space, compatibly
with the archimedean operators implementing $\omega(\gamma(t),\delta(t))$.
\item\label{pf:thmK:H:kreal} The leaf \ref{pf:thmK:L:kreal}.
\item\label{pf:thmK:H:siegel} Reduction theory, \cref{pf:thmK:siegel}.
\item\label{pf:thmK:H:eis} The meromorphic continuation, \cref{pf:thmK:eis}.
\item\label{pf:thmK:H:sw} The Siegel--Weil formula, \cref{pf:thmK:sw:thm}.
\item\label{pf:thmK:H:weilex} The existence of Weil representations, \cref{pf:thmK:weilex}.
\end{enumerate}
Then $\sfname{ThetaIntegralsAgree}$ and $\sfname{LocalMatching}$ hold.
\formallabel{PositivityTheory.Sph.Tr.closeSph}
\end{theorem}

\begin{proof}
By \cref{pf:thmK:agree} it suffices to establish the twenty-six leaves, \textup{(Dec)} and \textup{(Coi)} under these
hypotheses. \textup{(Coi)} follows from \ref{pf:thmK:H:dict} and the dictionary statement (\cref{pf:thmK:localmatching}). The
leaves \ref{pf:thmK:L:evzero}--\ref{pf:thmK:L:stab}, \ref{pf:thmK:L:levi}, \ref{pf:thmK:L:poldense} and, using
\ref{pf:thmK:H:siegel}, \ref{pf:thmK:L:weilconv}, \ref{pf:thmK:L:socont} and \textup{(Dec)} are \cref{pf:thmK:analytic}. The
towers \ref{pf:thmK:L:towerA}, \ref{pf:thmK:L:towerB} are \cref{pf:thmK:towers} (with \ref{pf:thmK:H:weilex});
\ref{pf:thmK:L:facne} is \cref{pf:thmK:facne}; \ref{pf:thmK:L:matchex} is \cref{pf:thmK:matchex}; the constant-term leaves
\ref{pf:thmK:L:splitex}--\ref{pf:thmK:L:ctsc} are \cref{pf:thmK:ctleaves}, where the general formulas \cref{pf:thmK:gen} are
applied to the collection consisting of the statements on mixed models (\lref{PositivityTheory.Mix.mixW_closed} and
\lref{PositivityTheory.Mix.mixX_closed}, which describe the effect of Levi elements on evaluation maps), the basic
properties of theta functions (\cref{pf:thmK:analytic}(a)), \ref{pf:thmK:L:evzero}, \ref{pf:thmK:L:evcone},
\ref{pf:thmK:L:stab}, the existence of models with Weil representations (by \ref{pf:thmK:H:weilex}), the existence of adapted
planes, the topological properties of the unitary groups, \ref{pf:thmK:L:weilconv} and the decay statements for
$\thmKSo$-data (by \ref{pf:thmK:H:siegel}). \ref{pf:thmK:L:anall} is \cref{pf:thmK:anall}; \ref{pf:thmK:L:polre} is
\cref{pf:thmK:polre} (with \ref{pf:thmK:H:sph}, \ref{pf:thmK:H:lie2}, \ref{pf:thmK:H:archlie}); \ref{pf:thmK:L:nonzeroan} is
\cref{pf:thmK:nonzeroan}, where \cref{pf:thmK:nonzerobot} uses \ref{pf:thmK:H:kreal}, \ref{pf:thmK:H:eis},
\ref{pf:thmK:H:sw}; \ref{pf:thmK:L:polzero} is \cref{pf:thmK:polzeroprop} (with \ref{pf:thmK:H:eis}, \ref{pf:thmK:H:sw});
\ref{pf:thmK:L:locA}, \ref{pf:thmK:L:locB} are \cref{pf:thmK:locAB} (with \ref{pf:thmK:H:so}, \ref{pf:thmK:H:ram},
\ref{pf:thmK:H:sph}, \ref{pf:thmK:H:lie2}, \ref{pf:thmK:H:archlie}); \ref{pf:thmK:L:endA}, \ref{pf:thmK:L:endB} are
\cref{pf:thmK:endsteps} (from \ref{pf:thmK:H:dict},
\ref{pf:thmK:H:finram}, \ref{pf:thmK:H:finsplit}, \ref{pf:thmK:H:fininert}, \ref{pf:thmK:H:lie2}, \ref{pf:thmK:H:archlie},
\ref{pf:thmK:H:siegel}--\ref{pf:thmK:H:weilex}); \ref{pf:thmK:L:kreal} is \ref{pf:thmK:H:kreal}; and \ref{pf:thmK:L:kfin} is
\cref{pf:thmK:localmatching} (with \ref{pf:thmK:H:finram}, \ref{pf:thmK:H:finsplit}, \ref{pf:thmK:H:fininert}).
\end{proof}

\begin{theorem}[Theorem K from the analytic inputs]\label{pf:thmK:head}
The modulus identity of \cref{pf:thmK:modst}, the meromorphic continuation of \cref{pf:thmK:eis} and the Siegel--Weil formula
of \cref{pf:thmK:sw:thm} imply $\sfname{ThetaIntegralsAgree}$ and $\sfname{LocalMatching}$; hence the latter two statements
imply them.
\formallabel{PositivityTheory.ThKR.theoremK_head2_of_mod}\formallabel{PositivityTheory.ThKR.theoremK_head2}
\end{theorem}

\begin{proof}
We verify the hypotheses of \cref{pf:thmK:closesph}. Hypotheses \ref{pf:thmK:H:dict}--\ref{pf:thmK:H:finsplit} are
\lref{PositivityTheory.BernsteinRamification.Close.kCoinv},
\lref{PositivityTheory.BernsteinRamification.Close.soCoinv},
\lref{PositivityTheory.BernsteinRamification.Close.ramifiedMultiplicityOne},
\lref{PositivityTheory.BernsteinRamification.Close.finRam} and
\lref{PositivityTheory.BernsteinRamification.Close.finSplit}, which prove exactly these statements, all from a filtration
of the local Weil representations at split places by the geometric lemma of Bernstein--Zelevinsky type
\cite{BernsteinZelevinsky1977}. Hypothesis \ref{pf:thmK:H:sph} is \lref{PositivityTheory.Sph.sphTri_of}, which deduces it
from the spherical analysis at split places (\lref{PositivityTheory.Sph.sphSplit_geo}) and the generation of the
coinvariants by the vacuum at inert places (\lref{PositivityTheory.Sph.vacInert}). Hypothesis \ref{pf:thmK:H:fininert} is
\lref{PositivityTheory.Sph.finInert_of}, applied as in \cref{pf:thmK:localmatching}; it is the only place where the modulus
identity enters, through \lref{PositivityTheory.Sph.sphFun_of_modSt}. Hypotheses \ref{pf:thmK:H:lie2},
\ref{pf:thmK:H:archlie}, \ref{pf:thmK:H:kreal} are \lref{PositivityTheory.ThKR.lieSecond2},
\lref{PositivityTheory.ThKR.archLie'} and \lref{PositivityTheory.ThKR.kReal'}, which prove them from the local
oscillator representation on Schwartz space (\lref{PositivityTheory.ThKR.localLift_holds}). Hypothesis
\ref{pf:thmK:H:siegel} is \cref{pf:thmK:siegel}; \ref{pf:thmK:H:eis} and \ref{pf:thmK:H:sw} are the hypotheses; and
\ref{pf:thmK:H:weilex} is \cref{pf:thmK:weilex}. For the second assertion, the modulus identity holds by
\cref{pf:thmK:modst}.
\end{proof}

\begin{proof}[Proof of \cref{pf:thmK:main}]
By \cref{pf:thmK:eis,pf:thmK:sw:thm} the meromorphic continuation and the Siegel--Weil formula hold; apply
\cref{pf:thmK:head}.
\end{proof}

% ===========================================================================
\subsection{Remarks on the formal statements}\label{sec:pf:thmK:remarks}

\begin{remark}[Unused hypothesis]\label{pf:thmK:rem:unused}
The formal statement of \cref{pf:thmK:main} takes as hypothesis the conjunction of three propositions (a local-global
compatibility for Galois representations attached to polarised pairs, which appears twice, and the existence of a local
Langlands correspondence). It is not used: the proof applies \cref{pf:thmK:head} to \cref{pf:thmK:eis,pf:thmK:sw:thm}, and the
formal statements supplying these two theorems do not use the hypothesis either. The suffix ``of literature'' of the formal name is
therefore not descriptive.
\end{remark}

\begin{remark}[Scope of the statements]\label{pf:thmK:rem:scope}
$\sfname{ThetaIntegralsAgree}$ and $\sfname{LocalMatching}$ are universally quantified over all data forming a triple; they say
nothing about the existence of such data. The conditions on a triple are strong ($E/F$ unramified at all finite places with
$E/\Q$ Galois, inert places odd of norm $>3$, $m_\alpha-m_\beta>a+a'$, $\lfloor(a-a')/2\rfloor$ and $a'$ larger than
$2(p+q)+1$). The step relations assert the existence of signs $(e_1,e_2,s)$ and of a non-zero constant $C$ depending on the
geometric data and the real maps, but not their values.
\end{remark}

\begin{remark}[Weil representations]\label{pf:thmK:rem:weil}
Weil representations are characterised by properties (\cref{pf:thmK:def:models}); no uniqueness is asserted, and the theorem
holds for every representation with these properties. The local actions at a place are extracted from the global operators
through the relation $T\colon\varphi\thmKacts{v}\varphi_1$, which holds vacuously when no admissible function $r$ exists; in
all uses $\varphi\in\thmKS$.
\end{remark}

\begin{remark}[Values of divergent expressions]\label{pf:thmK:rem:junk}
Theta series are unconditional sums and theta integrals are Bochner integrals, both $0$ when they diverge; convergence is
asserted separately in the decay condition. Partial $L$-functions are infinite products and occur only at integers $\ge2$,
where they converge absolutely. The zeta integrals in the zeta condition of \cref{pf:thmK:def:maps} are defined with a finite
sum over Cartan cells which is $0$ when infinitely many cells contribute, so that condition has its intended meaning only for
functions with finitely many contributing cells.
\formal{GlobalTheory.SplitFlip.shellSum}
\end{remark}

\begin{remark}[Hypotheses not used]\label{pf:thmK:rem:below}
In \ref{pf:thmK:L:endA} and \ref{pf:thmK:L:endB} the hypothesis $\mathcal B(p+q)$ is not used by the proof
(\cref{pf:thmK:endAB}). In \cref{pf:thmK:keyB,pf:thmK:keyA} and in the leaves \ref{pf:thmK:L:ctbz}--\ref{pf:thmK:L:ctsc},
\ref{pf:thmK:L:locA}, \ref{pf:thmK:L:locB}, the constant is taken to be $1$ when no triple exists on the configurations, and
the statements are then vacuous.
\end{remark}
